% concatenated from imeval.tex
\documentclass[a4paper, 12pt]{article}

\usepackage[sort&compress]{natbib}
\bibpunct{(}{)}{;}{a}{}{,} 

\usepackage{amsthm, amsmath, amssymb, mathrsfs, multirow, url, subfigure}
\usepackage{graphicx} 
\usepackage{ifthen} 
\usepackage{amsfonts}
\usepackage[usenames]{color}
\usepackage{fullpage}
\usepackage{tikz}

%\RequirePackage[OT1]{fontenc} 
%\RequirePackage{hypernat}
%\RequirePackage[colorlinks]{hyperref}
%\hypersetup{citecolor=blue, linkcolor=red, urlcolor=black}


%\numberwithin{equation}{section} 

\theoremstyle{plain} 
\newtheorem{thm}{Theorem}
\newtheorem{cor}{Corollary}
\newtheorem{prop}{Proposition}
\newtheorem{lem}{Lemma}
\newtheorem*{lem0}{Lemma}
\newtheorem*{thm0}{Theorem}
\newtheorem*{cor0}{Corollary}


\theoremstyle{definition}
\newtheorem{defn}{Definition}

\theoremstyle{remark} 
\newtheorem{ex}{Example}
\newtheorem{remark}{Remark}
\newtheorem*{remark0}{Remark}
\newtheorem{note}{Note}



\newcommand{\prob}{\mathsf{P}} 
\newcommand{\E}{\mathsf{E}}

\newcommand{\bel}{\mathsf{bel}}
\newcommand{\pl}{\mathsf{pl}}
\newcommand{\pval}{\mathsf{pval}}

\newcommand{\pois}{{\sf Pois}}
\newcommand{\bin}{{\sf Bin}}
\newcommand{\unif}{{\sf Unif}}
\newcommand{\nm}{{\sf N}}
\newcommand{\expo}{{\sf Exp}}
\newcommand{\gam}{{\sf Gamma}}
\newcommand{\stt}{{\sf t}}
\newcommand{\dir}{{\sf Dir}}
\newcommand{\chisq}{{\sf ChiSq}}

\newcommand{\A}{\mathcal{A}} 
\renewcommand{\AA}{\mathbb{A}}
\newcommand{\B}{\mathsf{B}} 
\newcommand{\RR}{\mathbb{R}}
\newcommand{\X}{\mathcal{X}} 
\newcommand{\U}{\mathcal{U}}
\newcommand{\Z}{\mathscr{Z}}
\newcommand{\LL}{\mathbb{L}}
\newcommand{\OO}{\mathbb{O}}
\newcommand{\XX}{\mathbb{X}}
\newcommand{\YY}{\mathbb{Y}}
\newcommand{\ZZ}{\mathbb{Z}}
\newcommand{\UU}{\mathbb{U}}
\newcommand{\VV}{\mathbb{V}}
\newcommand{\TT}{\mathbb{T}}
\newcommand{\T}{\mathcal{T}} 

\newcommand{\cN}{\mathcal{N}}

\newcommand{\Xbar}{\overline{X}}
\newcommand{\xbar}{\overline{x}}
\newcommand{\Ubar}{\overline{U}}
\newcommand{\ubar}{\overline{u}}

\newcommand{\bigmid}{\; \Bigl\vert \;}

\renewcommand{\L}{\mathcal{L}}
\newcommand{\C}{\mathcal{C}}

\renewcommand{\S}{\mathcal{S}}
\renewcommand{\SS}{\mathbb{S}}

\newcommand{\del}{\partial}
%\renewcommand{\phi}{\varphi} 
\newcommand{\eps}{\varepsilon}

\newcommand{\iid}{\overset{\text{\tiny iid}}{\,\sim\,}}
\newcommand{\ind}{\overset{\text{\tiny ind}}{\,\sim\,}}


\newcommand{\model}{\mathscr{P}}
\newcommand{\prior}{\mathsf{Q}}
\newcommand{\credal}{\mathscr{Q}}

\newcommand{\pbcred}{\mathscr{R}}
\newcommand{\pb}{\mathsf{R}}

\newcommand{\cred}{\mathscr{C}}

\newcommand{\lPi}{\underline{\Pi}}
\newcommand{\uPi}{\overline{\Pi}}

\newcommand{\lChi}{\underline{\mathrm X}}
\newcommand{\uChi}{\overline{\mathrm X}}

\newcommand{\lUpsilon}{\underline{\Upsilon}}
\newcommand{\uUpsilon}{\overline{\Upsilon}}

\newcommand{\lPhi}{\underline{\Phi}}
\newcommand{\uPhi}{\overline{\Phi}}

\newcommand{\lprior}{\underline{\mathsf{Q}}}
\newcommand{\uprior}{\overline{\mathsf{Q}}}


\newcommand{\lprob}{\underline{\mathsf{P}}}
\newcommand{\uprob}{\overline{\mathsf{P}}}
\newcommand{\uE}{\overline{\mathsf{E}}}

\newcommand{\uuprob}{\mathbf{\overline{P}}}

\newcommand{\ee}{\text{\sc e}}
\newcommand{\loss}{\text{\sc loss}}

\newcommand{\eval}{\mathfrak{e}}


\newcommand{\rmnote}[1]{\marginpar{\color{red} \scriptsize #1}}


\title{Regularized e-processes: anytime valid inference with knowledge-based efficiency gains}
\author{Ryan Martin\footnote{Department of Statistics, North Carolina State University, {\tt rgmarti3@ncsu.edu}}
}
\date{\today}


\begin{document}

\maketitle 

\begin{abstract}
Classical statistical methods have theoretical justification when the sample size is predetermined.  In applications, however, it's often the case that sample sizes are data-dependent rather than predetermined.  The aforementioned methods aren't reliable in this latter case, hence the recent interest in e-processes and methods that are anytime valid, i.e., reliable for any dynamic data-collection plan.  But if the investigator has relevant-yet-incomplete prior information about the quantity of interest, then there's an opportunity for efficiency gain.  This paper proposes a {\em regularized e-process} framework featuring a knowledge-based, imprecise-probabilistic regularization with improved efficiency.  A generalized version of Ville's inequality is established, ensuring that inference based on the regularized e-process are anytime valid in a novel, knowledge-dependent sense.  Regularized e-processes also facilitate possibility-theoretic uncertainty quantification with strong frequentist-like calibration properties and other Bayesian-like properties: satisfies the likelihood principle, avoids sure-loss, and offers formal decision-making with reliability guarantees. 

\smallskip

\emph{Keywords and phrases:} Credal set; decision-making; e-value; imprecise probability; inferential model; possibility theory; safety; uncertainty quantification. 
\end{abstract}


%\tableofcontents 


\section{Introduction}
\label{S:intro}

Most statistical and machine learning methods in the literature have theoretical justification that assumes the sample size is (or can safely be treated as) fixed, say, by the experimental or data-collection protocol.  But in many applications this assumption is violated---e.g., decisions to stop/continue experimentation are made dynamically while monitoring the data---yet the same fixed-sample-size methods are used for data analysis.  Use of methods when they lack theoretical justification raises doubts about their reliability.  This lack of reliability is surely at least a contributor to the widely-publicized replication crisis in science \citep[e.g.,][]{nuzzo2014, baker.nature.2016, camerer2018}.  

To address this concern, there's been a surge of effort to develop {\em safe} or {\em anytime valid} statistical methods, i.e., methods that can preserve their reliability even when applied in cases where the data-collection process is dynamic; see the recent survey by \citet{evalues.review} and Section~\ref{SS:eprocess} below.  At the heart of these new developments are {\em e-processes}, which satisfy a time-uniform probability bound, namely, {\em Ville's inequality}, that can be leveraged to construct procedures with error rate control guarantees that hold no matter how the data-collection process is stopped.  
%There are some guiding principles behind the construction of an e-process, and numerous different instantiations are now available in the literature (Section~\ref{SS:eprocess}).  
Different anytime valid methods are typically compared in terms of their efficiency, i.e., power of tests, size of confidence intervals, etc.  In terms of the e-process itself, efficiency corresponds to how quickly it grows when evaluated at the ``wrong'' hypotheses: fast growth of the e-process is desirable because it means that what's ``wrong'' can be detected sooner, hence practitioners can make safe and reliable decisions with less time/exposure and fewer resources.   

But e-processes have limits to how fast they can grow \citep[e.g.,][]{grunwald.safetest}.  What can be done to improve efficiency when there are no more e-process tweaks to be made? It's now statistical second-nature to leverage penalty functions or prior distributions for {\em regularization}, to encourage certain structure in estimates, usually for the sake of efficiency gain.  Regularization is relatively easy when the goal is asymptotic consistency or calibration, since many different regularization strategies work in that sense. But anytime validity is a finite-sample property, so virtually any tweak made to an e-process would jeopardize the anytime validity property that motivated its use in the first place.  

This paper considers the situation in which the investigator {\em knows something} about $\Theta$, the relevant quantity of interest, before the data are available for analysis.  To keep the names and roles clear, ``I'' refers to me, this paper's author, and ``You'' refers to the investigator, a catch-all term for the individual(s) carrying out a study, analyzing data, etc.  The simple phrase {\em knows something} is rather nuanced.  On the one hand, I mean {\em know} in the strong sense of absolute certainty whereas, on the other hand, the {\em something} that's known could be so substantial that it renders data irrelevant, so miniscule that it means practically nothing, or somewhere in the middle of these two extremes.  
%(To be clear, ignorance is a special type of knowledge, and it's covered by what I'll present herein.) 
According to \citet{levi1980}, Your {\em corpus of knowledge} consists of the (lower and upper) probabilities that You'd assign to hypotheses about the uncertain $\Theta$ based on the information available to You at the relevant moment in time.  
%Later I'll give a mathematical description of Your corpus of knowledge.  
The point I want to emphasize is that Your corpus of knowledge is what it is at the time You call on it; without new information that warrants revising Your corpus, there's no justification for adding to or subtracting from it. Modern data science, however, focuses on methods that can be applied off-the-shelf which, by design, are incapable of accommodating Your prior knowledge and, therefore, expect You to subtract from Your corpus of knowledge. 
%modern data science research focuses on automation: developing tools that You can directly apply to reveal insights in Your data.  Since these automated methods can't accommodate what You know {\em a priori}, they implicitly expect You to subtract from Your corpus of knowledge.  
The present paper shows how You can use precisely what You know and retain statistical reliability.  
%Some might argue that this subtraction is warranted because Your subjective judgments have no place in scientific investigations.  At the end of the day, however, You'll inevitably apply what You {\em know} to Your problem in some way, and rightfully so.  Without statistical guidance on how to properly express and incorporate this knowledge, it may be done improperly, thereby creating systematic risk of errors, or at least doubts about reliability.  This paper aims to offer such guidance.

%The previous paragraphs betray my subjectivist perspective, but what I'm suggesting here is different from subjective Bayes.    
%The typical criticism of subjectivism is that it's focused more on what's ``right'' in principle than what ``works'' in practice.
%\footnote{For example, \citet{speed.principles} writes: ``If  a  statistical analysis is clearly shown to be effective [...], it gains  nothing from being described as principled.''}  
%That is, subjective Bayes has strong justification based on core principles, but practical obstacles often get in the way of achieving this principled ideal.  By abandoning those principles, however, practitioners can often find simpler methods that give reasonable answers and can be justified based on their effectiveness.  The apparent gap between what's ``right'' and what ``works'' is troubling, so I'm working to close it.  
%My proposal first converts Your corpus of knowledge into what I'm calling a {\em regularizer}, and then combines this regularizer with Your favorite e-process to get a new knowledge-aware e-process for reliable and efficient inference.  But since ignorance is a special type of knowledge---one where You're completely unsure about everything---my proposal isn't strictly for subjectivists.  
%In fact, the results presented herein are general and cover both the ignorant and non-ignorant cases alike.  
%For that reason, some of the developments below (Section~\ref{S:reg.uq} in particular) can be interpreted as contributions to e-process theory.  
%To be clear, I'm not asking You to concoct a non-trivial corpus of knowledge before You can follow my proposal.  Instead, I'm offering an opportunity to use what You already know---which could be nothing or something---to improve upon well-established statistical methods.

The key insight is leveraging Your corpus of knowledge in two new and distinct ways.  The first is {\em regularization}, which amounts to directly manipulating the data-driven e-process using what You know.  Unless You're ignorant, Your corpus of knowledge casts a degree of doubt on some values $\theta$ of $\Theta$.  The regularizer appropriately discounts those $\theta$ values, thereby boosting the regularized e-process there.  A bigger e-process means more $\theta$ values can be excluded from consideration, hence more efficiency.  The second is introducing a generalized, {\em regularization-aware} notion of anytime validity with respect to which the regularized e-process is evaluated.  Since You can't doubt what You know, You'd be willing to evaluate the performance of Your regularized e-process based on a metric that depends on what You know.  Roughly, my proposal requires that the regularized e-process be anytime valid with respect to any $(\text{data}, \Theta)$-joint distribution compatible with both the data-generating model and Your corpus of knowledge., the latter being expressed in terms of a credal set of ``priors'' for $\Theta$.  The standard definition of anytime validity corresponds to the special---and most restrictive---case where the aforementioned credal set contains all priors for $\Theta$.  By moving away from this most-restrictive case, regularization allows for efficiency gains without ruining anytime validity. 

The chief technical development in the paper is the establishment of a generalized version of Ville's inequality for my proposed regularized e-process, one that depends in a certain way on Your corpus of knowledge.  This leads to provably reliable, non-asymptotic inference with (or without) regularization and opportunities for efficiency gain.  Beyond construction of test and confidence procedures, a full-blown framework for (possibilistic) uncertainty quantification, which facilitates both calibrated, data-dependent (imprecise) probabilistic reasoning about $\Theta$, and reliable, data-driven decision-making.  

% IMPORTANT BUT REMOVED:
%Is this knowledge-based, regularization-aware notion of anytime validity satisfactory justification for the use of Your more-efficient, regularized e-process?  That it's satisfactory to You is automatic---it's based on {\em Your} corpus of knowledge.  Whether the justification is satisfactory to others depends on if Your corpus of knowledge can withstand their scrutiny.  This is a matter for epistemologists, pertaining to questions about justified beliefs \citep[e.g.,][]{goldman1979, ramsey.knowledge} and justified credence \citep[e.g.,][]{dunn2015, tang2016, pettigrew2021}, so here is not the place to pursue this.  My point is that harnessing the full power of data science is a collaborative effort between statisticians and subject matter experts, and my proposed knowledge-based regularized e-process framework aims to exemplify this relationship.  Statisticians provide the best-available e-process and other technical know-how, subject matter experts provide the application-specific ``justified credence,'' and together they (mathematically and figuratively) lift each other up to contribute more than each could individually.  

The organization of the paper is as follows.  After some background in Section~\ref{S:background} on e-processes and imprecise probability, Section~\ref{S:reg.eprocess} incorporates Your corpus of knowledge into the data-driven, e-process-based analysis via regularization,   
%It turns out that the so-called regularizer itself closely resembles an e-process, so the proposed combination is rather straightforward.  
%After a concrete regularizer construction and some further details, 
%In particular, Sections~\ref{SS:reg.ville}--\ref{SS:implications} present the 
establishes an imprecise-probabilistic generalization of Ville's inequality, and investigates its statistical implications.  Illustrations of the possible efficiency gains are presented in Section~\ref{SS:gains}.  Focus shifts in Section~\ref{S:reg.uq} to broader uncertainty quantification about $\Theta$.  What I'm proposing is a brand of {\em inferential model} (IM), where the familiar probabilistic reasoning is replaced by provably reliable possibilistic reasoning. My proposed e-possibilistic IM framework inherits the anytime validity from the regularized e-process, which implies that my uncertainty quantification is calibrated in a strong sense, hence reliable.  Remarkably, in addition to these frequentist-like calibration properties, the e-possibilistic IM satisfies several other desirable Bayesian-like properties: satisfying the likelihood principle, avoiding sure-loss, and formal decision-making with strong reliability guarantees.  An application analyzing clinical trial data is presented in Section~\ref{S:ware}, illustrating how subject matter knowledge gets translated into a regularizer.  Concluding remarks are made in Section~\ref{S:discuss}, and further technical details are given in the supplementary materials. 




% Levi (and Peirce) talks about beliefs in the corpus of knowledge at any given time, and emphasize that there is no doubt about the truthfulness/validity of these beliefs in the agent's mind -- otherwise, it wouldn't be included in the corpus... The prior information I'm talking about here has this interpretation... 



% [Re: possibility is a very special case of IP]  It's the same with ordinary/precise probability: there are many different brands of probability \citep[e.g.,][]{fine.book} but most statisticians exclusively use Kolomogorov's version based on density/mass functions, which is the least expressive but also the simplest and most immediate to apply... It's all about striking the appropriate balance between expressibility and practicality... I hope to make a case that possibility theory strikes this balance in the context of (regularized) statistical inference... 


% There's more to science (and what statistics can contribute to science) than purely data-driven procedures.






\section{Background}
\label{S:background}

\subsection{Setup and notation}
\label{SS:setup}

Start with a baseline probability space $(\SS, \A, \prob)$, where $\A$ is a $\sigma$-algebra of subsets of $\SS$ and $\prob$ is a probability measure.  Let $Z: \SS \to \ZZ$ be a measurable function that takes values in a topological space $\ZZ$.  
%The goal is to reliably answer certain questions about the uncertain $\prob$, or relevant features thereof, based on one or more data point $Z$.  
I'll also write $\prob$ for the induced distribution of $Z$.  

Next, let $Z_1, Z_2, \ldots$ denote independent and identically distributed (iid) copies of $Z$ and, for each $n \geq 1$, write $Z^n = (Z_1,\ldots,Z_n)$.  
%The iid assumption is mostly for convenience of notation; the independent-but-not-iid case is a straightforward extension discussed, among other things, in Section~\ref{S:extensions}.  
%The iid assumption isn't necessary, but the advantage is that the joint distribution of $Z^\infty = (Z_1, Z_2, \ldots)$ is fully determined by the same $\prob$ introduced above.  
Define the filtration $\A_n$ to be the sequence of $\sigma$-algebras determined by the information in  $Z^n$.  A stopping time $N$ is a positive integer-valued random variable such that, for each $n$, the event $\{N=n\}$ is in $\A_n$.  
%This means that $N$'s value is determined by the process's history.  

%{\color{red} Define a set $\cN$ of candidate stopping times...? Maybe need to restrict to a manageable subcollection of stopping times that are, e.g., ``computable'' in Turing's/Dawid's sense...?} 

For the statistical applications that I have in mind, $\prob$ is {\em uncertain}.  
%I'll be more specific about what  ``uncertain'' means in Section~\ref{SS:ip}; here I'm primarily focused on notation.  
To facilitate this notion of an ``uncertain $\prob$,'' it'll help to introduce (the notation of) a model, namely, $\model = \{\prob_\omega: \omega \in \OO\}$, indexed by $\OO$.  Of course, this could be a familiar parametric model, but it could also be that $\OO$ is in one-to-one correspondence with the set of all relevant probability distributions, so there's no loss of generality in introducing the index $\omega$.  
%The advantage is that it's easier to describe mathematically an ``uncertain $\prob$'' by a more familiar type of uncertain variable $\Omega$ in the space $\OO$.  In any case, again, the transition from ``uncertain $\prob$'' to ``uncertain $\Omega$'' is without loss of generality.  

The goal then is to quantify uncertainty about the uncertain $\prob$ or, equivalently, about the uncertain index $\Omega$, based on observations $Z_1,Z_2,\ldots$ from the underlying process that depends on $\prob$ or $\Omega$.  It's often the case that only certain features of the uncertain ($\prob$ or) $\Omega$ are relevant to the application at hand---e.g., maybe one only needs to know about the mean survival time of patients---so I'll define this relevant feature as $\Theta = f(\Omega)$, taking values in the possibility space $\TT = f(\OO)$.  
%To be clear, this feature-extracting function $f: \OO \to \TT$ is fixed by the context of the problem; it's only the input $\Omega$ that's uncertain, which makes the output $\Theta$ also uncertain.  
Since $f$ could be the identity function, in which case $\Theta = \Omega$, this focus on quantifying uncertainty about $\Theta$ based on data $Z_1,Z_2,\ldots$ is without loss of generality.  

One mild technical condition I'll impose is that the topology on $\OO$ is sufficiently rich that {\em $f$ is continuous}.  Continuity isn't necessary for the developments here (see Remark~\ref{re:continuity} in Appendix~\ref{A:remarks}), but it's no serious practical constraint and it greatly simplifies the developments in Section~\ref{SS:pullback}.  
%there's also a relatively naive ways to bypass a continuity assumption, but this leads to some efficiency loss so I won't discuss it here.  Besides, 
For parametric statistical models, the relevant quantity would either be the parameter itself, or some interpretable feature thereof, e.g., one component of a parameter vector, so the function is almost always trivially continuous.  When no parametric model is assumed, as is often the case in machine learning applications, the relevant features $\Theta = f(\Omega)$ are often risk minimizers.  Let $L_t: \ZZ \to \RR$ be a loss function indexed by $t \in \TT$ and define the mapping $(\omega, t) \mapsto \prob_\omega L_t$, the expected loss corresponding to $t$ at $\omega$.  If this mapping is continuous, then Berge's {\em maximum theorem} \citep[e.g.,][Theorem~17.31]{infinite.dim.analysis} implies that $f(\omega) := \arg\min_t \prob_\omega L_t$ is continuous.  



% https://en.wikipedia.org/wiki/Maximum_theorem



\subsection{E-processes}
\label{SS:eprocess}

Start by fixing a particular $\omega \in \OO$.  A sequence $(M^n: n \geq 0)$ is a {\em supermartingale}, relative to $\prob_\omega$ and the filtration $(\A_n)$, if $\E_\omega(M^n \mid \A_{n-1}) \leq M^{n-1}$, where $\E_\omega$ denotes expected value with respect to $\prob_\omega$.  I have in mind a function $M(\cdot)$ that maps $\ZZ$-valued sequences to numbers, and $M^n = M(Z^n)$ for each $n$, with $M^0 = M(\square)$ the value of $M$ when applied to the empty sequence $\square$.  An $\omega$-dependent supermartingale $(M_\omega^n: n \geq 0)$ is a {\em test supermartingale} for $\omega$ if it's non-negative and $M_\omega^0 \equiv 1$ \citep[e.g.,][]{shafer.vovk.martingale}.  A sequence $(\eval_\omega^n: n \geq 0)$ is an {\em e-process} for $\omega$ if it's non-negative and upper bounded by a test supermartingale for $\omega$; see \citet{ramdas.nnm} and \citet{evalues.review} for details.  Again, $(\eval_\omega^n)$ is determined by a mapping $\eval_\omega(\cdot)$, i.e., $\eval_\omega^n = \eval_\omega(Z^n)$.  Test supermartingales and, hence, e-processes satisfy two key properties: the first \citep[e.g.,][Theorem~5.7.6]{durrett2010} is 
\[ \E_\omega ( M_\omega^N ) \leq 1 \quad \text{all $\omega \in \OO$, all stopping times $N$}, \]
and the second, known as {\em Ville's inequality} \citep[e.g.,][]{shafer.vovk.book.2019}, is
\[ \prob_\omega(M_\omega^N \geq \alpha^{-1}) \leq \alpha \quad \text{all $\alpha \in (0,1]$, all stopping times $N$}. \]
Results of this type have important statistical implications \citep[e.g.,][]{shafer.betting}.  One is that, if $(\eval_\omega^n)$ is an e-process for $\omega$, then the test that rejects the hypothesis ``$\Omega=\omega$'' based on data $Z^n$ if and only if $\eval_\omega(Z^n) \geq \alpha^{-1}$ controls the frequentist Type~I error at level $\alpha$, regardless of what stopping rule might be used to terminate the data-collection process.  This {\em anytime validity}, or {\em safety} \citep{grunwald.safetest}, of the e-process-based tests is a major advancement beyond the classical tests that are valid only for fixed $n$. Of course, if one has a collection of e-processes $(\eval_\omega^n)$, one for each $\omega \in \OO$, then the above testing procedure can be inverted to construct a confidence set for $\Omega$ which inherits the same anytime validity property: the frequentist coverage probability is no less than the nominal level independent of the choice of stopping rule. 

With one exception, my examples below use Savage--Dickey e-processes, i.e., Bayes factors that rely on ``default'' priors.  Such constructions are quite general, but difficulties can arise in more complex problems.  My proposal described in the following sections can be applied to any e-process construction, including those in \citet{evalues.review}.  

A practically important extension of the ideas presented above is to the case of composite hypotheses.  Even the simple-looking hypotheses I have in mind here are composite---that is, ``$\Theta=\theta$'' generally corresponds to a set of $\omega$ values for $\Omega$.
%, i.e., $f^{-1}(\{\theta\}) = \{ \omega \in \OO: f(\omega) = \theta \}$.  
Fortunately, there's an easy way to accommodate this more general case under certain measurability constraints: if  $(\eval_\omega^n)$ is an e-process for each $\omega \in \OO$, then, with a minor abuse of notation, the following is an e-process for $\theta$,
\[ \eval_\theta^n = \inf_{\omega \in \OO: f(\omega) = \theta} \eval_\omega^n, \quad \theta \in \TT. \]
Ignoring issues concerning the potential non-measurability of the above infima,
%\footnote{If it happens that the infimum is not measurable, then one can replace the expectations and probabilities in \eqref{eq:eval.bound} and \eqref{eq:ville} with outer expectation and outer probability, respectively.  For reasons that will become clear shortly, I don't want to deal with this technical complication here.} 
it follows immediately from the properties discussed above that 
\begin{equation}
\label{eq:eval.bound}
\sup_{\omega: f(\omega) = \theta} \E_\omega \{\eval_\theta(Z^N) \} \leq 1 \quad \text{all $\theta \in \TT$, all stopping times $N$}. 
\end{equation}
There's also a corresponding version of Ville's inequality:
%, from which the anytime validity of the corresponding tests and confidence sets for $\Theta$ are derived: 
\begin{equation}
\label{eq:ville}
\sup_{\omega \in \OO: f(\omega) = \theta} \prob_\omega\bigl\{ \eval_\theta(Z^N) \geq \alpha^{-1} \bigr\} \leq \alpha \quad \text{all $\alpha \in (0,1]$, all stopping times $N$}. 
\end{equation}
I should mention that constructing a e-process for $\Omega$ and then marginalizing to $\Theta$ via optimization as just described is not the only way.  Indeed, one can construct e-processes for $\Theta$ directly, bypassing $\Omega$ altogether; see \citet{dey.gue} and Section~\ref{SS:gue} below. 
%but one I'd like to mention specifically is {\em universal inference}, first put forward in \citet{wasserman.universal}.  Universal e-processes typically correspond to likelihood ratios where one of the likelihoods is ``cautiously maximized,'' following a general idea goes back at least to \citet[][Eq.~10.10]{wald.sequential}.  For recent applications to nonparametric problems, see \citet{gangrade2023sequential} and \citet{preprocess}.  A variation for when there's no likelihood is proposed in \citet{dey.gue}.  




\subsection{Imprecise probability}
\label{SS:ip}

%Certain aspects of imprecise probability theory will play a key role in this paper, so here I introduce the essential notation and concepts.  The goal, recall, is to quantify uncertainty about $\Theta$, to give a quantitative expression of Your corpus of knowledge.  
%For now, to keep things simple, I'll refer directly to $\Theta$, ignoring the (important) connection between $\Theta$ and $\Omega$ and/or $\prob$.  

The most familiar approach to quantifying uncertainty about $\Theta$ is to introduce a probability measure $\prior$, supported on subsets of $\TT$ and, then, for any relevant hypothesis ``$\Theta \in H$,'' where $H \subseteq \TT$, Your uncertainty about its truthfulness is quantified by $\prior(H)$, the $\prior$-probability of $H$. (Henceforth I'll refer to both $H$ and ``$\Theta \in H$'' as hypotheses about $\Theta$.) But it's easy to imagine having limited information and, hence, You can't precisely state the probability of $H$ for each hypothesis $H$.  This includes the extreme case of {\em vacuous} information, where You literally know nothing about $\Theta$, so all You can say is that $\prior(H) \in [0,1]$ for all $H \not\in \{\varnothing, \TT\}$.  More generally, You might be in between the two extremes of knowing nothing and knowing $\prior$ precisely.  For example, suppose ``You're 95\% sure that $\Theta$ is bigger than 7.''  This information imposes bounds on $\prior(H)$ for some $H$, e.g., $\prior(H) \leq 0.05$ for all $H \subseteq (-\infty, 7]$, which are satisfied by many different $\prior$, so this information fails to identify a single $\prior$.  This ambiguity alone isn't grounds for You to ignore the information altogether.

When the information available about $\Theta$ determines a single $\prior$, then Your uncertainty quantification is {\em precise}; otherwise, it's {\em imprecise}.  Mathematically, the latter case can be handled using {\em imprecise probabilities}.  I must emphasize that precise probability isn't superior to imprecise probability.  Your information about $\Theta$ is what it is, so artificially embellishing on that information to avoid addressing the inherent imprecision doesn't make Your uncertainty quantification better.  The goal is to represent Your corpus of knowledge as faithfully as possible; if that requires imprecision, then so be it. 
%; but see Note~\ref{no:lindley}.

A fairly general way to define imprecise probabilities is via a {\em credal set} $\credal$, i.e., a non-empty, closed, and convex collection of probabilities $\prior$ on $\TT$ \citep[e.g.,][Sec.~4.2]{levi1980}.  Intuitively, $\credal$ encodes Your knowledge about $\Theta$, so the collection $\credal(H) = \{ \prior(H): \prior \in \credal \}$ of credences quantifies uncertainty about the truthfulness of $H$.  Natural summaries of this collection include the lower and upper bounds:
\begin{equation}
\label{eq:lower.upper}
\lprior(H) = \inf\{ \prior(H): \prior \in \credal \} \quad \text{and} \quad \uprior(H) = \sup\{ \prior(H): \prior \in \credal \}. 
\end{equation}
Thanks to the structure in $\credal$, the lower and upper probabilities are linked together, 
%through a relationship called {\em conjugacy}, 
i.e., $\lprior(H) = 1 - \uprior(H^c)$ for all $H$.  For instance, if ``You're 95\% sure that $\Theta$ exceeds 7'' case above, then You get $\lprior(H)=0$ and $\uprior(H) = 0.05$ for all $H \subseteq (-\infty, 7]$, $\lprior(H)=0.95$ for all $H$ with $H \subseteq (7,\infty)$, and $\uprior(H) = 1$ for all $H$ with $H \cap (7,\infty) \neq \varnothing$.  In what follows, without loss of generality, I'll focus almost exclusively on upper probabilities.  De Finetti-style betting interpretations can be given to $\lprior$ and $\uprior$ (see  Remark~\ref{re:betting} in Appendix~\ref{A:remarks}), but this isn't necessary.  You can follow \citet{shafer1976} and others by interpreting $\lprior(H)$ and $\uprior(H)$ simply as Your quantitative degree of confidence/support and degree of plausibility, respectively, in the truthfulness of $H$. 

%{\color{red} Elements of the credal set are countably-additive or only finitely-additive...?} 

In precise probability theory, probabilities are extended to expected values via Lebesgue integration.  In imprecise probability theory, there's an analogous extension to lower and upper expected values, or lower and upper previsions \citep[e.g.,][]{walley1991, lower.previsions.book}.  Given $g: \TT \to \RR$, the lower and upper expected value is defined as 
\begin{equation}
\label{eq:envelope}
\lprior \, g = \inf_{\prior \in \credal} \E^{\Theta \sim \prior} \{g(\Theta)\} \quad \text{and} \quad \uprior \, g = \sup_{\prior \in \credal} \E^{\Theta \sim \prior} \{ g(\Theta) \}, 
\end{equation}
i.e., the lower and upper limits of the usual expected values of $g(\Theta)$ with respect to probability distributions $\prior$ in the credal set $\credal$.  
%Note that lower and upper expectation of indicator functions returns the original lower and upper probabilities.  
According to the definition above, computation of the lower and upper expectations appears to require non-trivial optimization over a potentially infinite-dimensional space $\credal$.  In certain cases, however, there are equivalent formulations with computationally simpler formulas; see, e.g., Equation~\eqref{eq:upper.e.poss} below and the general discussion of {\em Choquet integration} in Appendix~\ref{A:choquet}.   

A special case of imprecise probability is {\em possibility theory}, 
%dating back at least to \citet{shackle1961} and 
closely tied to fuzzy set theory \citep[e.g.,][]{zadeh1978}.  The seminal text on possibility theory is \citet{dubois.prade.book}; see, also
%Dominik Hose's PhD thesis 
\citet{hose2022thesis} for a review and, e.g., \citet{dubois2006} and \citet{imchar, martin.partial2} for connections to statistics.  What distinguishes possibility theory from other brands of imprecise probability is its simplicity.  Indeed, a possibility measure is determined by a real-valued function, analogous to how precise probabilities are determined by a probability density.  The difference is that, while probability theory uses integration, possibility theory is based on optimization.  Define a contour function $q: \TT \to [0,1]$ such that $\sup_{\theta \in \TT} q(\theta) = 1$, the condition analogous to a probability density function integrating to 1.  Then the corresponding possibility measure is the upper probability $\uprior$ defined via optimization:
\begin{equation}
\label{eq:upper.p.poss}
\uprior(H) = \sup_{\theta \in H} q(\theta), \quad H \subseteq \TT.
\end{equation}
%and, as above, the lower probability is defined via conjugacy, $\lprior(H) = 1 - \uprior(H^c)$.  
A simple illustration of a possibility contour $q$ and of the formula \eqref{eq:upper.p.poss} is shown in Figure~\ref{fig:toy}.  And the ``You're 95\% sure that $\Theta$ exceeds 7'' case can easily be encoded by a possibility measure $\uprior$ as in \eqref{eq:upper.p.poss} determined by the contour $q(\theta) = 0.05 \times 1(\theta \leq 7) + 1(\theta > 7)$.
%; this makes clear that the $\uprior(H)$ is determined by $H$'s ``most possible'' element.  
%In any case, Your choice of how specifically to construct and evaluate the lower and upper probabilities needn't affect its interpretation, so all that I said above applies here.  

\begin{figure}[t]
\begin{center}
\scalebox{0.6}{\includegraphics{toy}}
\end{center}
\caption{A possibility contour $q$, the hypothesis of interest $H=[3,5]$ (red), and the corresponding upper probability $\uprior(H)$ determined by optimization as in \eqref{eq:upper.p.poss}.}
\label{fig:toy}
\end{figure}

Two possibility-theoretic details deserve mention.  The first is a simple formula for $\uprior$'s extension to an upper expected value.  
%Since I won't need to consider negative functions in this paper, I'll focus specifically on the non-negative case.  
If $g: \TT \to \RR$ is a suitable non-negative function and $\uprior$ a possibility measure determined by contour $q$ as described above, then the corresponding possibilistic upper expected value (see Appendix~\ref{A:choquet})
%\ref{A:choquet} 
is given by
\begin{equation}
\label{eq:upper.e.poss}
\uprior \, g = \int_0^1 \Bigl\{ \sup_{\theta: q(\theta) > s} g(\theta) \Bigr\} \, ds.
\end{equation}
%The key observation is that, by definition, the left-hand side of \eqref{eq:upper.e.poss} is the supremum of ordinary $\prior$-expected values over an infinite-dimensional $\credal$, whereas the right-hand side of \eqref{eq:upper.e.poss} is a one-dimensional (Riemann) integral of a $\text{dim}(\TT)$-dimensional supremum.  
The second is a simple characterization \citep{cuoso.etal.2001, dubois.etal.2004} of the credal set $\credal$ associated with the possibility measure $\uprior$ determined by contour $q$:
%Previously, before we had any concrete imprecise probability formulation in mind, the credal set was just an abstract construct.  But with the possibility-theoretic structure laid out above, it's now possible to make the credal set more tangible, less abstract.  
\begin{equation}
\label{eq:credal.char}
\prior \in \credal \iff \prior\{ q(\Theta) \leq \alpha \} \leq \alpha \;\; \text{for all $\alpha \in [0,1]$}. 
\end{equation}
The reader may notice the similarity to p-values: $\prior$ belongs to the credal set if and only if $q(\Theta)$ is stochastically no smaller than $\unif(0,1)$ as a function of $\Theta \sim \prior$.  
%This characterization has other consequences---some good and some (often perceived to be) bad---but I'll postpone discussion of these points to {\color{red} where...?}.  



%even some that are often perceived as ``shortcomings'' of possibilistic uncertainty quantification.  In particular, since the contour $q$ is fixed, the probability distributions $\prior$ in $\credal$ are ``no more diffuse'' than $q$ in a certain sense.   This implies the following relationships:
%\[ \lprior(H) > 0 \implies \uprior(H) = 1 \quad \text{and} \quad \uprior(H) < 1 \implies \lprior(H) = 0. \]
%That is, for every $H$, at most one of the quantities in the pair $\{\lprior(H), \uprior(H)\}$ is a non-extreme value.  {\color{red} More...} 


% In the hypothetical/conceptual limit where prior information increases indefinitely, convergence would be to knowledge of the THE VALUE OF THETA, not to knowledge of THE DISTRIBUTION OF THETA... 











\section{Regularized e-processes}
\label{S:reg.eprocess}

%\subsection{Objective}
%
%This section develops a framework of {\em regularized e-processes} that offers knowledge-based efficiency gains compared the basic e-processes surveyed in Section~\ref{SS:eprocess} above.  The point is that one generally can't make non-trivial changes to an existing e-process without jeopardizing its anytime validity properties, so care is needed.  My strategy is to relax the definition of ``anytime valid'' in a principled, knowledge-aware manner that allows for greater flexibility and, in turn, creates opportunities for improved efficiency.  
%
%Recall that regularization is predicated on the availability of genuine, relevant information about the quantities of interest.  
%%Without this, it's unclear (a)~what kind of structure/simplification the regularization is aiming to promote, and (b)~why any chosen form of regularization is meaningful/justifiable.  
%%For example, in the now-standard ridge and lasso regression, the motivation behind shrinkage towards this origin comes from subject-matter background information that says most of the regression coefficients are small.  
%%Some might counter my previous claims by referring to the bias--variance trade-off that ridge/lasso offers, but there's no motivation for introducing the bias resulting from shrinkage towards the origin if one didn't have some reason to believe that the origin was a distinctly special, {\em a priori} plausible value.  
%My proposal for knowledge-based e-process regularization is, roughly, to make use of exactly the information that's available in Your corpus of knowledge, no more and no less, and then combine this with an existing e-process such that the combination tends to be larger and yet is still anytime valid in the relaxed sense.  More specifically, my proposal is two-fold: 
%\begin{enumerate}
%\item Encode Your corpus of knowledge, which is initially expressed as an imprecise probability as in Section~\ref{SS:ip}, in the form of a {\em regularizer} and then (easily) combine it with any data-driven e-process as discussed in Section~\ref{SS:eprocess} above, and 
%%I propose an imprecise-probabilistic quantification of the available prior information that can readily combine with any existing data-driven e-process like those reviewed in Section~\ref{SS:eprocess} above, and
%\vspace{-2mm}
%\item Instead of requiring e-processes to be anytime valid relative to a vacuous corpus of knowledge, I'll ask for a weaker version that's relative to the Your prior knowledge, and I'll offer a generalization of Ville's inequality that's satisfied by a collection of e-processes that tend to be larger, hence more efficient, than those in Section~\ref{SS:eprocess}.
%\end{enumerate}   



\subsection{Prior information and regularization}

Suppose Your knowledge about $\Theta$ takes the form of a (possibly vacuous) imprecise probability $\uprior$.  Introduce $\rho: \TT \to [0,\infty]$ such that $\rho(\theta)$ represents the weight of evidence (relative to $\uprior$) against the basic proposition ``$\Theta=\theta$.''  This is different from familiar notions of statistical evidence because there's no data, but there are still connections.  \citet[][Ch.~6.7]{good.evidence.book}, for example, proves a now-well-known result (which he attributes to Turing) that, in certain cases, the expectation of the Bayes factor weighing evidence against a true proposition is equal to 1.  More generally, this ``expected value is $\leq 1$'' property is key to all the recent e-value developments; see \eqref{eq:eval.bound} above.  To be a proper {\em regularizer}, the function $\rho$ must be similarly bounded in (upper) expectation.

\begin{defn}
\label{def:regularizer}
Let $\uprior$ be an upper probability/prevision on $\TT$.  A function $\rho: \TT \to [0,\infty]$ is a regularizer if its $\uprior$-upper expected value is bounded by 1, i.e., if $\uprior \rho \leq 1$. 
\end{defn}

%More concrete explanations will be given shortly; see, also, Section~\ref{SS:possibilistic.case} and {\color{red}Appendix~\ref{A:other}}.  

Note that any non-negative function $\rho$ with $0 \leq \rho \leq 1$ is a regularizer, but such a choice is useless.  Suppose, for example, that the prior is fully vacuous, i.e., if You literally know nothing about $\Theta$ {\em a priori}, which implies $\uprior\rho = \sup_{\theta \in \TT} \rho(\theta)$.  Then choosing $\rho \leq 1$ is the only choice that satisfies the condition in Definition~\ref{def:regularizer}.  So, a {\em trivial} choice of $\rho$ that's upper-bounded by 1 matches the no-prior-information case and, hence, can't offer any meaningful regularization.  Henceforth, except for the vacuous prior case where $\rho \equiv 1$ is the clear choice of regularizer, I'll focus on {\em non-trivial} regularizers, i.e., those that aren't upper-bounded by 1. 
%(A more explicit definition of ``non-trivial'' is needed only for the dominance result in Appendix~\ref{SS:product}.)  
Then the general goal is, for Your given $\uprior$, to find a function $\rho$ which can take values possibly much larger than 1 while still satisfying $\uprior\rho \leq 1$.  Naturally, for two regularizers $\rho$ and $\rho'$ with $\rho'(\theta) \leq \rho(\theta)$ for all $\theta$, the dominant regularizer $\rho$ is preferred.  A concrete transformation of prior information into a regularizer will be presented below; further details and examples are given in Appendix~\ref{A:other}. 

The quality of the regularizer is determined by how much of a knowledge-based boost it offers to an existing e-process $\eval_\theta(\cdot)$ for $\Theta$, as described in Section~\ref{SS:eprocess}.
%Just as a hop variety's value is determined by the tastiness of beer produced with it, the regularizer's value lies in the efficiency gains seen in the test and confidence set procedures derived from the resulting regularized e-process.  So, let $\eval_\theta(\cdot)$ denote an e-process as described in Section~\ref{SS:eprocess}, i.e., $\eval_\theta(\cdot)$ is a $\theta$-dependent function that maps sequences of $\ZZ$-valued random elements to sequences of non-negative random variables upper-bounded by a test martingale.  
My proposal is to define a {\em regularized e-process} by combining $\rho$ and $(\eval_\theta: \theta \in \TT)$ as follows: 
\begin{equation}
\label{eq:reg.eprocess}
\eval^\text{reg}(\cdot, \theta) = \rho(\theta) \times \eval_\theta(\cdot), \quad \theta \in \TT.
\end{equation}
%This regularized e-process is just a combination of a standard e-process $\eval_\theta(\cdot)$ and the e-process-like regularizer $\rho(\theta)$.  
The intuition is that, if $\theta$ is incompatible with prior knowledge, so that $\rho(\theta) \gg 1$, then $\eval^\text{reg}(\cdot,\theta) \gg \eval_\theta(\cdot)$, which creates an opportunity for increased efficiency.  But non-trivially manipulating the original e-process in an effort to boost efficiency will surely jeopardize anytime validity in the sense of Section~\ref{SS:eprocess}, hence care is needed.  In Section~\ref{SS:reg.ville} below, I'll establish my main claim: in a certain sense, the regularized e-process enjoys efficiency gains without jeopardizing anytime validity.  This ``certain sense'' involves what I'll argue is a meaningful imprecise-probabilistic relaxation of anytime validity as in Section~\ref{SS:eprocess}.

Finally, there are justifications for my choice to define the regularized e-process in \eqref{eq:reg.eprocess} via multiplication.  One is based on a Bayesian-like updating-coherence property and the other is based on a formal demonstration of dominance, similar to a result in \citet{vovk.wang.2021}; these details are presented in Appendix~\ref{SS:product}.
%Appendix~\ref{SS:product}



\subsection{Special case: possibilistic priors}
\label{SS:possibilistic.case}

Suppose Your prior knowledge about $\Theta$ is encoded by a possibility measure with contour function $q$; see Appendix~B for other options.  Then Your $\uprior$ is defined via optimization as in \eqref{eq:upper.p.poss}, and upper expectation is as in \eqref{eq:upper.e.poss}.  This is the simplest imprecise probability model, and simplicity is a virtue: Your prior knowledge about $\Theta$ is necessarily limited, so a more expressive imprecise probability model used to describe it makes elicitation more difficult.  The claim is that experts can offer statements like: ``I'd not be surprised at all if $\Theta=a$, I'd be a little surprised if $\Theta=b$, and I'd be very surprised if $\Theta=c$.'' ({\em Surprise} is due to \citealt{shackle1961}.)  Then the qualitative statement above could be made quantitative by introducing a possibility contour $q$ with $q(a)=1$, $q(b)$ smaller, $q(c)$ very small, etc.  This is often how penalties and priors are chosen.  For example, a Bayesian might take the least-surprising value to be the prior mode and then choose a density that's a decreasing function of surprise.  But there's a drastic difference between a probability density and a possibility contour, even if they have similar shapes: the former determines a precise probability thereby adding artificial information to Your corpus of knowledge.  

Recall that a useful regularizer $\rho$ must take values (potentially much) larger than 1 at $\theta$ values that are incompatible with the prior information.  In the present context, ``incompatibility'' corresponds to a small $q$ value, so $\rho(\theta)$ should be large when $q(\theta)$ is small.  As a first attempt, this could be accomplished by taking $\rho$ to be the reciprocal of $q$.  But $q$ is no more than 1, so such a $\rho$ would never be less than 1 and hence it can't be a regularizer in the sense of Definition~\ref{def:regularizer}.  Apparently $q$ is too small for the reciprocal to work, so it needs to be inflated first.  Following \citet[][Sec.~6]{shafer.vovk.martingale}, define a(n admissible) {\em calibrator}
%\footnote{The calibrators I'm referring to here are called {\em admissible calibrators} in \citet{shafer.vovk.martingale}. I'll have no need here for inadmissible calibrators with $\int_0^1 \gamma(u)^{-1} \, du < 1$.} 
as a non-decreasing function $\gamma: [0,1] \to (0, \infty]$ such that   
\begin{equation}
\label{eq:calibrate}
\int_0^1 \frac{1}{\gamma(s)} \, ds = 1. 
\end{equation} 
Of course, $\gamma$ can taken to be the reciprocal of any probability density function supported on $[0,1]$.  Then 
%the step needed before Your prior information and data-driven e-process can be combined is a calibration step.  
I'll define the {\em regularizer} in terms of the calibrated version of $q$: 
\begin{equation}
\label{eq:regularizer}
\rho(\theta) = \{ \gamma \circ q(\theta) \}^{-1}, \quad \theta \in \TT. 
\end{equation}
The role played by the calibrator is to inflate the contour ``just enough.'' 

\begin{prop}
\label{prop:prior.reg}
Let $\uprior$ be the possibility measure determined by the contour function $q$.  Then $\rho$ defined in \eqref{eq:regularizer} is a regularizer, i.e., has $\uprior$-upper expectation bounded by 1.
\end{prop}

%\begin{proof}
%By definition of the $\uprior$-upper expectation, 
%\[ \uprior \, \rho = \int_0^1 \Bigl\{ \sup_{\theta: q(\theta) > s} \rho(\theta) \Bigr\} \, ds = \int_0^1 \Bigl\{ \sup_{\theta: q(\theta) > s} \frac{1}{\gamma \circ q(\theta)} \Bigr\} \, ds. \]
%Since $\gamma$ is non-decreasing, 
%\[ q(\theta) > s \implies \frac{1}{\gamma \circ q(\theta)} \leq \frac{1}{\gamma(s)}, \]
%and from here the claim follows immediately from \eqref{eq:calibrate}.  
%\end{proof}

I'll follow \citet[][Appendix~B]{vovk.wang.2021} and suggest use of (the reciprocal of) a suitable beta mixture of beta density functions, which takes the form 
\[ \gamma(u) = \frac{u \, (-\log u)^{1 + \kappa}}{\kappa \times {\tt igamma}(-\log u, 1 + \kappa)}, \quad u \in [0,1], \quad \kappa > 0, \]
where ${\tt igamma}(z, \alpha) = \int_0^z t^{\alpha-1} \, e^{-t} \, dt$ is the incomplete gamma function. A particular advantage of the above calibrator is how rapidly it vanishes as $u \to 0$, which is directly related to how severely the corresponding regularizer penalizes those $\theta$ values incompatible with the prior information.  The gamma defined in the above display is the calibrator that I'll use for my illustrations in Section~\ref{SS:gains} and elsewhere.  

Although the context and form here are different, the message above is a familiar one to those who have experience with e-processes, etc.  Indeed, Equation~\eqref{eq:credal.char} implies that, roughly, $q(\Theta)$ is a p-value relative to $\uprior$.  And it's well-known \citep[e.g.,][]{sellkebayarriberger2001, vovk1993} that one must calibrate a p-value so that its reciprocal is an e-value.  

 



\subsection{Induced upper joint distributions}
\label{SS:pullback}

%A bit more background is needed before I can get into the regularized e-process's properties.  
Recall that $\Theta = f(\Omega) \in \TT$ is a function of $\Omega \in \OO$.  Since $\TT$ is defined as the image of $\OO$, the map $f: \OO \to \TT$ is trivially surjective.  There are applications in which $f$ would be a bijection, e.g., when $f$ is the identity mapping, so that $\Theta$ and $\Omega$ are equivalent in some sense.  
%This identity-mapping case might be relevant if $\{ \prob_\omega: \omega \in \OO \}$ were a parametric model, so that $\Omega$ in its entirety can be interpreted as the ``true'' state of the world.  
But it's typical that $\Theta = f(\Omega)$ is a lower-dimensional feature of $\Omega$.  Whether $f$ is or isn't a bijection becomes relevant when considering how knowledge about $\Theta$ translates to knowledge about the primitive $\Omega$ from which $\Theta=f(\Omega)$ is derived.  The {\em push-forward} operation that takes a probability for $\Omega$ to a corresponding probability for $\Theta=f(\Omega)$ is well-defined, whereas the reverse {\em pull-back} operation is not unique when $f$ is not a bijection.  Here I give only the necessary details concerning pull-backs. 
%Details concerning pull-back measures are somewhat technical, but these aren't particularly relevant here, so I give only the key points.
%\footnote{Throughout the paper, I'm working with non-singleton sets of probabilities.  So whatever non-uniqueness might emerge in the pull-back operation is easily accommodated by my set of pull-backs.}

%https://math.stackexchange.com/questions/4427017/explicit-expression-for-pullback-measure-via-a-surjective-mapping

%https://math.stackexchange.com/questions/3322870/pull-back-of-a-finite-measure

Given a probability distribution $\prior$ for $\Theta$ on $\TT$, the relevant question is under what conditions does there exist a corresponding probability distribution, say, $\pb$, for $\Omega$ on $\OO$ such that the distribution of $f(\Omega)$ is $\prior$.  The classical result of \citet[][Lemma~2.2]{varadarajan1963} establishes the existence of an $\pb$ corresponding to $\prior$ under very mild conditions.
%, easily met in the statistical applications I have in mind here.  
In fact, for a given $\prior$, there's a non-empty class $\pbcred_\prior$ of such distributions, i.e., 
\begin{equation}
\label{eq:R.Q}
\pbcred_\prior = \{ \pb: \text{if $\Omega \sim \pb$, then $f(\Omega) \sim \prior$} \}, 
\end{equation}
and, moreover, this class is easily shown to satisfy the properties of a credal set.  

\begin{prop}
\label{prop:pullback}
For any given $\prior$ and continuous $f: \OO \to \TT$, the collection $\pbcred_\prior$ in \eqref{eq:R.Q} is non-empty, convex, and closed with respect to the weak topology. 
\end{prop}

%\begin{proof}
%Take any fixed $\prior \in \credal$.  Non-emptiness of $\pbcred_\prior$ is a consequence of the classical results mentioned above.  To prove convexity, take two elements $\pb_1$ and $\pb_2$ in $\pbcred_\prior$ and a constant $\tau \in [0,1]$, and then define the mixture $\pb^\text{mix} = (1-\tau) \pb_1 + \tau \pb_2$; the goal is to show that $\pb^\text{mix} \in \pbcred_\prior$.  If $\Omega \sim \pb^\text{mix}$, then for any event $H$ in $\TT = f(\OO)$, 
%\begin{align*}
%\pb^\text{mix}\{ f(\Omega) \in H \} & = (1-\tau) \, \pb_1\{ f(\Omega) \in H\} + \tau \, \pb_2\{ f(\Omega) \in H \} \\
%& = (1-\tau) \, \prior(H) + \tau \, \prior(H) \\
%& = \prior(H).
%\end{align*}
%This implies that $\pb^\text{mix} \in \pbcred_\prior$ if $\pb_1$ and $\pb_2$ are and, therefore, that $\pbcred_\prior$ is convex.  Finally, to prove that $\pbcred_\prior$ is closed with respect to the weak topology, consider a sequence $(\pb_t: t \geq 1)$ in $\pbcred_\prior$ with a weak limit $\pb_\infty$; the goal is to show that $\pb_\infty \in \pbcred_\prior$.  By the continuous mapping theorem, if $\Omega_t \sim \pb_t$, then $f(\Omega_t) \to f(\Omega_\infty)$ in distribution as $t \to \infty$.  But the distribution of $f(\Omega_t)$ is $\prior$ for all $t$ and, consequently, the distribution of $f(\Omega_\infty)$ is also $\prior$.  This implies $\pb_\infty$ is contained in $\pbcred_\prior$ and, hence, the latter collection is closed. 
%\end{proof} 

An advantage to this credal set characterization is that it allows for the construction of a coherent upper joint distribution for $(Z^N, \Omega)$ based solely on the model $\model$ and Your prior information in $\credal$ about $\Theta$.  
%Furthermore, this construction is conceptually straightforward: just maximize over all the ordinary joint distributions for $(Z^N, \Omega)$ in the above-defined credal set.  
Specifically, let $\uuprob$ denote this induced upper joint distribution for $(Z^N, \Omega)$, which is defined as via its upper expectation as 
\begin{equation}
\label{eq:uuprob}
\uuprob \, g = \sup_{\prior \in \credal} \sup_{\pb \in \pbcred_\prior} \underbrace{\E^{\Omega \sim \pb} \bigl[ \E^{Z \sim \prob_\Omega}\{ g(Z^N, \Omega) \} \bigr]}_{\text{joint distribution expectation}}, 
\end{equation}
where $g: \ZZ^\infty \times \OO \to \RR$ is a suitable function, and the highlighted term is just the usual expectation of $g(Z^N, \Omega)$ with respect to the joint distribution determined by the (marginal) distribution $\pb$ for $\Omega$ and the (now-interpreted-as-a-conditional) distribution $\prob_\omega$ for $Z^N$, given $\Omega=\omega$.  As mentioned in Section~\ref{SS:ip}, $\uuprob$-upper probabilities correspond to the right-hand side above with the appropriate indicator function plugged in for $g$.  
%It's with respect to this upper joint distribution that the regularized e-process has desirable properties, as I demonstrate next.  




\subsection{Regularized Ville's inequality}
\label{SS:reg.ville}

My claim is that the regularized e-process satisfies prior knowledge-dependent versions of property \eqref{eq:eval.bound} and of Ville's inequality \eqref{eq:ville}.  This leads to a corresponding ``regularized'' version of anytime validity discussed further in Section~\ref{SS:implications} below. 

\begin{thm}
\label{thm:reg.ville}
Suppose the available prior information is encoded in the upper probability $\uprior$, which determines an upper joint distribution $\uuprob$ as in \eqref{eq:uuprob}.  If $\rho$ is a regularizer in the sense of Definition~\ref{def:regularizer} relative to $\uprior$, then for any e-process $\{\eval_\theta(\cdot): \theta \in \TT\}$, the corresponding regularized version $\eval^\text{\rm reg}$ in \eqref{eq:reg.eprocess} satisfies 
\begin{equation}
\label{eq:reg.ville0}
\uuprob\bigl[ \eval^\text{\rm reg}\{Z^N,f(\Omega)\} \bigr] \leq 1, \quad \text{for all stopping times $N$}, 
\end{equation}
and the following regularized Ville's inequality holds:
\begin{equation}
\label{eq:reg.ville}
\uuprob \bigl[ \eval^\text{\rm reg}\{Z^N, f(\Omega)\} > \alpha^{-1} \bigr] \leq \alpha \quad \text{all $\alpha \in [0,1]$, all stopping times $N$}. 
\end{equation}
\end{thm}

%\begin{proof}
%The first claim is just a direct computation using \eqref{eq:uuprob} when the function $g = \eval^\text{reg}$ factors as $g(\cdot, \omega) = \rho(f(\omega)) \, \eval_{f(\omega)}(\cdot)$:
%\begin{align*}
%\uuprob \, \eval^\text{reg} & = \sup_{\prior \in \credal} \sup_{\pb \in \pbcred_\prior} \E^{\Omega \sim \pb} \bigl[ \E^{Z \sim \prob_\Omega}\{ \eval^\text{reg}(Z^N, f(\Omega)) \} \bigr] \\
%& = \sup_{\prior \in \credal} \sup_{\pb \in \pbcred_\prior} \E^{\Omega \sim \pb} \bigl[ \rho( f(\Omega) ) \,  \underbrace{\E^{Z \sim \prob_\Omega}\{ \eval_{f(\Omega)}(Z^N) \}}_{\text{$\leq 1$, by \eqref{eq:eval.bound}}} \bigr] \\
%& \leq \sup_{\prior \in \credal} \sup_{\pb \in \pbcred_\prior} \E^{\Omega \sim \pb} \bigl[ \rho\{ f(\Omega) \} \bigr] \\
%& = \sup_{\prior \in \credal} \E^{\Theta \sim \prior} \{ \rho(\Theta) \} \\
%& \leq 1,
%\end{align*}
%where the last inequality follows by definition of the regularizer $\rho$, and the penultimate equality follows by definition of $\pbcred_\prior$: if $\Omega \sim \pb \in \pbcred_\prior$, then the distribution of $f(\Omega)$ is the same as that of $\Theta$ under $\prior$.  The second claim, the regularized Ville's inequality in \eqref{eq:reg.ville}, follows from the first claim and an application of Markov's inequality inside the upper-probability calculation.  If I write ``$(Z,\Omega) \sim \prob_\bullet \otimes \pb$'' to represent the joint distribution of $(Z,\Omega)$ under the model where $\Omega \sim \pb$ and $(Z \mid \Omega=\omega) \sim \prob_\omega$, then:
%\begin{align*}
%\uuprob \bigl[ \eval^\text{\rm reg}\{Z^N, f(\Omega)\} > \alpha^{-1} \bigr] & = \sup_{\prior \in \credal} \sup_{\pb \in \pbcred_\prior} \prob^{(Z,\Omega) \sim \prob_\bullet \otimes \pb}\bigl[ \eval^\text{reg}\{Z^N, f(\Omega) \} > \alpha^{-1} \bigr] \\
%& \leq \sup_{\prior \in \credal} \sup_{\pb \in \pbcred_\prior} \alpha \, \E^{\Omega \sim \pb} \bigl[ \E^{Z \sim \prob_\Omega}\{ \eval^\text{reg}(Z^N, f(\Omega)) \} \bigr] \\
%& = \alpha \, \uuprob \, \eval^\text{reg} \\
%& \leq \alpha,
%\end{align*}
%where the first ``$\leq$'' above is by the usual Markov's inequality and the last line is by the claim \eqref{eq:reg.ville0} proved above.  
%\end{proof}

%{\color{magenta}
%I'll end this subsection with two technical remarks, both concerning Markov's inequality.  The statistical consequences of Theorem~\ref{thm:reg.ville} are discussed in Section~\ref{SS:implications} below. 
%\begin{itemize}
%\item The above proof demonstrates that the $\uuprob$-upper probability defined in \eqref{eq:uuprob} satisfies a version of Markov's inequality.  This isn't a coincidence---many of the familiar probability inequalities have imprecise-probabilistic versions based on properties of the defining Choquet integrals \citep[see, e.g.,][]{wang2011.choquet}.  
%\item 
The above result is related to the``uniformly-randomized Markov inequality'' in \citet[][Theorem~1.2]{ramdas.randomized.markov}.  They prove that, if $X$ and $U$ are independent random variables, with $X \geq 0$ and $U$ stochastically no smaller than $\unif(0,1)$, then $\prob(X / U \geq a^{-1}) \leq a \,\E(X)$ for $a > 0$. 
%The key observation is that the random variable $X$ is being boosted by the division by $U \leq 1$ but the same Markov inequality bound holds, thereby sharpening the usual Markov inequality claims.
%; there are other enhancements discussed in \citet{ramdas.randomized.markov} that I won't mention here.  
The connection between their randomized Markov inequality and the result in Theorem~\ref{thm:reg.ville} is most clear in the special case where $\rho$ is defined in terms of the possibility contour $q$ as in Section~\ref{SS:possibilistic.case}.  In such a case, the random variable $q(\Theta)$ is no smaller than $\unif(0,1)$ when $\Theta \sim \prior$, for any prior $\prior$ in the credal set $\credal$.  Since the calibration step $\gamma \circ q(\Theta)$ boosts it further, the proposed regularization is simply dividing the usual e-process by a no-smaller-than-uniform random variable.  
%Theorem~\ref{thm:reg.ville} is not a special case of Ramdas and Manole's results, however, since they assume independence that doesn't hold in the present case. 
%\end{itemize} 
%}


\subsection{Statistical implications}
\label{SS:implications}

Theorem~\ref{thm:reg.ville}'s most important take-away is its implication concerning safe, anytime valid inference.  Indeed, when relevant, non-vacuous prior information is available,  Theorem~\ref{thm:reg.ville} shows how that knowledge can be used to enhance an e-process in such a way that anytime validity is preserved but efficiency is generally gained.  This enhancement is achieved through the incorporation of a regularizer as in \eqref{eq:reg.eprocess} that inflates and deflates the original e-process when the latter is large and small, respectively.  If the goal is to test ``$\Theta \in H$,'' then the regularized e-process-based testing procedure  
\[ \text{reject ``$\Theta \in H$'' based on data $z^n$} \iff \inf_{\theta \in H} \eval^\text{reg}(z^n,\theta) > \alpha^{-1} \]
will be anytime valid (relative to Your prior knowledge) in the sense that 
\begin{align*}
\uuprob(\text{Type~I error}) := & \; \uuprob\{ \text{$f(\Omega) \in H$ and test of `$f(\Omega) \in H$' rejects} \} \\
= & \; \uuprob\Bigl\{ f(\Omega) \in H \text{ and } \inf_{\theta \in H} \eval^\text{reg}(Z^N,\theta) > \alpha^{-1} \Bigr\} \\
\leq & \; \uuprob\bigl\{ \eval^\text{reg}(Z^N,f(\Omega)) > \alpha^{-1} \bigr\} \\
\leq & \; \alpha,
\end{align*}
where the last line follows by \eqref{eq:reg.ville}.  This is a non-trivial generalization of the familiar frequentist Type~I error control so it warrants some remarks.  In the statistics literature, the prevailing viewpoint is that the quantity of interest $\Theta$ is a fixed unknown, so the statement ``$\Theta \in H$'' is absolutely true for some $H$ and absolutely false for others, but nothing more can be said.  This aligns with the vacuous-prior case where, with the exception of the trivial hypotheses $\varnothing$ and $\TT$, Your corpus of knowledge makes exactly the same statements about every hypothesis: $\lprior(H) = 0$ and $\uprior(H)=1$.  That is, there's no evidence supporting the truthfulness of either ``$\Theta \in H$'' or ``$\Theta \not\in H$.'' With this extreme form of prior knowledge, Your corresponding regularizer is $\rho \equiv 1$---so that $\eval^\text{reg}(\cdot, \theta) = \eval_\theta(\cdot)$---and $\uuprob$ simplifies (see Remark~\ref{re:vac.supremum} in Appendix~\ref{A:remarks}) to yield the following:
\begin{align}
\uuprob(\text{Type~I error}) := & \; \uuprob\{ \text{$f(\Omega) \in H$ and test of `$f(\Omega) \in H$' rejects}\} \notag \\
= & \; \sup_{\omega: f(\omega) \in H} \prob_\omega \Bigl\{ \inf_{\theta \in H} \eval_\theta(Z^N) > \alpha^{-1} \Bigr\} \label{eq:vac.supremum}. 
\end{align}
Then the anytime Type~I error control property of the original e-process-based test, derived from \eqref{eq:ville}, is recovered as a special case of Theorem~\ref{thm:reg.ville}.  
%Moreover, the familiar and intuitive ``supremum-over-the-null'' is simply what appears when the imprecise-probabilistic framework formulated here is applied to the special vacuous-prior case.  
Beyond this extreme case, the different $H$'s have varying degrees of belief/plausibility and this is taken into account in the evaluation of the test via the conjunction: ``$f(\Omega) \in H$'' {\em and} ``test of `$f(\Omega) \in H$' rejects.'' In particular, monotonicity of $\uuprob$ implies that $\uuprob(\text{Type~I error}) \leq \uprior(H)$, which means that the test doesn't have to be particularly good for those $H$ with small $\uprior(H)$ to control the (generalized) Type~I error rate.  If (generalized) Type~I error control is easy for those {\em a priori} implausible $H$'s, then that gives the e-process an opportunity to focus its effort on handling the plausible $H$'s more efficiently.  

Similarly, a nominal $100(1-\alpha)$\% confidence region for $\Theta$, based on the regularized e-process applied to data $z^n$, is 
\[ C_\alpha^\text{reg}(z^n) = \bigl\{ \theta \in \TT: \eval^\text{reg}(z^n, \theta) \leq \alpha^{-1} \bigr\}, \]
i.e., the collection of simple null hypotheses that the regularized e-process-based test would not reject at level $\alpha$.  Then the non-coverage (upper) probability is 
\begin{equation}
\label{eq:reg.noncvg}
\uuprob\{ C_\alpha^\text{reg}(Z^N) \not\ni f(\Omega) \} = \uuprob\{ \eval^\text{reg}(Z^N, f(\Omega)) > \alpha^{-1} \} \leq \alpha, 
\end{equation}
which implies that the corresponding lower probability of coverage, i.e., of $C_\alpha^\text{reg}(Z^N) \ni f(\Omega)$, is bounded from below by $1-\alpha$.  This justifies calling $C_\alpha^\text{reg}$ a (generalized) ``anytime $100(1-\alpha)$\% confidence region.'' Like in the discussion above, this is generally different from the usual notion of coverage probability of a confidence set.  In the case of vacuous prior information, however, the regularizer is $\rho \equiv 1$, the corresponding confidence region is $C_\alpha(z^n) = \{ \theta \in \TT: \eval_\theta(z^n) \leq \alpha^{-1} \}$, and inequality  \eqref{eq:reg.noncvg} reduces to 
\[ \sup_\omega \prob_\omega\{ C_\alpha^\text{reg}(Z^N) \not\ni f(\omega) \} \leq \alpha. \]
This, of course, is exactly the familiar anytime coverage probability guarantees offered by e-process-based confidence sets, via \eqref{eq:ville}.  

%Finally, I want to (re-)emphasize that there's more to this proposal than anytime validity preservation: regularization creates opportunities for efficiency gain.  Recall that Your knowledge encoded in $\credal$ or $\uprior$ is ``real''---no reason to include it otherwise!---so You expect $\rho(\Theta)$ to be small; see Definition~\ref{def:regularizer}.  Therefore, in the absence of data, if $\rho(\theta)$ is large, then You'd infer that ``$\Theta=\theta$'' isn't true.  So, You'd be satisfied if $\theta$'s present in the unregularized e-process confidence region $C_\alpha(z^n)$ happened to be absent from Your $C_\alpha^\text{reg}(z^n)$ if it's because a large regularizer $\rho(\theta)$ value knocked them out.  It's the removal of $\theta$'s that are sufficiently incompatible with the prior knowledge that leads to the improved efficiency.  While it's not generally true that $C_\alpha^\text{reg}(z^n)$ is contained in $C_\alpha(z^n)$ for all $\alpha$ and (almost) all $z^n$, the regularized confidence regions do tend to be narrower than their unregularized counterpart, as demonstrated in Section~\ref{SS:gains} below.  

To summarize, my proposal involves two key steps: first, a regularizer is created and fused with a given e-process in such a way that the resulting regularized e-process tends to be larger at points incompatible with prior knowledge and hence more efficient than the given e-process alone; second, the definition of ``anytime validity'' is correspondingly relaxed, via the prior knowledge-dependent $\uuprob$, to accommodate the broadly larger and more efficient regularized e-process.  Of course, whether a property like ``the method's answer is wrong with small $\uuprob$-probability'' carries weight with readers, peer reviewers, collaborators, etc.~hinges on whether Your prior knowledge incorporated in $\uuprob$ is justified.  The same can be said for any kind of model-based analysis: the conclusions drawn can be no more convincing than the justification given for the assumed model.  But compared to a Bayesian framework that requires prior knowledge be encoded as a precise probability distribution, here I'm placing no serious constraints on the form Your prior knowledge takes; You're free to use whatever $\uprior$ You can justify, even if it's vacuous.  Your target audience is scientifically obligated to scrutinize Your justification, but if Your beliefs warrant being called {\em prior knowledge}, You must have strong justification for them and so passing this scrutiny should only be a matter of explaining Your reasoning clearly.  




\subsection{Efficiency gains}
\label{SS:gains}

The goal here is to demonstrate the efficiency that can be gained through the incorporation of partial prior information via my proposed regularized e-process framework.  I'll take a simple model and e-process, so that effort can be focused on the possibilistic prior, the corresponding regularizer, and its effect on efficiency.  I must emphasize that {\em I'm not recommending off-the-shelf use of any particular regularizer}.  It's Your responsibility to determine what, if any, relevant prior information about $\Theta$ is available and how to quantify it.  I'm in no position to say what You should or shouldn't believe about $\Theta$.  

To set the scene, consider a normal mean model with known variance equal to 1.  I'm assuming a parametric model, hence the uncertain model index, $\Omega$, and the quantity of interest, $\Theta$, are the same.  So I'll write $\model = \{ \prob_\theta: \theta \in \TT\}$, where $\prob_\theta = \nm(\theta, 1)$ and $\TT = \RR$.  For the e-process, I'll consider the Bayes factor
\[ \eval_\theta(\cdot) = \frac{\int L_t(\cdot) \, \xi(t) \, dt}{L_\theta(\cdot)}, \]
where $L_\theta$ is the usual Gaussian likelihood function and, in the numerator, $\xi$ is a mixing probability density function on $\TT$.  In what follows, I'll take $\xi$ to be the $\nm(0, v)$ density function, with $v=10$.  This is a version of the so-called Savage--Dickey e-process \citep[e.g.,][]{grunwald.epost}.  The integration can be done in closed-form, so
\[ \eval_\theta(z^n) = (nv + 1)^{1/2} \exp\bigl\{ -\tfrac{n}{2} (\theta - \bar z_n)^2 + \tfrac12 ( \tfrac{n}{nv + 1} ) \, \bar z_n^2 \bigr\}, \quad \theta \in \TT. \]
In words, the e-process is a ratio of the (Bayesian marginal) likelihood under the model where $\Theta$ is different from $\theta$---where ``different from $\theta$'' is quantified by $\xi$---to the likelihood under the model where $\Theta$ equals $\theta$.  It's important to point out that $\xi$ isn't genuine prior information, it's just a (default) choice that's made to define the e-process.  
%Virtually all e-processes have a likelihood ratio form as above, often involving mixtures.  
The mixing distributions may be relatively diffuse, hence my choice of $\xi$'s variance as $v=10$.  

Here I'll describe the construction of a partial prior that takes the form of a possibility measure, based on an assessment of {\em surprise}, as mentioned briefly above.  This is not the only kind of partial prior one can consider, and some additional details and examples are shown in Appendix~\ref{S:emp.efficiency}.  The strategy here, which is quite common, is to assign Your surprise by analogy, by making comparison to familiar probability assignments.  For example, suppose that, for each $t > 0$, You'd be as surprised to learn that $|\Theta| > t$ as you would be to observe $|Y| > t$, where $Y \sim \nm(0,K)$ for a specified $K$.  Matching Your surprise to the latter probability amounts to choosing a prior contour 
\[ q(\theta) %= \prob\{ {\tt dnorm}(Y, 0, K^{1/2}) \leq {\tt dnorm}(\theta, 0, K^{1/2})\} 
= 1 - {\tt pchisq}(\theta^2 / K, {\tt df}=1), \quad \theta \in \TT. \]
This is a principled possibility measure construction, and further details can be found in Remark~\ref{re:prob.to.poss} of Appendix~\ref{A:remarks}. 
%called the {\em probability-to-possibility transform} \citep[e.g.,][]{dubois.etal.2004}, a very powerful tool.  
Importantly, what I just described is entirely different from You taking $\nm(0,K)$ as a prior distribution for $\Theta$, since $\nm(0,K)$ is just one of the many probability measures contained in the possibilistic prior's credal set. 

%The examples that follow consider three different forms of partial prior information and the corresponding regularization of the above simple e-process.  So, all that I'll do here is offer examples of prior knowledge that might be available in an application, to show how that information can be encoded as a possibility contour and converted into a regularizer and to demonstrate the effects that incorporation of this information can have on the results, in particular, as it pertains to efficiency.  

Figure~\ref{fig:gauss} plots the regularized (and unregularized) log-transformed e-process as a function of $\theta$ for the ``Gaussian prior'' described above, for $K \in \{0.1, 0.2, 0.4, 0.8\}$, and for three different values of the observed sample mean $\bar z$ based on a sample of size $n=5$:
\begin{itemize}
\item $\bar z=0.25$ is consistent with all the priors;
\vspace{-2mm}
\item $\bar z=0.5$ is marginally inconsistent with the priors, and;
\vspace{-2mm}
\item $\bar z=1$ is rather inconsistent with the priors.
\end{itemize}
The black line corresponds to the unregularized e-process, and the four colored lines correspond to the different values of $K$; the dashed horizontal line corresponds to $-\log0.05 \approx 3$, which is the cutoff that determines an e-process's 95\% confidence interval.  As expected, the stronger the prior information (i.e., smaller $K$), the greater the efficiency gains.   

\begin{figure}[t]
\begin{center}
\subfigure[$\bar z=0.25$]{\scalebox{0.6}{\includegraphics{nm_gauss_x025}}}
\subfigure[$\bar z=0.5$]{\scalebox{0.6}{\includegraphics{nm_gauss_x05}}}
\subfigure[$\bar z=1$]{\scalebox{0.6}{\includegraphics{nm_gauss_x1}}}
\end{center}
\caption{Plot of $\theta \mapsto \eval^\text{reg}(z^n, \theta)$ for three different data sets $z^n$.}
\label{fig:gauss}
\end{figure}

Another perspective on the efficiency gain concerns the number of sample points required before a false hypothesis can be rejected.  Here I'll consider the hypothesis ``$\Theta=0.7$'' and track the (log) regularized e-process's growth as the sample size increases.  The fewer number of sample points required to reject the false hypothesis, the more efficient the procedure is.  Figure~\ref{fig:gauss.grow} plots the path $n \mapsto \text{avg}\{\log \eval^\text{reg}(z^n, 0.7)\}$, where the average is taken over 1000 replications of the data-generating process.  Clearly, the strongest regularization ($K=0.1$) makes it possible to detect the discrepancy of 0.7 with generally fewer samples, whereas the other priors, which offer weaker regularization, take more samples on average to reach the desired conclusion.  
%Except for the weakest prior information ($K=0.8$), regularization results in greater efficiency.
%; see Section~\ref{S:discuss}.  

\begin{figure}[t]
\begin{center}
\scalebox{0.6}{\includegraphics{nm_gauss_grow}}
\end{center}
\caption{Plot of $n \mapsto \text{avg}\{\log \eval^\text{reg}(z^n, 0.7)\}$ when data are sampled from $\nm(0,1)$.}
\label{fig:gauss.grow}
\end{figure}





\section{Regularized e-uncertainty quantification}
\label{S:reg.uq}

\subsection{Objective}

So far, I've treated inference as the construction of suitable test and confidence set procedures.  In this section, however, I have a more ambitious goal of reliable, data-driven, uncertainty quantification about the uncertain $\Theta$.  
%Following Fisher and others, my goal here is to uncover what Efron called the ``Holy Grail of statistics.'' 

It's a mathematical fact that uncertainty arising in statistical inference generally cannot be quantified reliably using ordinary probability.  Indeed, the {\em false confidence theorem} \citep{balch.martin.ferson.2017} says, roughly, every probabilistic quantification of uncertainty---Bayes, fiducial, etc.---tends to assign high confidence, i.e., high posterior probability, to certain false hypotheses \citep{martin.nonadditive, imchar, martin.belief2024}.  False confidence creates a risk of unreliability, or ``systematically misleading conclusions'' \citep{reid.cox.2014}.  
%Fortunately, alternatives to probability theory already exist, such as Dempster--Shafer theory \citep{dempster2008, shafer1976}, possibility theory \citep{dubois.prade.book}, and the theory of lower previsions \citep{walley1991, lower.previsions.book}.  
{\em Inferential models} (IMs), originally developed in \citet{imbasics, imbook}, offer a new framework for possibilistic statistical reasoning that's reliable in the sense of avoiding false confidence; see the recent review in \citet{imreview}.   Remark~\ref{re:why.possibility} in Appendix~\ref{A:remarks} explains why the possibilistic form is right for statistical problems. 

%First I give a brief overview of the recent formulation of {\em possibilistic IMs}, and then I'll discuss a new twist on this IM formulation that involves e-processes \citep{martin.basu}.  Partial prior knowledge can also be incorporated, leading to a framework of regularized e-process-based uncertainty quantification.  The corresponding e-possibilistic IM and its reliability properties are developed in Section~\ref{SS:euq.properties}, and certain unexpected-but-welcome behavioral consequences of this reliability in the context are established in Section~\ref{SS:euq.behavior}.  

%In particular, Theorem~\ref{thm:no.sure.loss} below establishes that the proposed e-possibilistic IM, interpreted as a prior-to-posterior beliefs updating rule, avoids sure-loss (and more); and Theorem~\ref{thm:imdec} generalizes and strengthens the decision-theoretic reliability results in \citet{imdec} and \citet{grunwald.epost}.  




\subsection{Possibilistic IMs}
\label{SS:ims}

My current efforts \citep{martin.partial2, martin.partial3, martin.basu} focus on the construction of {\em possibilistic IMs}, where the IM output takes the mathematical form of a possibility measure as reviewed briefly in Section~\ref{SS:ip}.  The specific proposal put forward in the above references starts with defining the IM's possibility contour, based on observed data $Z^n=z^n$, as
\[ \pi_{z^n}(\theta) = \sup_{\omega: f(\omega) = \theta} \prob_\omega\{ r(Z^n,\theta) \leq r(z^n,\theta)\}, \quad \theta \in \TT, \]
where $r(z^n,\theta)$ provides a ranking of the parameter value $\theta$ in terms of its compatibility with $z$---large values indicate higher compatibility.  For example, in \citet{martin.partial2, martin.partial3}, I recommended taking $r$ to be the relative profile likelihood 
\[ R(z^n,\theta) = \frac{\sup_{\omega: f(\omega)=\theta} L_{z^n}(\omega)}{\sup_\omega L_{z^n}(\omega)}, \quad \theta \in \TT, \]
with $\omega \mapsto L_{z^n}(\omega)$ the likelihood function corresponding to the ``iid $\prob_\omega$'' model. Then the possibilistic IM's upper probability output is, as in \eqref{eq:upper.p.poss}, given by optimization: 
\begin{equation}
\label{eq:upi}
\uPi_{z^n}(H) = \sup_{\theta \in H} \pi_{z^n}(\theta), \quad H \subseteq \TT. 
\end{equation} 
The interpretation is that a small $\uPi_{z^n}(H)$ means there's strong evidence in data $z^n$ against the truthfulness of $H$.  That is, a small upper probability assigned to $H$ implies doubt and, therefore, I'd be inclined to ``reject'' $H$.  But how small is ``small''?  Clearly some calibration of the IM's numerical output is needed in order for this line of reasoning to be reliable, i.e., to ensure that I don't systematically doubt true hypotheses or buttress false hypotheses.  The possibilistic IM's calibration property is
\begin{equation}
\label{eq:valid}
\sup_{\omega: f(\omega) \in H} \prob_\omega\{ \uPi_{Z^n}(H) \leq \alpha \} \leq \alpha, \quad \alpha \in [0,1], \quad H \subseteq \TT. 
\end{equation}
In words, \eqref{eq:valid} says that the IM assigning small upper probability values to true hypotheses about $\Theta$ is itself a small probability event.  Since there's an explicit link between the two notions of ``small,'' this offers the calibration needed to avoid systematically erroneous inferences.  
%Compare this to what \citet{walley2002} calls the ``fundamental frequentist principle.''  
The analogous property concerning the IM's lower probabilities---recall the conjugacy relationship, $\lPi_{z^n}(H) = 1 - \uPi_{z^n}(H^c)$, from Section~\ref{SS:ip}---reads as follows: 
\[ \sup_{\omega: f(\omega) \not\in H} \prob_\omega\{ \lPi_{Z^n}(H) \geq 1-\alpha \} \leq \alpha, \quad \alpha \in [0,1], \quad H \subseteq \TT. \]
This latter result says that the IM assigning large lower probability---or ``confidence''---to false hypotheses is a small-probability event; now it should be clear why I say that the IM is safe from the risk of false confidence.  

It's based on the property \eqref{eq:valid} that I say the IM is {\em valid} and, in turn, that uncertainty quantification based on the IM output is reliable.  Among other things, it's straightforward to establish that test and confidence procedures derived from the IM output achieve the usual frequentist error rate control guarantees.  
%But there's more to the IM output than the statistical procedures derived from it \citep{martin.partial2, cella.martin.probing}. 

%Statistical methods---from classical to the possibilistic IM described above---often assume the data-generating process is fully known.  This can be an issue when the data are collected sequentially and the (possibly data-dependent) rule by which the collection stops isn't explicitly stated.  Consider (say) an iid sequence $Z_1,Z_2,\ldots$ from $\prob_\omega$, with $Z^n = (Z_1,\ldots,Z_n)$ the first $n$ instances along the sequence, and define a filtration $\F_n = \sigma(Z^n)$, $n \geq 1$.  A stopping time $N$ is an integer-valued random variable such that $\{N \leq n\} \in \F_n$, $n \geq 1$.  Then a realization of the data $Z^N$ is of the form $z^n$, with $n$ the observed value of $N$ and $z^n=(z_1,\ldots,z_n)$ are the observed $Z$ values.  If the data analyst isn't privy to which stopping time $N$ was employed, then he/she might just assume that $N \equiv n$ was fixed at the observed value in advance.  But this could be very misleading, since the relevant sampling distributions of statistics based on $Z^n$ could be drastically different that based on $Z^N$, thus jeopardizing the statistical method's reliability.  I'll say that a statistical method is {\em safely reliable} if its reliability holds uniformly over stopping times.  How can the above IM be made safely reliable? 



\subsection{Connection to e-processes}
\label{SS:im.e}

The discussion above focused on the case of a fixed sample size.  It's well-known, however, that the fixed-$n$ sampling distribution properties---such as validity in \eqref{eq:valid}---generally fail when the stopping rule $N$ is data-dependent.  
%For example, if the value of $N$ is determined by continuing to collect data until the IM's upper probability assigned to $H$ is more than 0.1, then $\prob_\omega\{ \uPi_{Z^N}(H) > 0.1 \} \equiv 1$ for all $\omega$, which violates \eqref{eq:valid}.  
If the stopping rule $N$ that's used is {\em known}, then this can be easily incorporated into the IM construction such that validity is achieved.  
%Indeed, for any given ranking $r$, the contour 
%\[ \pi_{z^n}(\theta) = \sup_{\omega: f(\omega)=\theta} \prob_\omega\{ r(Z^N, \theta) \leq r(z^n, \theta) \}, \quad \theta \in \TT, \]
%determines a possibilistic IM that's valid in the sense of 
%\[ \sup_{\omega: f(\omega) \in H} \prob_\omega\{ \uPi_{Z^N}(H) \leq \alpha \} \leq \alpha, \quad \alpha \in [0,1], \quad H \subseteq \TT. \]
%Note that the contour in penultimate display takes the sampling distribution of $Z^N$ into account, so the calibration properties of $\uPi_{Z^N}$ follow from exactly the same arguments as when the sample size $n$ was fixed in advance.  
The trouble, of course, is that the $N$ employed in the data-collection process is often {\em unknown}.
%, e.g., the original investigators might not explain all the study details in their paper.  
\citet{martin.basu} suggested a conceptually simple work-around that interprets $N$ as a ``nuisance parameter.'' Then the general rules in \citet{martin.partial3} for handling nuisance parameters immediately suggest a new IM contour 
\[ \pi_{z^n}(\theta) = \sup_N \sup_{\omega: f(\omega)=\theta} \prob_\omega\{ r(Z^N,\theta) \leq r(z^n,\theta)\}, \quad \theta \in \TT, \]
where the outermost supremum is over all the stopping rules in consideration (that are consistent with the observation ``$N=n$'').  Of course, direct computation of the right-hand side above can be challenging, depending on the complexity of the set of stopping rules in consideration.  It's here that e-processes come in handy.  

Recall that the ranking function $r$ in the possibilistic IM construction is quite flexible.  One alternative to the proposal above is to take $r$ as the reciprocal of an e-process: $r(z^n, \theta) = \eval_\theta(z^n)^{-1}$. In that case, the usual Ville's inequality gives 
\begin{align*}
\pi_{z^n}(\theta) & = \sup_N \sup_{\omega: f(\omega)=\theta} \prob_\omega\{ r(Z^N,\theta) \leq r(z^n,\theta)\} \\
& = \sup_N \sup_{\omega: f(\omega)=\theta} \prob_\omega\{ \eval_\theta(Z^N) \geq \eval_\theta(z^n) \} \\
& \leq 1 \wedge \eval_\theta(z^n)^{-1}.
\end{align*}
Set $\pi_{z^n}^\eval(\theta) = 1 \wedge \eval_\theta(z^n)^{-1}$ to be the upper bound, so that 
\[ \pi_{z^n}(\theta) \leq \pi_{z^n}^\eval(\theta) := 1 \wedge \eval_\theta(z^n)^{-1}, \quad \theta \in \TT. \]
This upper bound corresponds to what \citet{grunwald.epost} calls his capped {\em e-posterior}.  But there are two further observations about this bound that are worth noting.  First, the bound itself is typically (see Remark~\ref{re:typically} in Appendix~\ref{A:remarks}) a possibility contour.  So, there's a corresponding possibilistic IM, with coherent upper probability $\uPi_{z^n}^\eval$ defined via optimization as in \eqref{eq:upi}.  This determines an {\em e-possibilistic IM} for uncertainty quantification about $\Theta$.  Second, that $\pi_{z^n}^\eval$ is an upper bound of the contour $\pi_{z^n}$, and that the latter determines an anytime valid IM, immediately implies that the IM corresponding to the former is anytime valid too.  This discussion is summarized as

\begin{thm}
\label{thm:valid.vac}
Given an e-process $\eval$, the corresponding e-possibilistic IM with upper probability determined by optimization of the contour $\pi_{z^n}^\eval$, i.e.,  
\[ \uPi_{z^n}^\eval(H) = \sup_{\theta \in H} \pi_{z^n}^\eval(\theta), \quad H \subseteq \TT, \]
is anytime valid in the sense that 
\begin{equation}
\label{eq:valid.vac}
\sup_{\omega: f(\omega) \in H} \prob_\omega\{ \uPi_{Z^N}^\eval(H) \leq \alpha \} \leq \alpha, \quad \text{all $\alpha \in [0,1]$, all $N$, all $H \subseteq \TT$}. 
\end{equation}
\end{thm}

%\begin{proof}
%Since the IM with contour $\pi^\eval$ dominates that with contour $\pi$ defined earlier, the anytime validity of the latter implies that of the former.  For concreteness, however, I'll give a direct proof of anytime validity of $\uPi^\eval$.  
%
%For any data set $z^n$, the IM's possibility measure output $H \mapsto \uPi_{z^n}(H)$ is monotone.  Therefore, if $\omega$ is such that $f(\omega) \in H$, then 
%\[ \uPi_{z^n}^\eval(H) \geq \uPi_{z^n}^\eval(\{ f(\omega) \}) = \pi_{z^n}^\eval(f(\omega)), \]
%and, consequently, for any $\alpha \in [0,1]$, 
%\[ \uPi_{z^n}(H) \leq \alpha \implies \pi_{z^n}^\eval(f(\omega)) \leq \alpha \iff \eval_{f(\omega)}(z^n) \geq \alpha^{-1}, \]
%where the right-most property is by definition of the IM's contour function in terms of the e-process's reciprocal.  It follows that 
%\[ \sup_{\omega: f(\omega) \in H} \prob_\omega\{ \uPi_{Z^N}^\eval(H) \leq \alpha \} \leq \sup_{\omega: f(\omega) \in H} \prob_\omega\{ \eval_{f(\omega)}(Z^N) \geq \alpha^{-1} \} \leq \alpha, \]
%with the last inequality due to \eqref{eq:ville}, thus completing the proof.
%\end{proof}

As in Section~\ref{SS:ims}, this calibration property is important because, without it, there'd be no meaningful justification for any particular interpretation of the numerical values an e-possibilistic IM assigns to various inputs $(z^n, H)$.  But in light of Theorem~\ref{thm:valid.vac}, the same reliability-guaranteeing ``no-false-confidence'' statement made above for the simple possibilistic IM also holds here for the new e-possibilistic IM.  
%\citet{grunwald.epost} worked with the same contour function $\pi_{z^n}^\eval$ as above but he didn't treat his ``e-posterior'' as a possibility distribution for reliable, anytime valid uncertainty quantification in the broader sense considered here.   

One important consequence of anytime validity is that one can derive test procedures from the IM and these will inherit the desired frequentist properties.  Indeed, for a given hypothesis $H \subset \TT$ about $\Theta$, and a given level $\alpha \in [0,1]$, the test procedure 
\[ \text{reject ``$\Theta \in H$'' if and only if $\uPi_{Z^N}^\eval(H) \leq \alpha$} \]
controls the frequenist Type~I error rate at the specified level $\alpha$.  This result is ``obvious'' because $\uPi_{Z^N}^\eval(H)$ can be small if and only if $\eval_\theta(Z^N)$ is large for all $\theta \in H$: in particular, 
\[ \uPi_{Z^N}^\eval(H) \leq \alpha \iff \inf_{\theta \in H} \eval_\theta(Z^N) \geq \alpha^{-1}. \]
If the hypothesis is true, i.e., if $\Theta = f(\Omega) \in H$, then the right-most event in the above display implies $\eval_{f(\Omega)}(Z^N) \geq \alpha^{-1}$, which, by Ville's inequality \eqref{eq:ville}, has $\prob_\Omega$-probability no more than $\alpha$, uniformly over stopping times $N$.  

For the case of vacuous prior information, there's a consequence of anytime validity \eqref{eq:valid.vac} that deserves mention; this property is used in the proof of Theorem~\ref{thm:valid.vac}.  For the more general prior information cases to be considered below, this new property is actually stronger than \eqref{eq:valid.vac}, but here the two are equivalent.  

\begin{cor}
\label{cor:strong.valid.vac}
Given an e-process $\eval$, the corresponding e-possibilistic IM is strongly anytime valid in the sense that its contour satisfies 
\begin{equation}
\label{eq:strong.valid.vac}
\sup_{\omega \in \OO} \prob_\omega\bigl[ \pi_{Z^N}^\eval\{ f(\omega) \} \leq \alpha \bigr] \leq \alpha, \quad \text{all $\alpha \in [0,1]$, all $N$}.
\end{equation}
\end{cor}

%\begin{proof}
%Apply Theorem~\ref{thm:valid.vac} with $H$ running over all the singletons in $\TT = f(\OO)$. 
%\end{proof}

Two comments related to Corollary~\ref{cor:strong.valid.vac} are in order.  First, similar to the points about the construction of anytime valid testing procedures, it's straightforward to construct confidence sets too.  Given a significance level $\alpha \in [0,1]$, a $100(1-\alpha)$\% confidence set derived from the e-process-based IM is the $\alpha$-level set defined by the contour:
\[ C_\alpha^\eval(Z^N) = \{ \theta \in \TT: \pi_{Z^N}^\eval(\theta) > \alpha\}. \]
Then strong anytime validity implies that this is, indeed, a genuine anytime valid confidence set in the sense that $\sup_{\omega \in \OO} \prob_\omega\{ C_\alpha^\eval(Z^N) \not\ni f(\omega) \} \leq \alpha$ for all stopping times $N$.
%, \quad \text{all stopping times $N$}. \]
But, by definition of $\pi^\eval$, it's easy to see that $C_\alpha^\eval$ is identical to
%can be re-expressed as 
%\[ C_\alpha^\eval(Z^N) = \{ \theta \in \TT: \eval_\theta(Z^N) < \alpha^{-1}\}, \]
%which the reader will recognize as 
the (unregularized) confidence set determined by the e-process $\eval$.  Then the above non-coverage probability bound follows immediately from the relevant properties of $\eval$ presented above.  

Second, to see why this strong anytime validity \eqref{eq:strong.valid.vac} is ``stronger'' than anytime validity \eqref{eq:valid.vac}, note that monotonicity of possibility measures implies
\begin{equation}
\label{eq:monotone}
\uPi_{Z^N}^\eval(H) \geq \uPi_{Z^N}^\eval(\{f(\omega)\}) \equiv \pi_{Z^N}^\eval(f(\omega)) \quad \text{for all $H$ with $H \ni f(\omega)$}.
\end{equation} 
Interestingly, the above holds even for hypotheses $H$ that depend on data $Z^N$ in some way, e.g., if an adversary with access to the data $Z^N$ is trying to dupe the statistician by selecting a ``most difficult'' hypothesis $H$ post hoc.  This suggests that strong anytime validity is equivalent to a {\em uniform-in-hypotheses} version of anytime validity, and the following corollary states this explicitly; see, also, \citet{cella.martin.probing}.  Note the difference between \eqref{eq:uniform.valid.vac} below and \eqref{eq:valid.vac} above: the former has varying $H$ on the {\em inside} of the probability statement whereas the latter has varying $H$ on the {\em outside}.  

\begin{cor}
\label{cor:uniform.vac}
The e-process-based possibilistic IM constructed above is uniformly anytime valid in the sense that 
\begin{equation}
\label{eq:uniform.valid.vac}
\sup_{\omega \in \OO} \prob_\omega\{ \uPi_{Z^N}^\eval(H) \leq \alpha \text{ for some $H$ with $H \ni f(\omega)$} \} \leq \alpha, \quad \text{all $\alpha \in [0,1]$, all $N$}. 
\end{equation}
\end{cor}

%\begin{proof}
%The proof is based on the following observation:
%\begin{equation}
%\label{eq:equivalence} 
%\uPi_{Z^N}^\eval(H) \leq \alpha \text{ for some $H$ with $H \ni f(\omega)$} \iff \pi_{Z^N}^\eval(f(\omega)) \leq \alpha. 
%\end{equation}
%The ``$\Longleftarrow$'' direction is obvious since $H := \{f(\omega)\}$ is a hypothesis that contains $f(\omega)$.  The ``$\Longrightarrow$'' is similarly obvious by the monotonicity property as stated in \eqref{eq:monotone}. Since Corollary~\ref{cor:strong.valid.vac} says the right-most event in the above display has probability no more than $\alpha$, uniformly in $\omega$ and in $N$, the same must be true of the equivalent left-most event.  
%\end{proof}

Therefore, the e-possibilistic IM constructed here offers reliable, anytime valid uncertainty quantification that's safe not just against fixed hypotheses as in \eqref{eq:valid.vac}, but also against possibly adversarial or otherwise data-dependent hypotheses as in \eqref{eq:uniform.valid.vac}.  This is important for at least two reasons.  First, investigators feeling ``publish-or-perish'' pressures might succumb to temptations to explore for hypotheses that are incompatible with their data, to secure a statistically significant result.  A commitment to using e-possibilistic IMs prevents these sociological temptations from causing harm. Second, as \citet{mayo.book.2018} articulates, investigators yearn for more than null hypothesis significance tests---they want to probe 
%.  Conclusions about a ``null hypothesis'' mark the beginning of the scientific investigation, not the end, so it's imperative that investigators can probe 
for hypotheses that might be compatible with their data, without risking unreliability.  Multiplicity corrections can't accommodate this.
%Multiplicity corrections can't accommodate this kind of probing; see \citet{cella.martin.probing}.





\subsection{Partial priors and regularization}
\label{SS:euq.properties}

%I imagine that it's clear to the reader where this is going.  
Section~\ref{S:reg.eprocess} showed how prior knowledge can be encoded as a regularizer and then combined with a given e-process such that the resulting regularized e-process is anytime valid in a relaxed sense.  The next step is to flip this regularized e-process into a regularized e-possibilistic IM and to establish the corresponding anytime validity properties in the context of uncertainty quantification.  The remainder of this subsection details these next steps.  Section~\ref{SS:euq.behavior} below deals with some more nuanced behavioral properties.
%questions about the behavioral properties of the proposed possibilistic IM.  

%Following the developments in Section~\ref{SS:im.e} above, 
Define the regularized e-possibilistic IM contour function for $\Theta$, given $z^n$, as 
\begin{equation}
\label{eq:contour.reg}
\pi_{z^n}^{\eval \times \rho}(\theta) = 1 \wedge \eval^\text{reg}(z^n, \theta)^{-1}, \quad \theta \in \TT. 
\end{equation}
The corresponding possibility lower and upper probabilities are denoted as $\lPi_{z^n}^{\eval \times \rho}$ and $\uPi_{z^n}^{\eval \times \rho}$, respectively, with the latter defined via optimization and the former via conjugacy, as usual.  By definition of $\pi^{\eval \times \rho}$ above, the level sets 
\[ C_\alpha^{\eval \times \rho}(z^n) = \{ \theta \in \TT: \pi_{z^n}^{\eval \times \rho}(\theta) > \alpha \}, \quad \alpha \in [0,1], \]
are identical to the regularized e-process-based confidence regions $C_\alpha^\text{reg}$ in Section~\ref{SS:implications}.  Therefore, the sets $C_\alpha^{\eval \times \rho}$ inherit the same non-coverage upper probability bound as in \eqref{eq:reg.noncvg}, hence can be referred to as (generalized) ``anytime valid confidence regions.'' 

Theorem~\ref{thm:strong.valid} below extends Corollary~\ref{cor:strong.valid.vac} to the regularized e-possibilistic IM case.  It's an immediate consequence of Theorem~\ref{thm:reg.ville} and the definition of $\pi^{\eval \times \rho}$.  Recall the upper joint distribution $\uuprob$ for $(Z^N,\Omega)$, defined in \eqref{eq:uuprob} above.  

\begin{thm}
\label{thm:strong.valid}
Given an e-process $\eval$ and regularizer $\rho$, the corresponding regularized e-possibilistic IM with contour \eqref{eq:contour.reg} is strongly anytime valid in the sense that it satisfies 
\begin{equation}
\label{eq:strong.valid}
\uuprob\bigl[ \pi_{Z^N}^{\eval \times \rho}\{ f(\Omega) \} \leq \alpha \bigr] \leq \alpha, \quad \text{all $\alpha \in [0,1]$, all $N$}.
\end{equation}
\end{thm}

This is a generalization of Corollary~\ref{cor:strong.valid.vac} because, if the prior information is vacuous and $\rho \equiv 1$, then $\pi_{Z^N}^{\eval \times \rho} \equiv \pi_{Z^N}^\eval$ and the upper joint distribution $\uuprob$ reduces to the supremum of the $\prob_\omega$-probabilities over all $\omega \in \OO$ as in \eqref{eq:strong.valid.vac}. 

%Theorem~\ref{thm:strong.valid} has a couple immediate consequences that are perhaps easier to interpret from a statistical perspective.  These are summarized in the following corollary.

\begin{cor}
\label{cor:strong.valid}
For the same regularized e-possibilistic IM considered in Theorem~\ref{thm:strong.valid}, the following hold for all thresholds $\alpha \in [0,1]$ and all stopping times $N$:
\begin{align*}
\uuprob\bigl\{ \text{there exists $H$ with $H \ni f(\Omega)$ and $\uPi_{Z^N}^{\eval \times \rho}(H) \leq \alpha$} \bigr\} & \leq \alpha \\
\uuprob\bigl\{ \text{$H \ni f(\Omega)$ and $\uPi_{Z^N}^{\eval \times \rho}(H) \leq \alpha$} \bigr\} & \leq \alpha, \quad H \subseteq \TT. 
\end{align*}
\end{cor}

The interpretation is, as before, in terms of reliability.  The second line says that, for any fixed $H$, the upper probability that $H$ is true {\em and} the IM assigns it small upper probability is small.  The first line is stronger, it says that the probability the IM assigns small upper probability to any true hypothesis about $f(\Omega)$ is small.  
%And just as before, there's an analogous result for the regularized e-possibilistic IM's lower probability.  
To make sense of this, suppose data $z^n$ is observed and, for a relevant hypothesis, it happens that $\uPi_{z^n}^{\eval \times \rho}(H)$ is smaller than a threshold that You deem  to be ``sufficiently small.''  Then Corollary~\ref{cor:strong.valid} offers justification for You to draw the inference that $H$ is false---if it were true, then it would be a ``rare event'' that the IM assigned it such a small upper probability.  

Finally, since
%I'll end this subsection on a slightly different note.  The following comments apply to both the regularized and unregularized e-possibilistic IM, but I'll focus attention on the regularized version since the unregularized version is a special case.  Recall 
e-processes are likelihood ratios of some sort, the IM output $(\lPi_{z^n}^{\eval \times \rho}, \uPi_{z^n}^{\eval \times \rho})$ depends on the data $z^n$ only through the likelihood function and no other specifics concerning a ``model'' are used for drawing inferences.
%; e.g., the threshold is $\alpha$ not some $\alpha$-quantile of a model-dependent sampling distribution.  
This implies that the IM satisfies both the {\em likelihood principle} \citep[e.g.,][]{birnbaum1962, basu1975, bergerwolpert1984} and the strong frequentist-like properties derived above.  This is remarkable because the likelihood principle is Bayesians' turf, but I've shown that both Bayesian-like uncertainty quantification is in reach without sacrificing frequentist-like reliability. 

%it's *possible* (pun intended) to quantify uncertainty in a Bayesian-like way that satisfies the likelihood principle without sacrificing frequentist-like reliably. 
%\begin{quote}
%{\em I, myself, came to take... Bayesian statistics... seriously only through recognition of the likelihood principle.} (Savage---discussion of \citealt{birnbaum1962})
%\end{quote}
%Here I've shown that it's *possible* (pun intended) to quantify uncertainty in a Bayesian-like way that satisfies the likelihood principle without sacrificing frequentist-like reliably.  


\subsection{Illustration}
\label{SS:gue}

As a quick illustration, consider inference on the median, $\Theta$, of an underlying distribution; of course, other quantiles can be handled similarly.  Suppose that the observable data are scalars with fully unknown distribution $\prob_\Omega$---that is, the index set $\OO$ is in one-to-one correspondence with the set of distributions on $\RR$.  One way to characterize the mapping from a generic distribution $\prob_\omega$ to the median $\theta$ is through the introduction of a loss function $L_t(z) = |z-t|$, and then defining 
\[ \theta = f(\omega) := \arg\min_t \prob_\omega L_t, \]
the minimizer of the expected absolute difference.  \citet{dey.gue} developed a general e-process construction, extending the universal inference framework of \citet{wasserman.universal} to cases where the quantity of interest is the minimizer of an expected loss.  Their proposed e-process is a sort of ``likelihood ratio,'' where the individual ``likelihood'' terms aren't genuine likelihoods, but suitable quasi-likelihoods constructed from the provided loss function.  Specifically, they define 
\[ \eval_\theta^n = \exp\Bigl[ -\eta \sum_{i=1}^n\{ L_{\hat\theta^{i-1}}(z_i) - L_\theta(z_i)\} \Bigr], \]
where $L_t(z) = |z-t|$ is as above, $\eta > 0$ is a tuning parameter to be discussed further below, and $\hat\theta^k = \hat\theta(z^k) = \text{\tt median}(z_1,\ldots,z_k)$, for $k \geq 1$, with $\hat\theta^0$ any fixed constant.  Then there exists $\eta > 0$ such that $\eval_\theta^n$ is an e-process for the median $\Theta$.  

\citet{dey.gue} describe a data-driven approach for choosing $\eta$ but, for simplicity, I'll take a fixed $\eta$ that depends on a mild assumption about the data-generating process.  By the triangle inequality, the absolute loss function is 1-Lipschitz, i.e., 
\[ \bigl| L_t(z) - L_\theta(z) \bigr| \leq |t - \theta|, \quad \text{all $(z,t,\theta)$}. \]
Then the loss difference $L_t(Z) - L_\theta(Z)$ is a bounded random variable for each pair $(t,\theta)$, hence it's subgaussian and, in particular, it satisfies 
\[ \sup_{\omega: f(\omega) = \theta} \prob_\omega e^{-\eta(L_t - L_\theta)} \leq e^{-\eta \cdot \inf_{\omega: f(\omega)=\theta} \prob_\omega (L_t - L_\theta) + \eta^2 (t - \theta)^2 / 4}. \]
Then the right-hand side is upper-bounded by 1 if and only if
\[ \eta \leq \frac{4 \cdot \inf_{\omega: f(\omega)=\theta} \prob_\omega (L_t - L_\theta)}{(t-\theta)^2}. \]
If $\prob_\omega$ is absolutely continuous, then the expected loss can be approximated locally as
\[ \prob_\omega(L_t - L_\theta) \approx \tfrac12 \, B(\omega) \, (t-\theta)^2, \]
where $B(\omega)$ is the second derivative of $t \mapsto \prob_\omega L_t$ at the median $t=\theta$.  For $L_t(z) = |z-t|$, this second derivative is the density of $\prob_\omega$ evaluated at the median $\theta=f(\omega)$, which is non-negative.  So, if I restrict attention to those (absolutely continuous) $\prob_\omega$'s whose density at $f(\omega)$ is at least $\eps > 0$, then choosing $\eta \leq 2\eps$ will make $\eval_\theta^n$ an e-process.  %In the numerical example that follows, I take $\eps = 0.1$.  

For illustration, suppose that Your prior knowledge is an embellished version of the ``95\% sure the median is positive'' belief, i.e., 
\[ q(\theta) = \{0.05 \times 1(\theta < 0) + 1(\theta \geq 0)\} (1 + |\theta|)^{-1}, \quad \theta \in \RR. \]
That is, beyond satisfying $\lprior(\Theta > 0) = 0.95$, posited values of the median further away from 0 are strictly less plausible than values closer to 0; see the gray curve in Figure~\ref{fig:gue}.  Note that only prior knowledge about $\Theta$ is involved here, nothing is said about other features of $\prob_\Omega$.  The other curves in Figure~\ref{fig:gue} represent the unregularized and regularized e-possibility contours---i.e., $\pi_{z^n}^\eval$ and $\pi_{z^n}^{\eval \times \rho}$, respectively---as described above for three different data sets.  Each data set is based on a common sample of size $n=25$ from a Student- distribution with degrees of freedom 2; the true density function at 0 is larger than the threshold $\eps = 0.1$ that I set here for determining $\eta$.  This common data set is shifted so that the sample median is to the left, on top of, and to the right of 0, to represent different degrees of compatibility with prior knowledge.  Here, like in Section~\ref{SS:gains}, the take-away message is that regularization leads to greater efficiency gains when the data shows mild signs of incompatibility with prior knowledge, as in Panel~(a), since those parameter values mildly compatible with the data but at least mildly incompatible with the prior knowledge get downweighted by the e-possibilistic contour. 

\begin{figure}[t]
\begin{center}
\subfigure[$\text{\tt median}(z^n) = -0.5$]{\scalebox{0.6}{\includegraphics{gue_left}}}
\subfigure[$\text{\tt median}(z^n) = 0$]{\scalebox{0.6}{\includegraphics{gue_middle}}}
\subfigure[$\text{\tt median}(z^n) = 0.5$]{\scalebox{0.6}{\includegraphics{gue_right}}}
\end{center}
\caption{Plot of the e-possibility contour for three different data sets $z^n$.}
\label{fig:gue}
\end{figure}


\subsection{From uncertainty quantification to behavior}
\label{SS:euq.behavior}

There are (at least) two different behavioral aspects that can be considered.  One pertains to the subjective interpretation of the e-possibilistic IM's output as lower and upper bounds on prices You would assign to gambles.  In that context, a relevant question is if adopting such a pricing scheme put You at risk of being made a sure loser.  The answer is {\em No}---anytime validity implies no sure loss---and these details are in Appendix~\ref{S:euq.coherence}.  

Another behavioral aspect, to be explored in detail here, concerns formal decision-making and the assessment of relevant actions in terms of (lower and upper) expected loss.  For a generic action space $\AA$, let $\ell_a(\theta) \geq 0$ represent the (non-negative) 
%\footnote{Non-negativity of the loss is not necessary, it's just that the theory below is more complicated otherwise.  If the loss is bounded away from $-\infty$, then a positive constant can be added to it to make it non-negative without affecting the shape, the expected loss minimizer, etc.}) 
loss associated with taking action $a \in \AA$ when the relevant state of the world is $\theta \in \TT$.  Under ideal, perfect-information settings, where the true state of the world $\Theta$ is known, the best decision would correspond to choosing $a^\star$ to minimize the loss $a \mapsto \ell_a(\Theta)$; note that the case of maximizing utility can be recovered from this by defining utility as, say, $u_a = -\ell_a$.  Unfortunately, $\Theta$ is uncertain, so the aforementioned minimization exercise can't be carried out.  But an e-possibilistic IM $(\lPi_{z^n}^{\eval \times \rho}, \uPi_{z^n}^{\eval \times \rho})$ is available for data-driven uncertainty quantification, so I propose to mimic the von Neumann--Morgenstern proposal and evaluate the corresponding upper and lower expected loss:
\[ \uPi_{z^n}^{\eval \times \rho}(\ell_a) = \int_0^1 \Bigl\{ \sup_{\theta: \pi_{z^n}^{\eval \times \rho}(\theta) \geq s} \ell_a(\theta) \Bigr\} \, ds \qquad \text{and} \qquad \lPi_{z^n}^{\eval \times \rho} \, \ell_a = -\uPi_{z^n}^{\eval \times \rho}(-\ell_a). \]
Since the loss is non-negative, the upper and lower expectations surely exist, though they could be $+\infty$ and $-\infty$, respectively, at least for some $a$.  They're both finite at a given $a$ if $\theta \mapsto \ell_a(\theta)$ is previsable with respect to $(\lPi_{z^n}^{\eval \times \rho}, \uPi_{z^n}^{\eval \times \rho})$, e.g., if the loss is bounded.  In any case, these two functionals together determine the IM's assessment of the quality of action $a$.  Since the goal is to make the loss small in some sense, it's reasonable to define the ``best'' action, given data $z^n$, as the minimizer of the upper expected loss:
\begin{equation}
\label{eq:a.hat}
\hat a(z^n) := \arg\min_{a \in \AA} \uPi_{z^n}^{\eval \times \rho}(\ell_a). 
\end{equation}
%Of course, if there are multiple minimizers, then any one could be chosen. 
Further details can be found in \citet{denoeux.decision.2019} and \citet{imdec}.  

For a simple illustration, reconsider the Gaussian example from Section~\ref{SS:gains}, with  
\[ \eval_\theta(z^n) = (nv + 1)^{1/2} \exp\bigl\{ -\tfrac{n}{2} (\theta - \bar z_n)^2 + \tfrac12 ( \tfrac{n}{nv + 1} ) \, \bar z_n^2 \bigr\}, \quad \theta \in \TT, \]
where $v > 0$ is a variance hyperparameter that can take any value; in my numerical examples here and above, I'll take $v=10$.  Starting with vacuous prior information, so that $\rho \equiv 1$, the e-possibilistic IM's contour $\pi_{z^n}^\eval$ is plotted in Figure~\ref{fig:contour.gauss}.  I take the loss function to be squared error: $\ell_a(\theta) = (\theta-a)^2$, also plotted in Figure~\ref{fig:contour.gauss} for several values of $a$.  It's easy to see, based on the symmetry of both the e-process and the loss, that the minimizer of the upper expected loss is $\hat a(z^n) = n^{-1} \sum_{i=1}^n z_i$, the sample mean.  For reasons that will become clear below, it's of interest to evaluate that upper expected loss at the aforementioned minimizer (see Appendix~\ref{A:choquet} for details):
\begin{align*}
\uPi_{z^n}^{\eval} \, \ell_{\hat a(z^n)} & =  \int_0^1 \Bigl\{ \sup_{\theta: \pi_{z^n}^{\eval}(\theta) \geq s} \ell_{\hat a(z^n)}(\theta) \Bigr\} \, ds = n^{-1}\bigl\{ 2 + \log(nv + 1) + (\tfrac{n}{nv+1}) \bar z_n^2 \bigr\}.
\end{align*}
%The detailed calculation is presented in {\color{red}Appendix~A}. 
%Note that, for fixed $v$, as a function of $n$, the asymptotic order of the upper expected loss is $n^{-1} \log n$.  This is within a log-factor of the usual rate $n^{-1}$ associated with the Bayes or frequentist minimax risks---the log-factor is the price one pays for the benefits of anytime validity.  

Stick with the same Gaussian illustration, but this time consider regularization corresponding to the partial prior described in Section~\ref{SS:gains}, the strongest of the three partial priors consider there.  Figure~\ref{fig:contour.gauss.reg}(a) shows the possibility contours for the four different values of the partial prior hyperparameter as in Section~\ref{SS:gains}, along with the vacuous prior contour (same as in Figure~\ref{fig:contour.gauss}).  As expected, the regularized contours are a bit more concentrated than the unregularized contour, with more concentration corresponding to the stronger prior knowledge.  From this picture, it's pretty clear that the regularized e-possibilistic IM's risk-minimizing action is closer to 0 than that for the unregularized possibilistic IM.  The introduction of regularization complicates some of the calculations, so I opt for a numerical solution.  Figure~\ref{fig:contour.gauss.reg}(b) shows plots of the e-possibilistic IM's risk assessment as a function of the action $a$, again for different values of the partial prior hyperparameter.  This confirms the intuition that the more the partial prior supports ``$\Theta$ near 0,'' the closer the IM's risk-minimizing rule will be to 0. 

\begin{figure}[t]
\begin{center}
\scalebox{0.6}{\includegraphics{nm_unreg_x05}}
\end{center}
\caption{Plot of the (unregularized) e-possibilistic IM's contour (solid), based on data $z^n$ with $n=5$ and $\bar z_n = 0.5$. Plots of the squared error loss function for several different values of $a$ are overlaid (dotted lines), including for $\hat a = \bar z_n$ (dashed line).}
\label{fig:contour.gauss}
\end{figure}

\begin{figure}[t]
\begin{center}
\subfigure[Contour]{\scalebox{0.6}{\includegraphics{nm_reg_x05}}}
\subfigure[IM's risk assessment]{\scalebox{0.6}{\includegraphics{nm_riskplot_x05}}}
\end{center}
\caption{Panel (a): Plots of the (regularized) e-possibilistic IM's contour, unregularized and regularized, based on data $z^n$ with $n=5$ and $\bar z_n = 0.5$; colors correspond to those in Figure~\ref{fig:gauss}. Plots of the squared error loss function for several different values of $a$ are overlaid (dotted lines), including for $\hat a = \bar z_n$ (dashed line). Panel (b): Plots of the e-possibilistic IM's risk assessment $a \mapsto \uPi_{z^n}^{\eval \times \rho}\,\ell_a$ for five different partial priors---the vacuous prior (black) is minimized at $\bar z_n = 0.5$ while the others shrink toward 0.}
\label{fig:contour.gauss.reg}
\end{figure}

%The IM framework offers a new assessment of various actions.  
%What does this new framework for assessment have to offer?  
%The results in Section~\ref{SSS:euq.coherence} would be relevant here, just as they're relevant in the traditional Bayesian case, but I'll not focus on this here.  
Given the IM's strong reliability properties presented above, it makes sense to ask if these carry over to formal decision-making.  The next result establishes that, indeed, the IM's assessment of an action is expected to not be too optimistic compared to the oracle's. This generalizes a result in \citet{grunwald.epost} for capped e-posteriors.

\begin{thm}
\label{thm:imdec}
Given a loss function $\ell \geq 0$, the the regularized e-possibilistic IM's assessment of actions satisfies 
\begin{equation}
\label{eq:dec.bound}
\uuprob\Bigl\{ \sup_{a \in \AA} \frac{\ell_a(f(\Omega))}{\uPi_{Z^N}^{\eval \times \rho}(\ell_a)} \Bigr\} \leq 1 \quad \text{for all $N$}. 
\end{equation}
In the special case of vacuous partial prior information, \eqref{eq:dec.bound} specializes to 
\[ \sup_{\omega \in \OO} \E_\omega \Bigl\{ \sup_{a \in \AA} \frac{\ell_a(f(\omega))}{\uPi_{Z^N}^{\eval}(\ell_a)} \Bigr\} \leq 1 \quad \text{for all $N$}. \]
\end{thm}

%\begin{proof}
%Write $\Theta$ instead of $f(\Omega)$ for now.  First, observe that the integrand 
%\[ s \mapsto \sup\{\ell_a(\theta): \pi_{Z^N}^{\eval \times \rho}(\theta) \geq s\} \]
%in $\uPi_{Z^N}^{\eval \times \rho}(\ell_a)$ is non-decreasing.  Next, I proceed by considering two separate cases.  
%\begin{enumerate}
%\item Case: $\eval^\text{reg}(Z^N,\Theta) > 1$.  In this case,  $\pi_{Z^N}^{\eval \times \rho}(\Theta) = 1$ so, by the monotonicity property mentioned above, the pair $(Z^N,\Theta)$ is such that 
%\[ \sup_{\theta: \pi_{Z^N}^{\eval \times \rho}(\theta) \geq s} \ell_a(\theta) \geq \ell_a(\Theta) \quad \text{for all $s \in [0,1]$}. \]
%Of course, if the integrand is bounded below, then the integral, over all of $[0,1]$, is also bounded below by the same value.  Therefore, 
%\[ \uPi_{Z^N}^{\eval \times \rho}(\ell_a) \geq \ell_a(\Theta), \]
%which implies 
%\begin{equation}
%\label{eq:risk.bound}
%\frac{\ell_a(\Theta)}{\uPi_{Z^N}^{\eval \times \rho}(\ell_a)} \leq 1 < \eval^\text{reg}(Z^N,\Theta). 
%\end{equation}
%\item Case: $\eval^\text{reg}(Z^N,\Theta) \leq 1$.  In this case, $\pi_{Z^N}^{\eval \times \rho}(\Theta) \leq 1$, so I can lower-bound the upper expected loss by truncating the range of integration as follows:
%\begin{align*}
%\uPi_{Z^N}^{\eval \times \rho}(\ell_a) & = \int_0^1 \Bigl\{ \sup_{\theta: \pi_{z^n}^{\eval \times \rho}(\theta) \geq s} \ell_a(\theta) \Bigr\} \, ds \\
%& \geq \int_0^{\pi_{Z^N}^{\eval \times \rho}(\Theta)} \Bigl\{ \sup_{\theta: \pi_{z^n}^{\eval \times \rho}(\theta) \geq s} \ell_a(\theta) \Bigr\} \, ds \\
%& \geq \pi_{Z^N}^{\eval \times \rho}(\Theta) \, \ell_a(\Theta),
%\end{align*}
%where the last inequality is again by the monotonicity property highlighted above.  On rearranging, and using the fact that $\pi_{Z^N}^{\eval \times \rho}(\Theta) = \eval_\Theta(Z^N) \, \rho(\Theta)$ in this case, I get 
%\[ \frac{\ell_a(\Theta)}{\uPi_{Z^N}^{\eval \times \rho}(\ell_a)} \leq \eval^\text{reg}(Z^N,\Theta), \]
%which is the same bound as in \eqref{eq:risk.bound}.  
%\end{enumerate} 
%Next, it's clear that the common bound derived in the two separate cases above holds uniformly in the actions, i.e., 
%\[ \sup_{a \in \AA} \frac{\ell_a(\Theta)}{\uPi_{Z^N}^{\eval \times \rho}(\ell_a)} \leq \eval^\text{reg}(Z^N,\Theta). \]
%Plugging $f(\Omega)$ back in for $\Theta$, taking $\uuprob$-expectation on both sides, and then applying \eqref{eq:reg.ville0}, establishes the bound in  \eqref{eq:dec.bound}. 
%\end{proof}

My first remark concerns the interpretation of Theorem~\ref{thm:imdec}: there is no (possibly data-dependent) action $a$ such that the e-possibilistic IM's assessment of $a$ doesn't tend to be overly optimistic (i.e., significantly smaller) than that of an oracle who knows the true state of the world $\Omega$.  To understand this, consider the oracle and IM assessments: $a \mapsto \ell_a(\Theta)$ and $a \mapsto \uPi_{Z^N}^{\eval \times \rho}(\ell_a)$. Imagine a plot of these two functions.  The former, the oracle assessment, will take small values for $a$ near the unattainable ``best'' action $a^\star = \arg\min_a \ell_a(\Theta)$.  Since the data's informativeness is limited, it's unrealistic to expect that the IM's assessment would tend to be small near $a^\star$ too.  Consequently, if the IM's assessment of some $a$ is much more optimistic than the oracle's, then it's likely that this $a$ is different from $a^\star$.  And if the IM's assessment is favoring actions away from $a^\star$, then there'd be a risk of suffering a large loss by taking actions suggested by the IM, e.g., $\hat a$ in \eqref{eq:a.hat}.  Theorem~\ref{thm:imdec} excludes this possibility, so the IM-based assessment is reliable.  

Secondly,the uniform bound on an expectation can be turned into an uniform probability bound.  In particular, 
\begin{equation}
\label{eq:dec.ville}
\uuprob\Bigl\{ \sup_{a \in \AA} \frac{\ell_a(f(\Omega))}{\uPi_{Z^N}^{\eval \times \rho}(\ell_a)} \geq \alpha^{-1} \Bigr\} \leq \alpha \quad \text{for all $\alpha \in [0,1]$, all $N$}. 
\end{equation}
%Compare this to the result in, e.g., Theorem~2 of \citet{imdec}.  
The above says that existence of a triple $(Z^N, \Theta, a)$ such that the IM's risk assessment is drastically smaller than the oracle's is a $\uuprob$-rare event, hence reliability.  
%This helps drive home the point that Theorem~\ref{thm:imdec} implies the IM's risk assessment is reliable.  

Third, the supremum on the inside of the $\uuprob$-expectation implies that the fixed $a$ can be replaced by any data-dependent $a$, i.e., $a(Z^N)$, leading to 
\[ \uuprob\Bigl\{ \frac{\ell_{a(Z^N)}(f(\Omega))}{\uPi_{Z^N}^{\eval \times \rho}(\ell_{a(Z^N)})} \Bigr\} \leq 1 \quad \text{for all $N$}. \]
In particular, the above holds with IM-based rule $\hat a(Z^N)$ in \eqref{eq:a.hat}.  With this focus on a single, but data-dependent action, say, $\hat a(Z^N)$, there's a bit more that can be said concerning the probability bound \eqref{eq:dec.ville}.  Indeed, \eqref{eq:dec.ville} can be unwrapped as 
\[ \uuprob\Bigl\{ \ell_{\hat a(Z^N)}(f(\Omega)) \geq \alpha^{-1} \uPi_{Z^N}^{\eval \times \rho}\,\ell_{\hat a(Z^N)} \Bigr\} \leq \alpha. \]
That is, it's a $\uuprob$-rare event that the realized loss $\ell_{\hat a(Z^N)}(\Theta)$---corresponding to the uncertain $(Z^N,\Theta)$---incurred by taking the IM's suggested action $\hat a(Z^N)$ is a large multiple of the IM's internal, data-dependent assessment of the risk.  

%To conclude this section, I'll briefly compare the setup and result here to that in \citet{grunwald.epost}.  His capped e-posterior (Definition~1) is exactly my unregularized e-possibilistic IM contour, but he doesn't use it to develop a possibility-theoretic approach to uncertainty quantification and inference.  Instead, he proceeds in somewhat of an ad hoc manner, proposing the follow {\em e-risk assessment} of an action $a$, 
%\[ \text{\sc er}_a(Z^N) =  \sup_\theta \bigl\{ \ell_a(\theta) \, \pi_{Z^N}^\eval(\theta) \bigr\}. \]
%The right-hand side above does weight the risk relative to the capped e-posterior, but since it's not a Choquet integral, the direct connection with the von Neumann--Morganstern theory is lost.  But he does establish a result analogous to that in Theorem~\ref{thm:imdec} above: using my notation, Gr\"unwald's Proposition~2 says that, under a condition,  
%\[ \sup_{\omega \in \Omega} \E_\omega\Bigl\{  \frac{\ell_{\tilde a(Z^N)}(f(\omega))}{2 \times \text{\sc er}_{\tilde a(Z^N)}(Z^N)} \Bigr\} \leq 1, \]
%where $\tilde a(Z^N)$ is the minimizer of the e-risk assessment above.  One clear difference between my result and Gr\"unwald's is the factor ``2'' in the above display.  But how do the various risk assessments compare?  Figure~\ref{fig:im.gr} suggests the following relationships:
%\begin{equation}
%\label{eq:grunwald}
%\text{\sc er}_a(z^n) \leq \uPi_{z^n}^{\eval}\,\ell_a \leq 2 \times \text{\sc er}_a(z^n), \quad a \in \AA. 
%\end{equation}
%The first inequality holds in some generality,\footnote{That is, if $\theta \mapsto \ell_a(\theta)$ is convex and if the level sets $\{\theta: \pi_{z^n}^\eval(\theta) \geq s\}$, for $s \in [0,1]$ are convex, then the first inequality in \eqref{eq:grunwald} holds.} but I've not been able to confirm the second inequality beyond just with plots like the one here.  If the second inequality also holds, then my results are stronger than Gr\"unwald's in the sense that it's easier to control a ratio if his denominator is bigger than mine.  The other difference is that his result holds only under what he calls ``Condition~0,'' but I won't discuss this here.  The point of this comparison is simply that the results obtained here and in \citet{grunwald.epost} are similar, but my e-possibilistic IM result holds with and without regularization, and is all done within a proper, coherent, possibility-theoretic framework.  
%
%\begin{figure}[t]
%\begin{center}
%\scalebox{0.65}{\includegraphics{im_gr_risk}}
%\end{center}
%\caption{Plots of the (unregularized) e-possibilistic IM's risk assessment $a \mapsto \uPi_{z^n}^{\eval}\,\ell_a$ (black) and Gr\"unwald's two risk assessments $a \mapsto \text{\sc er}_a(z^n)$ (red) and $a \mapsto 2 \times \text{\sc er}_a(z^n)$ (red dashed) based on data $z^n$ with $n=5$ and $\bar z_n = 0.5$.}
%\label{fig:im.gr}
%\end{figure}











%One notable observation is that my proposed possibilistic e-UQ can be used for safely reliable decision-making.  Let $\ell_a(\theta) \geq 0$ denote the loss associated with taking action $a \in \mathbb{A}$ when the world is in state $\theta$.  Then the possibilistic IM's upper expected loss associated with an action $a$ is a Choquet integral, which is given by 
%\[ \uPi_{z^n}^\ee \, \ell_a = \int_0^1 \sup_{\theta: \pi_{z^n}^\ee(\theta) > s} \ell_a(\theta) \, ds. \]
%The claimed reliability of my IM's upper expected loss corresponds to the following extension of the results in \citet{grunwald.epost, imdec}:
%\[ \sup_N \sup_{\omega \in \OO} \E_\omega\left\{ \sup_{a \in \mathbb{A}} \frac{\ell_a(f(\omega))}{\uPi_{Z^N}^\ee \, \ell_a} \right\} \leq 1. \]
%In words, the IM's assessment of an action $a$ won't be more optimistic than that of an oracle who knows the true $\theta=f(\omega)$, and this holds uniformly over actions, true states of the world, and stopping rules. More details will be presented elsewhere.


%Open questions include how to incorporate incomplete or partial prior information about $\Theta = f(\Omega)$ and/or about the stopping rule $N$.  Indeed, there's an opportunity for efficiency gain if we don't need to be safe relative to {\em all} stopping rules.  With the inclusion of partial prior information, one could try to generalize the above upper expected loss result and, moreover, try to prove a version of the classical Bayesian result, namely, that the possibilistic IM's optimal action---the one that minimizes the upper expected loss $a \mapsto \uPi_{z^n}^\ee \ell_a$---is admissible.  It's relatively straightforward to incorporate partial prior information coming in the form of a possibility measure, and I'll report on this elsewhere.  But given the connection to e-values, game-theoretic probability, etc., it's of interest to consider prior information that takes the form of gambles about the uncertain value of $\Theta$ that are {\em a prior} acceptable to the data analyst, i.e., gambles that aren't expected to multiply an opponent's capital by a large factor. 



%\subsection{Efficiency gains, cont.}



\section{Application}
\label{S:ware}

Statistical, scientific, and ethical challenges emerge when comparing two medical treatments when one is potentially far superior to the other.  The goal is to demonstrably prove that the one treatment is superior, so that the inferior treatment can be safely removed from use, thereby saving/improving lives.  On the one hand, if one treatment is far superior, then, this should be easy to prove statistically with sufficient data.  On the other hand, if one treatment is far superior, then it's borderline unethical
%---the Hippocratic Oath says ``Do no harm''---
to continue randomizing patients to the inferior treatment just to get ``sufficient data.'' 
%with the planned randomized trial that exposes some patients to the inferior treatment with much higher mortality rates.  
This context clearly highlights the importance of (a)~accommodating adaptive data-collection schemes and (b)~using reliable and efficient statistical methods from which justifiable conclusions can be drawn while minimizing patients' exposure to an inferior treatment.  

One well-known study is presented in \citet{ware1989}.  His investigation concerned persistent pulmonary hypertension in newborns.  The standard treatment for many years, called the conventional medical therapy (CMT), had a mortality rate of at least 80\%.  But by 1985, a new potential treatment had emerged, namely, extracorporeal membrane oxygenation (ECMO), for which reported survival rates were at least 80\%.  Unfortunately, the empirical support for this latter claim was rather thin: only one randomized trial had been performed and, in that trial, only one patient was assigned to CMT.  Despite ECMO's strong performance in these studies, the combination of its statistically-limited support and its potential for side-effects gave some medical researchers pause.  Ware's paper describes a trial that's based on randomization (within blocks) until a prespecified number of deaths---namely, 4---are observed in {\em each group}.  Ware's data is as follows:
\begin{equation}
\label{eq:ware.data}
\text{CMT: 10 patients with 4 deaths} \qquad \text{ECMO: 9 patients with 0 deaths}. 
\end{equation}
%I'll focus my analysis below on the data in \eqref{eq:ware.data}.  
It's important to note that the investigators {\em didn't carry out the design as planned}, i.e., they didn't continue following patients until a fourth death in the ECMO group was observed.  This emphasizes the need for anytime valid statistical procedures. 
%\footnote{It's worth noting that Ware's team did carry out a second, non-randomized phase, in which another 20 patients were given ECMO and only 1 died.  But there were other changes made in the patient admission process, which created some concerns about homogeneity between the first and second phases of the trial.  As such, like Ware and also Walley, I'm excluding the Phase~2 data from my analysis.}  


Assuming independence, and that the uncertain survival probabilities---denoted by $\Theta_\text{cmt}$ and $\Theta_\text{ecmo}$---are constant across patients, the likelihood function based on \eqref{eq:ware.data} is 
\[ L_{z^n}(\theta) \propto \theta_\text{cmt}^6 \, (1 - \theta_\text{cmt})^4 \, \theta_\text{ecmo}^9, \qquad \theta=(\theta_\text{cmt}, \theta_\text{ecmo}), \]
where ``$z^n$'' is just the symbol I'll use for the data in \eqref{eq:ware.data}.  For the e-process, I'll use a slightly modified version of the recommendation in \citet{turner.grunwald.2023},
\[ \eval_\theta(z^n) = \frac{\hat\theta_{\text{cmt},\beta}^6 \, (1 - \hat\theta_{\text{cmt},\beta})^4 \, \hat\theta_{\text{ecmo},\beta}^9}{\theta_\text{cmt}^6 \, (1 - \theta_\text{cmt})^4 \, \theta_\text{ecmo}^9}, \qquad \theta = (\theta_\text{cmt}, \theta_\text{ecmo}), \]
where 
\[ \hat\theta_{\text{cmt},\beta} = \frac{6 + \beta}{10 + 2\beta} \quad \text{and} \quad \hat\theta_{\text{ecmo},\beta} = \frac{9 + \beta}{9 + 2\beta}, \]
with $\beta=0.18$ as suggested in \citet[][Sec.~3]{turner.grunwald.2023}. The results of my analysis based on this e-process will be presented below; see, e.g., Figure~\ref{fig:ware.2d}(a). 

The reader can glean from the above discussion that, while the data in this particular study might be rather limited, there is relevant ``prior information'' available.  But what to do with it?  The frequentist analysis in \citet{ware1989} formally ignores prior information; informally, however, the prior information is used in ad hoc ways.  Bayesian solutions face the problem that prior information is insufficient to determine a precise prior distribution for $\Theta$.  A common strategy in such cases is a prior sensitivity analysis based on a few selected priors that are ``consistent'' with the information available \citep{kass.greenhouse.1989, kass1992}.  Clearly, none of these priors are ``right,'' so the best one can hope for is that the sensitivity analysis reveals that the posterior isn't sensitive to the prior.  
%But that happens only when the data is at least moderately informative, in which case, the likelihood alone is enough---the prior and Bayesian analysis isn't really needed.  
\citet[][Sec.~5]{walley1996} offers a generalized Bayesian analysis that also virtually ignores the prior information; but instead of literally ignoring it, he models the near-ignorance by a prior credal set that contains all independent beta distributions for $(\Theta_\text{cmt}, \Theta_\text{ecmo})$.  In cases like the present CMT/ECMO study, where data is limited, my goal is to leverage, rather than ignore, real-but-necessarily-incomplete prior information in order to {\em strengthen} the analysis and justification for its conclusions.  

Next, I'll present my encoding of the available partial prior information---with minimal embellishment on the information given in \citet{ware1989}---as a possibility distribution to be used in my subsequent analysis.  To be clear, I have no medical expertise; so, my analysis is sure to be overly simplistic and, hence, is only for illustrative purposes.  That said, I think my formulation and conclusions drawn are quite reasonable.  
%There's bound to be some degree of subjectivity in these considerations, but that's nothing to be afraid of; subjectivity isn't a dirty word.  In high-stakes contexts involving uncertainty, the reason we give more credibility to experts than to novices is that we trust the former's insights and judgments more than that latter's, not because the one knows more objective facts than the other. 
%One of my non-technical goals in this paper is to emphasize both the need for and potential benefits from sound subjective judgments in statistical analyses.  So, please feel free to disagree with my prior encoding, analysis, etc.~and put forward your own; I welcome the opportunity to learn.

There are two statements---``$\leq 0.2$'' and ``$\geq 0.8$''---that stand out in Ware's report, which are best interpreted as ``prior limits'' for $\Theta_\text{cmt}$ and $\Theta_\text{ecmo}$, respectively.  What's missing are quantitative statements about the confidence or degree of belief in these limits, but qualitative statements are made in the text.  This is where some subjective judgment becomes necessary.  As Ware explains, there's reasonably strong support for the prior limit ``$\Theta_\text{ecmo} \geq 0.8$'' based on historical data.  Importantly, Ware argues that the study inclusion protocol, etc.~are such that the patients in these previous studies and those in his study are more-or-less homogeneous.  The support for the prior limit ``$\Theta_\text{cmt} \leq 0.2$,'' on the other hand, is considerably weaker: there was only one case involving a patient that could've received ECMO but was randomly assigned to CMT instead.
%; this casts doubt on whether the difference between survival rates in prior studies can be solely attributed to the CMT/ECMO treatment.  
Based on my interpretation of Ware's report, I assign confidence as follows:
\begin{itemize}
\item 90\% confident in ``$\Theta_\text{ecmo} \geq 0.8$'' and
\vspace{-2mm}
\item 50\% confident in ``$\Theta_\text{cmt} \leq 0.3$.'' 
\end{itemize}
Aside from the difference in confidence levels, which I'll explain shortly, notice that I stretch out the limit for $\Theta_\text{cmt}$ a bit; this is because Ware used ``0.2'' as a sort of prior mode for $\Theta_\text{cmt}$, so it might be prudent to have the mode closer to the middle of the range to which I assign some non-trivial degree of confidence.  The confidence levels above are based on my judgment that researchers are quite confident in the performance of ECMO but far less---say, roughly {\em half} as---confident about CMT on similar populations of patients.  Mathematically, I opt to encode this vaguely-stated quantification of uncertainty as a possibility measure for $\Theta=(\Theta_\text{cmt}, \Theta_\text{ecmo})$.  Marginally, the two contours are
\begin{align*}
q_\text{ecmo}(\theta_\text{ecmo}) & = 0.1 + 0.9 \cdot 1(\theta_\text{ecmo} \geq 0.8) \\
q_\text{cmt}(\theta_\text{cmt}) & = 0.5 + 0.5 \cdot 1(\theta_\text{cmt} \leq 0.3).
\end{align*}
The aforementioned ``confidence levels'' correspond to the properties 
\[ \uprior_\text{cmt}(\Theta_\text{cmt} \leq 0.3) = 0.5 \quad \text{and} \quad \lprior_\text{ecmo}(\Theta_\text{ecmo} \geq 0.8) = 0.9. \]
Following Walley and others, I'll treat $\Theta_\text{cmt}$ and $\Theta_\text{ecmo}$ as independent {\em a priori} and then take the joint possibility measure $\uprior$ for $\Theta=(\Theta_\text{cmt}, \Theta_\text{ecmo})$ to have contour $q$ equal to the product of $q_\text{ecmo}$ and $q_\text{cmt}$ above. This $q$ is a piecewise constant function on the unit square, taking value 1 on the rectangle $[0,0.3] \times [0.8, 1]$ and smaller values in the three other rectangles.  While others might disagree to some extent with the particular confidence levels that I chose, I believe that this possibilistic prior is consistent with the information Ware presented.  In particular, the credal set determined by $\uprior$ contains independent products of certain beta distributions, among others.  

Plots of the unregularized and regularized e-processes are shown in Figure~\ref{fig:ware.2d}; for the regularized e-process, I'm using the same calibrator $\gamma$ as described in Section~\ref{SS:possibilistic.case}.  As expected, the e-process contour bottoms out at the value corresponding to the simple sample proportions, namely, $(0.6, 1)$.  
%As $\theta$ moves away from that minimizer, the e-process increases and creates contours of a shape  determined by the data and the specific e-process formulation.  
Note that the unregularized e-process contours are smooth, whereas those of the regularized version's are rough in some places; this roughness is due to the relatively large jump discontinuity in the prior contour there.  The heavy red line marks the confidence sets $C_\alpha$ and $C_\alpha^\text{reg}$ for $\alpha=0.05$, and there are two notable observations.  First, as it pertains to $\Theta_\text{ecmo}$, for which prior information is relatively strong, the limits are much tighter in the regularized case compared to the unregularized.  Second, for $\Theta_\text{cmt}$, the limits are a bit looser for the regularized case compared to the unregularized.  The latter point might seem disappointing, but this is exactly what should happen: prior knowledge that's somewhat incompatible with data ought to result in more conservative inference.  

\begin{figure}[t]
\begin{center}
\subfigure[Unregularized]{\scalebox{0.55}{\includegraphics{ware_eprocess}}}
\subfigure[Regularized]{\scalebox{0.55}{\includegraphics{ware_reprocess}}}
\end{center}
\caption{Plots of the unregularized and regularized e-processes based on Ware's CMT/ECMO data.  Heavy red lines mark the corresponding 95\% confidence sets.}
\label{fig:ware.2d}
\end{figure}

A relevant question is if EMCO is more effective than CMT.  One way to answer this is to test the hypothesis ``$\Theta_\text{ecmo} \leq \Theta_\text{cmt}$'' versus ``$\Theta_\text{ecmo} > \Theta_\text{cmt}$.'' The diagonal line through the two plots in Figure~\ref{fig:ware.2d} represents the boundary between these two propositions and, since the red contour curves intersect with the diagonal line, the corresponding e-process-based tests cannot reject the hypothesis ``$\Theta_\text{ecmo} \leq \Theta_\text{cmt}$'' at level $\alpha=0.05$.
%; but they could reject at a slightly larger $\alpha$ levels.  
For comparison, Ware presents an analysis that using Fisher's exact test, leading to a p-value of 0.054, which is consistent with my conclusions based on Figure~\ref{fig:ware.2d}(a).  My results are based on procedures proved to be anytime valid, so this provides additional comfort given that Fisher's exact test is not anytime valid.
%---it's conditional on relevant marginal counts.  

For uncertainty quantification, the possibility contours, $\pi^\eval$ and $\pi^{\eval \times \rho}$, in this case look identical to those in Figure~\ref{fig:ware.2d}, just with different numerical labels on the contours.  So, the confidence regions and test conclusions derived from $\pi^\eval$ and $\pi^{\eval \times \rho}$ give exactly the same results as those obtained from Figure~\ref{fig:ware.2d}.  But there's more that can be done with the IM's possibilistic uncertainty quantification.  First, consider a feature $\Delta = \Theta_\text{ecmo} - \Theta_\text{cmt}$, the difference in survival rates between EMCO and CMT.  The possibility-theoretic {\em extension principle} 
%\citep[e.g.,][]{zadeh1975a} 
determines a (marginal) possibility contour 
%\footnote{Note, this is the Greek letter ``upsilon,'' not the Latin letter $v$ or the Greek letter $\nu$.} 
for $\Delta$ from that for $\Theta$:
\[ \phi_{z^n}^\eval(\delta) = \sup_{\theta \in [0,1]^2: \, \theta_\text{ecmo}-\theta_\text{cmt}=\delta} \pi_{z^n}^\eval(\theta), \quad \delta \in [-1,1]. \]
The regularized e-possibilistic marginal IM contour $\phi_{z^n}^{\eval \times \rho}$ is defined analogously. The corresponding upper probabilities $\uPhi_{z^n}^\eval$ and $\uPhi_{z^n}^{\eval \times \rho}$ are defined via optimization as usual.  The dashed line at $\alpha=0.05$ determines a marginal anytime valid confidence interval for $\Delta$ and, at least for the regularized version (in red), this more-or-less agrees with the generalized Bayes 95\% credible interval presented in \citet[][Fig.~1]{walley1996}.  The message from this plot is that, while ``$\Delta > 0$'' is highly possible, nearly the entire range $[-1,1]$ is ``sufficiently possible'' based on Ware's data.  Admittedly, the extension principle is conservative, but it's the most direct way to marginalize while preserving validity.
%, without going back to the original data and formulating a e-process directly for $\Delta$.  

\begin{figure}[t]
\begin{center}
\scalebox{0.65}{\includegraphics{ware_marginal}}
\end{center}
\caption{Plots of the regularized and unregularized e-possibilistic IM's marginal contour for $\Delta = \Theta_\text{ecmo} - \Theta_\text{cmt}$ based on Ware's CMT/ECMO data.}
\label{fig:ware.1d}
\end{figure}

A formal decision-theoretic approach can be considered, using suitable lower and upper expected loss/utility based on the e-possibilistic IM output.  \citet{walley1996} argues that of utmost importance---or {\em utility}---is a patient's survival.  That is, the desired goal is for the patient to survive with the treatment they were given, all other factors are irrelevant.  So, if $Y$ is a binary indicator that the patient in question survives, then the utilities associated with CMT and ECMO are 
\[ u_\text{cmt}(y) = u_\text{ecmo}(y) = y, \quad y \in \{0,1\} \equiv \{\text{dies}, \text{survives}\}. \]
In this case, the expected utilities are $\Theta_\text{cmt}$ and $\Theta_\text{ecmo}$.  Since $(\Theta_\text{cmt}, \Theta_\text{ecmo})$ is uncertain, I can't simply pick which of CMT and ECMO has higher expected utility.  But I can evaluate the lower and upper expectation of the difference between the two expected utilities, namely, $\Theta_\text{ecmo} - \Theta_\text{cmt}$, based on the full uncertainty quantification provided by the e-possibilistic IM.  I've already obtained the marginal IM for the difference $\Delta$ in expected utilities, so I just need to evaluate the lower and upper expectations.  For the unregularized e-possibilistic IM, the upper and lower expectation of $\Delta$ are 
\begin{align*}
\uPhi_{z^n}^\eval(\Delta) = \int_0^1 \sup\{ \delta: \phi_{z^n}^\eval(\delta) > s\} \, ds \quad \text{and} \quad \lPhi_{z^n}^\eval(\Delta) = -\uPhi_{z^n}^\eval(-\Delta), 
\end{align*}
and analogously for the regularized version.  Numerically, I get:
\[ \bigl[ \lPhi_{z^n}^\eval(\Delta), \uPhi_{z^n}^\eval(\Delta) \bigr] = [-0.042, 0.739] \quad \text{and} \quad \bigl[ \lPhi_{z^n}^{\eval \times \rho}(\Delta), \uPhi_{z^n}^{\eval \times \rho}(\Delta) \bigr] = [0.086, 0.820]. \]
That the first, unregularized interval contains 0 means that I can't rule out, based on Ware's data, that CMT is preferred to ECMO in terms of utility.  The regularized interval, however, is strictly to the right of the origin, which implies that ECMO is the preferred treatment in terms of expected utility, albeit just barely.  For comparison, Walley reports his version the ``lower and upper expectation of $\Delta$'' as $[0.152, 0.5]$, so he too concludes that ECMO is the demonstrably better treatment.  
%One has to decide whose analysis is more convincing, mine or Walley's, but my IM solution actually uses prior information and my expected utility offers a strong, decision-making reliability guarantee via Theorem~\ref{thm:imdec}. 


%{\color{red} In this case, it might not be necessary for the analysis to be ``anytime valid'' in the sense above.  Talk about possible improvements if the particular study design is taken into account, but highlight the utility of ``anytime validity'' for clinical trials involving dynamic stop-continue decisions, etc...} 



%\section{Extensions}
%\label{S:extensions}
%
%Here I'd like to briefly mention a few possible extensions of the formulation here.  Some of these I more-or-less know how to do, and others I don't.  
%
%First, I've focused here on the case where, given $\Omega$ (or $\prob$), the data are iid.  There's no conceptual difficulty in extending this to the case of independent but not iid, e.g., if data points depend on covariates.  This case was treated informally in Section~\ref{S:ware} where the treatment assignment---CMT versus ECMO---plays the role of a covariate.  Dependent data sequences is a bigger challenge, a topic for future investigation.  
%
%%Second, my general setup here is quite accommodating, but my illustrations focused on cases involving a parametric statistical model.  Outside this case, there are new e-process developments in, e.g., \citet{dey.gue}, that take care of one part of the regularized e-process construction when $\Theta$ can be interpreted as a risk minimizer.  An understanding of how to elicit partial prior information for risk minimizers (see Section~\ref{S:discuss} below) and empirical investigations into how the corresponding regularized e-processes perform, again, are topics for future investigation. 
%
%Second, the literature on anytime valid inference---including this paper---has focused exclusively on the case where $N$ can be literally anything (subject to the mathematical constraint that it be a stopping time).  That's why the above results all have ``for all $N$'' appended to the statements.  In applications, however, it's hard to imagine that literally {\em any} stopping rule is in consideration.  If some stopping rules can be excluded, so that ``for all $N$'' can be replace by ``for such-and-such $N$,'' then that creates an opportunity to leverage those constraints and improve the (regularized) e-process's efficiency.  Even more generally, You might not know which $N$ was actually used, but You might be able to quantify Your uncertainty with a possibility contour or other kind of imprecise probability.  The relevant questions are if partial prior information about $N$ can be incorporated; if so, then how; and then what kind of efficiency gains can be expected. 
%
%Finally, while the problem formulation is quite general, my illustrations have only considered cases involving uncertain $\Theta$ with fixed dimension.  In anticipation of applications to high-dimensional problems, it's important to recognize that an efficiency gain can be realized only if the regularization is allowed to adapt to the dimension of $\Theta$.  In the classical statistics and machine learning literature, this is accommodated by tying the dimension to the sample size and then allowing the penalty/prior to depend explicitly on the sample size.  This is a little awkward in the present context, for multiple reasons.  Perhaps the biggest challenge is that the data are streaming, and the stopping rule is dynamic, so if dimension and sample size are linked, then it's as if the inference problem doesn't have a well-defined target.  Another challenge is dealing with the fact that $\credal$ is assumed to be a {\em real} representation of Your partial knowledge and it might be awkward, at least initially, to think about knowledge depending on dimension.  As I see it, however, the intuition behind standard penalties/priors---e.g., ``most of the entries in $\Theta$ are near-zero''---is inherently imprecise/possibilistic and notions like ``most'' are relative to the number of things in question, which in this case is the dimension of $\Theta$.  This seems doable, but the question remains: specifically how should $\credal$ adapt to dimension? 





\section{Conclusion}
\label{S:discuss}

This paper develops the new concept of and theory associated with {\em regularized e-processes}, from which I further develop reliable---i.e., anytime valid and efficient---inference and uncertainty quantification.  On the technical side, the regularized e-process involves the combination of a fully data-driven e-process and prior knowledge about $\Theta$ that You, the investigator, might have.  Importantly, Your prior information need not be---and typically won't be---sufficiently complete to pinpoint a single prior as a Bayesian analysis would require, so my proposal explicitly draws on aspects of imprecise probability theory.  My proposed regularization discounts those values of $\Theta$ that are incompatible with the available prior knowledge, making it easier to ``reject'' such values compared to the purely data-driven e-process, hence the efficiency gains.  The critical point, however, is that any such non-trivial discounting like described above jeopardizes the original e-process's inherent anytime reliability; therefore, the proposed regularization requires care.  

By encoding the partial prior information rigorously as a coherent imprecise probability, I can build a collection of joint distributions for $(\text{data}, \Theta)$ compatible with the assumed model and available prior information.  Then I generalize the now-familiar anytime validity property in a sound way, accounting for the available prior information, and similarly generalize Ville's inequality to prove that the proposed regularized e-process remains anytime valid in this more general sense.  This reformulation offers mathematical justification for You to accept the efficiency gains offered by regularization.  

%From standard frequentist-style of inference based on (anytime) valid tests and confidence sets, I go on to develop a Bayesian-like framework for broader uncertainty quantification based on e-processes.  This falls under the general umbrella of {\em inferential models}, or IMs.  The proposed brand of uncertainty quantification is possibility-theoretic, which means that the proposed e-possibilistic IM operates with a different albeit simple and intuitive calculus: optimization instead of integration.  Aside from these small technical differences, the approach proposed here at least superficially resembles Bayesian inference.  What distinguishes the proposed e-possibilistic IM framework is that it exactly (not just asymptotically) enjoys the calibration necessary to ensure that the uncertainty quantification derived from it is reliable.  This kind of reliability immediately implies that certain summaries of the e-possibilistic IM provide anytime valid test and confidence set procedures; it also implies more than the usual frequentist guarantees, e.g., reliable probing for hypotheses supported by data (Corollaries~\ref{cor:uniform.vac} and \ref{cor:strong.valid}).  Remarkably, the e-possibilistic IM construction leads to a number of other desirable properties, things often associated with Bayesian inference.  This includes: satisfying the likelihood principle, avoiding sure-loss, and formal decision-making with strong reliability guarantees.  

%An illustration of the proposed framework in the context of a clinical study was presented in Section~\ref{S:ware}.  While this was mainly just to illustrate what the proposed framework has to offer, the results were reasonable and computation was completely trivial.  A next step is to consider other applications, of greater complexity, to build up experience and showcase what regularized e-processes can do.  Since having incomplete prior knowledge is the norm, and not the exception, the proposal here should immediately replace the tools that practitioners are using on a daily basis, ones that require investigators to unjustifiably add to or delete from their corpus of knowledge.  

A number of interesting and challenging questions remain open, I'll end here with a few remarks about two of them.  First, to put the theory presented here into practice, You must convert what You know about the uncertain $\Theta$ into an imprecise probability.  For example, in the clinical trial data illustration of Section~\ref{S:ware}, some degree of belief attached to bounds on $\Theta$ was converted into a possibility contour.  The point is that You need to know (a)~how to tease relevant quantitative details out of Your mostly qualitative corpus of knowledge and (b)~how to properly map those relevant details to a suitable imprecise probability.  On top of this, the second point, is that a possibilistic formulation of the prior information requires choice of a calibrator as I described in Section~\ref{SS:possibilistic.case}.  The elicitation and encoding of partial prior information is an interesting question, more psychological than statistical, something that I plan to explore.  A thorough comparison of how different calibrators perform in this context across different model/data settings is also needed. 

%Second, the reader may have noticed in Section~\ref{SS:gains} that there are instances (i.e., model and data combinations) in which the regularized e-process based on weak prior information apparently gives slightly wider confidence limits than the unregularized e-process that assumes vacuous prior information.  This is counter-intuitive because a vacuous prior encodes strictly less information than even a weak partial prior---my intuition was that regularization would always help if data and prior are compatible.  That intuition holds if I could regularize via multiplying the e-process by the prior possibility contour's reciprocal.  But this reciprocal, on its own, can't be a regularizer.  So the relevant question is: for a given calibrator, is some prior information ``too weak'' that You're better off ignoring it than trying to pass it through a calibrator before regularization?  It's hard for me to accept that relevant information should ever be ignored, so let me restate the previous question from a different perspective: should the choice of calibrator depend on the strength of the partial prior information and, if so, then how?  

Second, my illustrations only considered cases involving $\Theta$ of fixed dimension.  For applications to high-dimensional problems, it's important to recognize that an efficiency gain can be realized only if the regularization adapts to the dimension of $\Theta$.  This is typically accommodated by tying the dimension and sample size together and then allowing the penalty/prior to depend explicitly on the sample size.  This is a little awkward in the present context, for several reasons.  One is that the data are streaming, and the stopping rule is dynamic, so if dimension and sample size are linked, then it's as if the inference problem doesn't have a well-defined target.  Another challenge is dealing with the fact that $\credal$ is assumed to be a {\em real} representation of Your prior knowledge and it might seem strange that Your knowledge depends on dimension.  The intuition, however, behind standard penalties/priors---e.g., ``most of the entries in $\Theta$ are near-zero''---is inherently imprecise and notions like ``most'' are relative to the number of things in question, which in this case is the dimension of $\Theta$.  This seems doable, but the question remains: specifically how should $\credal$ adapt to dimension? 


\section*{Acknowledgments}

This work is partially supported by the U.S.~National Science Foundation, grants SES--2051225 and DMS--2412628. 


%\bibliographystyle{apalike}
%\bibliography{/Users/rgmarti3/Dropbox/Research/mybib}


%\pagebreak 

\appendix


\section{Technical remarks}
\label{A:remarks}

Below are several remarks offering further insights on and explanation of some specific points made in the main text. 

\begin{remark}
\label{re:continuity}
If one encounters a rare case in which the mapping from $\Omega$ to $\Theta=f(\Omega)$ is discontinuous, then there's a simple work-around.  Just carry out the developments below as if $\Omega$ is the quantity of interest, and then marginalize down to $\Theta$ at the end, e.g., by stating the relevant hypotheses about $\Theta$ in terms of the corresponding hypotheses about $\Omega$.  All of the theory presented in the paper supports such an approach. The only downside is that carrying the higher-dimensionality of $\Omega$ from beginning to (near) end can cause efficiency to be lost compared to dropping the dimension to that of $\Theta$ at the outset. 
\end{remark}

\begin{remark}
\label{re:betting}
Following de Finetti, it is common to interpret lower and upper probabilities $(\lprior,\uprior)$ in terms of buying/selling prices for gambles:
\begin{equation}
\label{eq:prices}
\begin{split}
\lprior(H) & = \text{Your supremum buying price for $\$1(\Theta \in H)$} \\
\uprior(H) & = \text{Your infimum selling price for $\$1(\Theta \in H)$}.
\end{split}
\end{equation}
That is, You'd be willing to buy a ticket that pays \$1 if ``$\Theta \in H$'' from me for no more than $\$\lprior(H)$ and, similarly, You'd be willing to sell a ticket to me that pays \$1 if ``$\Theta \in H$'' to me for no less than $\$\uprior(H)$.  Intuitively, You shouldn't be willing to buy a gamble for a price higher than You're willing to sell it for---that would put You at risk of a sure loss---so the setup would be unsatisfactory if it didn't rule out this case.  Indeed, the assumed structure of the credal set implies that $\lprior(H) \leq \uprior(H)$ for all $H$.  Moreover, $\uprior$ is not just a summary or derivative of $\credal$, the two are in one-to-one correspondence:
\begin{equation}
\label{eq:credal}
\credal = \{\prior: \prior(H) \leq \uprior(H) \text{ for all measurable $H$} \}. 
\end{equation}
That is, $\uprior$ is exactly the upper envelope of $\credal$ and, consequently, credal set-driven imprecise probabilities are {\em coherent} \citep[e.g.,][]{walley1991}, generalizing the notion put forth by de Finetti, Savage, Ramsey, and others for precise probability.
\end{remark}

\begin{remark}
\label{re:vac.supremum}
The simplification in \eqref{eq:vac.supremum} may not be immediately obvious, so here's an explanation.  Expectation of a random variable is linear in the distribution and, therefore, it's supremum over a closed and convex set of distributions is attained at the extremes, on the boundary.  If the prior is vacuous, so that the credal set contains all distributions, then the boundary consists of point mass distributions. Hence, the supremum expectation over all distributions is attained at a point mass distribution.
\end{remark}

\begin{remark}
\label{re:prob.to.poss}
The choice of prior contour in Section~\ref{SS:gains} is an application of the general {\em probability-to-possibility transform} described in, e.g., \citet{dubois.etal.2004}, \citet{hose.hanss.2021}, and \citet{hose2022thesis}.  Mathematically, the transformation is rather straightforward.  If $Y$ is a random vector taking values in $\YY$ with distribution $\prob_Y$ and corresponding density/mass function $f_Y$, then the probability-to-possibility transform results in a possibility measure with contour defined via  
\[ \psi(y) = \prob_Y\{ f(Y) \leq f(y) \}, \quad y \in \YY. \]
That $\psi$ is a possibility contour is easy to see: if $f_Y$ has mode $y^\text{mode}$, then $\psi(y^\text{mode})=1$; the unbounded-$f_Y$ case can be handled similarly.  In Section~\ref{SS:gains} of the main text, the contour that involved the chi-square distribution function is obtained by applying the above formula with $f_Y$ a Gaussian density.  What's notable about this particular possibility measure construction is that the credal set corresponding to $\psi$ is the smallest of all those credal sets that contain $\prob_Y$ and whose upper envelope is a possibility measure.  So, if, like in Section~\ref{SS:gains}, ones prior knowledge consists of only a surprise-assessment that agrees with a Gaussian probability, then the choice of prior possibility contour recommended there is the best choice in the sense that it's corresponds to the smallest (possibilistic) credal set consistent with the available prior knowledge.  
\end{remark}

\begin{remark}
\label{re:why.possibility}
A reasonable question is: why are {\em possibilistic} IMs appropriate?  To me, the most compelling justification comes from the uniform validity result in Corollary~\ref{cor:uniform.vac}.  The equivalence---see Equation~\eqref{eq:equivalence} below---in the proof holds for all imprecise-probabilistic IMs, but what I'm calling the ``contour'' $\pi^\eval$ has properties under a possibilistic formulation that it doesn't have under other formulations.  For a generic, data-dependent upper probability $\uPi_{Z^N}$ on $\TT$, the function $\theta \mapsto \uPi_{Z^N}(\{\theta\})$ doesn't completely determine $\uPi_{Z^N}$ and, moreover, it can happen that $\theta \mapsto \uPi_{Z^N}(\{\theta\})$ is always small, no matter what $Z^N$ is.  Then the probability (with respect to the sampling distribution of $Z^N$) that it's less than some $\alpha < 1$ could be large---perhaps even equal to 1.  It's unique to the possibilistic framework that the contour fully determines the upper probability and, moreover, takes values arbitrarily close to 1.  Without this special structure, strong validity and, hence, uniform validity can't be attained.  If it could be attained by some other, non-possibilistic IM construction, then Lemma~1 in \citet{martin.partial2} says that there's a possibilistic IM that's no worse in terms of efficiency.  The conclusion is that, if strong validity and the safety it offers is a priority, which it is to me, then the possibilistic formulation is without loss of generality/efficiency.  
\end{remark}

\begin{remark}
\label{re:typically}
The function $\pi_{z^n}^\eval$ fails to be a possibility contour at a given $z^n$ if and only if $\eval_\theta(z^n)$ is strictly greater than 1 for all $\theta$.  But e-processes have expected value upper-bounded by 1, so it'd be exceptionally rare, although not impossible, for $z^n$ to not be particularly compatible with any $\theta$, so that $\eval_\theta(z^n)$ is everywhere greater than 1 and, hence, $\pi_{z^n}^\eval$ is everywhere below 1. This doesn't affect the statistical properties, only the interpretation of the IM; see Appendix~\ref{S:euq.coherence} below. 
\end{remark}



\section{Choquet integration}
\label{A:choquet}

%IMPORTANT -- MOVE TO APPENDIX?
%Roughly, the {\em upper envelope theorem} in \citet[][Theorem~4.38]{lower.previsions.book} connects the bounds in \eqref{eq:envelope} to the so-called {\em natural extension} of $\uprior$, and then their Theorem~6.14 links the natural extension to {\em Choquet integration} \citep{choquet1953}.  I'll present the Choquet integral formula for one special case below; see Appendix~A for further details.

Choquet integration has an important role to play in the main paper's developments.  Basically, Choquet integration plays the same role in imprecise probability theory as Lebesgue integration does in ordinary/precise probability theory.  That is, just as Lebesgue integration is the go-to mathematical framework for defining expectation with respect to probability measures, Choquet integration is the appropriate way to extend lower/upper probabilities, at least for the kind of models in consideration here, to more general lower/upper expectations.  The path to making this connection isn't exactly direct, and the details are too involved to present here, but I think a relatively brief overview would be beneficial.  My summary here is based primarily on details presented much more thoroughly and rigorously in \citet{lower.previsions.book}.  

Let $g: \TT \to \RR$ be a function, which I'll assume to be non-negative only for simplicity; if $g$ can take both positive and negative values, then one apply the developments here to the difference between the positive and negative parts of $g$.  If $\uprior$ is a general capacity---a normalized, monotone set function---supported on subsets of $\TT$, then the Choquet integral of $g$ with respect to $\uprior$ is defined as 
\begin{equation}
\label{eq:choquet}
\mathcal{I}_\text{\sc choq}(g) := \int_0^\infty \uprior\{ \theta \in \TT: g(\theta) \geq t \} \, dt, 
\end{equation} 
where ``$\int$'' on the right-hand side is a Riemann integral, which is well-defined since the integrand is a monotone non-increasing function of $t$.  In some sense, there's nothing particularly special about defining an ``integral''---anyone can do it.  The challenge is defining an integral that represents something relevant.  In the present context, the most meaningful notion of an expected value of $g(\Theta)$ with respect to a coherent upper probability $\uprior$ (defined on all subsets $H$ of $\TT$) is the upper envelope 
\[ \uprior\,g = \sup_{\prior \in \credal} \E^{\Theta \sim \prior}\{ g(\Theta) \}, \]
which is the second expression given in Equation~(6) of the main paper; recall that $\credal$ here is the set of all probabilities dominated by $\uprior$, as in Equation~(5).  How is this connected to the Choquet integral?

What links the Choquet integral above to the upper expectation is a deep result of Walley's, concerning the so-called {\em natural extension} of $\uprior$ from an upper probability to an upper expectation.  In a purely mathematical sense, one may have a function $f$ defined on a domain $\XX$ with certain properties, and the relevant question is if $f$ can be extended from $\XX$ to a function $f^\star$ defined on a larger domain $\XX^\star$, such that there's agreement on $\XX$, i.e., $f^\star(x) = f(x)$ for $x \in \XX$, and $f^\star$ maintains $f$'s relevant properties on $\XX^\star \setminus \XX$.  In the present context, the upper envelope can be viewed as a functional $1_H \mapsto \uprior\,1_H := \uprior(H)$ defined on the collection $\{\theta \mapsto 1_H(\theta): H \subseteq \TT\}$ of indicator functions/gambles.  This functional has a coherence property, by assumption, so the question is if it can be extended to a broader class of (bounded\footnote{Walley's developments focus on bounded gambles, but Part~II of \citet{lower.previsions.book} generalizes Walley's results to certain unbounded gambles.}) gambles without sacrificing coherence.  \citet[][Ch.~3]{walley1991} answers this question in the affirmative, with what he calls the {\em natural extension}---the extension that imposes the least additional structure while preserving coherence.  On the importance of natural extension, \citet[][p.~121--122]{walley1991} writes:
\begin{quote}
{\em ...natural extension may be seen as the basic constructive step in statistical reasoning; it enables us to construct new previsions from old ones.}
\end{quote}
The formula for the natural extension is rather complicated and not necessary for the present purposes.  The relevant point here is that the {\em upper envelope theorem}\footnote{As the title of their book suggests, \citet{lower.previsions.book} focus almost exclusively on {\em lower} previsions, and what they prove is a {\em lower} envelope theorem.  There is, however, an analogous result for the upper prevision and that's what I'm referring to here as the {\em upper envelope theorem}.} in \citet[][Theorem~4.38]{lower.previsions.book} links the upper expectation $\uprior\,g$ to Walley's natural extension of $\uprior$ to (bounded) gambles.  Then they follow up \citep[][Theorem~6.14]{lower.previsions.book} by linking the natural extension of $\uprior$ to the Choquet integral in \eqref{eq:choquet}.  Therefore, the ``$\mathcal{I}_\text{\sc choq}(g)$'' notation can be dropped---the Choquet integral and the upper expectation $\uprior\,g$ are the same, so the latter notation is sufficient.

Then the formula given in Equation~(8) for the upper expectation with respect to a possibility measure $\uprior$ determined by contour $q$ follows immediately---or at least {\em almost} immediately.  Using the definition $\uprior(H)$ via optimization as in Equation~(7), the Choquet integral formula \eqref{eq:choquet} above reduces to
\[ \uprior\,g = \int_0^\infty \Bigl\{ \sup_{\theta: g(\theta) \geq t} q(\theta) \Bigr\} \, dt. \]
This expression looks similar to the formula given in Equation~(8), but it's not the same; the latter roughly has the roles of $g$ and $q$ in the above expression reversed.  This apparent ``symmetry'' in the roles of $g$ and $q$, and the corresponding alternative form of the Choquet integral as I advertised in Equation~(8), is established in Proposition~7.14 (and Proposition~C.8) of \citet{lower.previsions.book} for the case of bounded $g$; this is generalized to certain unbounded gambles $g$ in, e.g., their Proposition~15.42. 

For a bit of practice with the possibility-theoretic Choquet integral, I'll first demonstrate that Choquet integration formula in Equation~(8) for a possibility measure $\uprior$ reduces to the definition of upper probability in Equation~(7), via optimization of the contour $q$, when the function $g$ is an indicator, i.e., $g(\theta) = 1(\theta \in H)$ for some $H \subseteq \TT$.  For such a case, the integrand in Equation~(8) is given by 
\[ s \mapsto \sup_{\theta: q(\theta) \geq s} 1(\theta \in H) = \begin{cases} 1 & \text{if $s \leq \sup_{\theta \in H} q(\theta)$} \\ 0 & \text{otherwise}. \end{cases} \]
Then it's clear that 
\begin{align*}
\uprior\,g & = \int_0^1 \sup_{\theta: q(\theta) > s} 1(\theta \in H) \, ds \\
& = \int_0^{\sup_{\theta \in H} q(\theta)} 1 \, ds + \int_{\sup_{\theta \in H} q(\theta)}^1 0 \, ds \\
& = \sup_{\theta \in H} q(\theta) \\
& = \uprior(H), 
\end{align*}
as was to be shown.  Warm-up complete. 

Next, I have a slightly more ambitious goal of verifying the formula below that was presented without proof in Section~4.5.2 of the main paper:
\begin{align*}
\uPi_{z^n}^{\eval} \, \ell_{\hat a(z^n)} =  \int_0^1 \Bigl\{ \sup_{\theta: \pi_{z^n}^{\eval}(\theta) \geq s} \ell_{\hat a(z^n)}(\theta) \Bigr\} \, ds = n^{-1}\bigl\{ 2 + \log(nv + 1) + (\tfrac{n}{nv+1}) \bar z_n^2 \bigr\}.
\end{align*}
In this case, the contour is $\pi_{z^n}^\eval(\theta) = 1 \wedge \eval_\theta(z^n)^{-1}$, where 
\[ \eval_\theta(z^n) = (nv + 1)^{-1/2} \exp\bigl\{ \tfrac{n}{2} (\theta - \bar z_n)^2 - \tfrac12 ( \tfrac{n}{nv + 1} ) \, \bar z_n^2 \bigr\}, \quad \theta \in \RR. \]
Then 
\begin{align*}
\pi_{z^n}^\eval(\theta) \geq s & \iff 1 \wedge \eval_\theta(z^n)^{-1} \geq s \\
& \iff \eval_\theta(z^n) \leq s^{-1} \\
& \iff (\theta - \bar z_n)^2 \leq n^{-1} \bigl\{ 2\log(s^{-1}) + \log(nv+1) + (\tfrac{n}{nv+1}) \bar z_n^2 \bigr\}.
\end{align*}
It just so happens that $\ell_{\hat a(z^n)}(\theta) = (\theta - \bar z_n)^2$, and from this it's clear that 
\[ \sup_{\theta: \pi_{z^n}^{\eval}(\theta) \geq s} \ell_{\hat a(z^n)}(\theta) = n^{-1} \bigl\{ 2\log(s^{-1}) + \log(nv+1) + (\tfrac{n}{nv+1}) \bar z_n^2 \bigr\}. \]
So it remains to integrate the right-hand side above with respect to $s$ over the interval $[0,1]$.  Using the identity 
\[ \tfrac{d}{ds} (s - s \log s) = \log(s^{-1}), \]
and the fundamental theorem of calculus, it follows that 
\begin{align*}
\uPi_{z^n}^{\eval} \, \ell_{\hat a(z^n)} & = \int_0^1 n^{-1} \bigl\{ 2\log(s^{-1}) + \log(nv+1) + (\tfrac{n}{nv+1}) \bar z_n^2 \bigr\} \, ds \\
& = n^{-1}\bigl\{ 2 + \log(nv + 1) + (\tfrac{n}{nv+1}) \bar z_n^2 \bigr\}, 
\end{align*}
as was to be shown. 


\section{More details about regularizers}
\label{A:other}

Towards a more concrete understanding of what it takes to achieve the condition in Definition~\ref{def:regularizer} in the main text, consider the case of finite $\TT$.  Then Your prior credal set $\credal$ can be interpreted as a collection of probability mass functions, and 
\[ \uprior\rho = \sup_{\prior \in \credal} \underbrace{\sum_{\theta \in \TT} \rho(\theta) \, \prior(\{\theta\})}_{\prior\rho := \E^{\Theta \sim \prior}\{ \rho(\Theta)\}} \leq \sum_{\theta \in \TT} \rho(\theta) \, \sup_{\prior \in \credal} \prior(\{\theta\}). \]
Since the functional $\prior \mapsto \prior\rho$ is linear and the domain $\credal$ is closed and convex, it's well-known (related to the Krein--Milman theorem)
%, e.g., \citealt{bauer1958}) 
that the supremum is attained on the extreme points of $\credal$.  Then it's clear that, for each (sub-)probability mass function\footnote{A sub-probability mass function $\eta \geq 0$ satisfies $\sum_{\theta \in \TT} \eta(\theta) \leq 1$. In the present case, there's no reason not to take $\eta$ as a genuine probability mass function whose sum is exactly equal to 1.} 
$\eta$ on $\TT$, the following function $\rho_\eta$ is a regularizer:
\[ \rho_\eta(\theta) = \frac{\eta(\theta)}{\sup_{\prior \in \credal} \prior(\{\theta\})}, \quad \theta \in \TT. \]
Note that the form $\rho_\eta$ above mimics the universal inference formulation of an e-process in \citep[e.g.,][]{wasserman.universal}, albeit in a different context.  Two special cases deserve mention.  First, if $\credal$ is a singleton, then all the admissible regularizers would be of the form $\rho_\eta$ above, for some probability mass function $\eta$.  Second, if $\credal$ is vacuous, i.e., if it contains all the probability mass functions on $\TT$, then the denominator is constant equal to 1 and, since $\eta$ can't exceed 1, neither can $\rho_\eta$, hence, it's trivial.  So, clearly, it's a bad idea for You to plug a ``vacuous prior'' into the regularizer construction.  If Your prior information is genuinely vacuous, then just use $\rho \equiv 1$, so that Your regularized e-process matches the original, unregularized e-process.  

In the main paper I focused exclusively on regularizers for the case where Your prior information is encoded as a possibility measure with contour $q$.  I did so for two reasons: (a)~I think the possibilistic formulation makes sense, and (b)~it's relatively concrete to handle compared to other formulations.  But this isn't the only option, so here I want to briefly mention a couple other strategies.  This is absolutely not intended to be an exhaustive list of alternatives, nor is it my intention for this to be a tutorial on how to construct a regularizer in a non-possibilistic setup.  My goal is simply to expose the reader to some other avenues to pursue if the possibilistic version I presented in the main paper isn't fully satisfactory; this also demonstrates my point that the possibilistic formulation is simpler and more concrete. 

Arguably one of the most common imprecise probabilities models are the so-called {\em contamination classes}, {\em gross error models}, or {\em linear--vacuous mixtures}, often found in the literature on robust statistics \citep[e.g.,][]{huber1973.capacity, huber1981, walley1991, walley2002, wasserman1990}.  This corresponds to a choice of centering probability $\prior_\text{cen}$ and a weight $\eps \in (0,1)$.  Then the credal set is given by 
\[ \credal = \bigl\{ (1-\eps) \, \prior_\text{cen} + \eps \, \prior: \, \text{$\prior$ is any probability on $\TT$} \bigr\}. \]
The basic idea is that You think the precise probability $\prior_\text{cen}$ is a pretty good assessment of Your uncertainty about $\Theta$, but You don't fully trust the information that lead to this assessment; so $\eps$ is like the ``chance You were misled,'' and if You were mislead, then literally anything might be true.  It's relatively clear that the corresponding upper probability/prevision, say $\uprior$, has upper expectation of $\rho(\Theta)$ given by 
\[ \uprior\,\rho = (1-\eps) \, \prior_\text{cen}\,\rho + \eps \times \sup_{\theta \in \TT} \rho(\theta). \]
Then the goal is to choose $\rho$ such that the right-hand side above is (less than or) equal to 1.  There's no simple formula for this, like in the case of possibilistic prior information, but the idea is take $\rho$ such that it's relatively small where You ``expect'' $\Theta$ to be, i.e., in the support of $\prior_\text{cen}$, and then somewhat large elsewhere.  If, for example, $\prior_\text{cen}$ has bounded support, so that $\rho$ can be defined independently on that support and elsewhere, then $\rho$ can take values as large as $\eps^{-1}(1-\eps) m$ outside that support, where $m = \prior_\text{cen}\,\rho$. 

A second kind of imprecise probability model is what \citet[][Sec.~2.9.4]{walley1991} calls the {\em constant odds ratio} model.  Similar to that above, this is indexed by a precise probability distribution, which I'll denote again as $\prior_\text{cen}$, and a weight $\tau \in (0,1)$.  Walley explains this model in the context of a risky investment, where $\tau$ represents the rate at which You're taxed on said investment.  Setup and context aside, the lower and upper probability of a hypothesis/event $H$ under this model is 
\[ \lprior(H) = \frac{(1-\tau) \prior_\text{cen}(H)}{1 - \tau \prior_\text{cen}(H)} \quad \text{and} \quad \uprior(H) = \frac{\prior_\text{cen}(H)}{1-\tau \prior_\text{cen}(H)}. \]
The name ``constant odds ratio model'' comes from the fact that the lower and upper odds of $H$ versus $H^c$ are 
\[ \frac{\lprior(H)}{\uprior(H^c)} = \frac{(1-\tau) \prior_\text{cen}(H)}{\prior_\text{cen}(H^c)} \quad \text{and} \quad \frac{\uprior(H)}{\lprior(H^c)} = \frac{\prior_\text{cen}(H)}{(1-\tau)\prior_\text{cen}(H^c)}, \]
respectively, so the rato of lower to upper odds is constant in $H$:
\[ \frac{\lprior(H) / \uprior(H^c)}{\uprior(H) / \lprior(H^c)} = \cdots = (1-\tau)^2. \]
The upper expectation of $\rho(\Theta)$ under this model is not so straightforward as for the contamination model above, but Walley shows that $\uprior\,\rho$ solves the equation $f(x)=0$, where $f(x) = \tau \, \prior_\text{cen}(\rho - x)^+  + (1-\tau) \, ( \prior_\text{cen}\,\rho - x )$, with $w^+ = \max(w, 0)$ the positive part of $w \in \RR$.  For a given $\rho$, one can numerically solve for $x$ as a function of $(\tau, \prior_\text{cen}\,\rho)$, to obtain $\uprior\,\rho$.  But choosing $\rho$ such that this numerical solution is $\leq 1$ requires care. 

Other kinds of imprecise models can be considered, including monotone capacities \citep[e.g.,][]{huber1973.capacity, wasserman.kadane.1990, sundberg.wagner.1992}, belief functions \citep[e.g.,][]{denoeux1999, shafer1976, dempster1967}, and probability-boxes \citep[e.g.,][]{destercke.etal.pbox, ferson.etal.pbox}, and formulas for their upper expectations can be found in, e.g., Chapters~6--7 of \citet{lower.previsions.book}.  

%{\color{red} Talk about the generalization to other kinds of prior assessments... Basically, once we have an upper probability $\uprior$ for $\Theta$ on $\TT$, not necessarily a possibility measure, we just need to define a regularizer gamble $\rho$ such that $\uprior \, \rho \leq 1$... Then all the same theory goes through... The difficulty is that the construction of $\rho$ isn't so concrete in these other contexts... but see the coherent upper previsions in Section~2.9 of Walley 1991 (p.~92--)...}



%\section{Credal set pull-backs}
%\label{A:pullback}




\section{Justification of the product form in Eq.~\eqref{eq:reg.eprocess}}
\label{SS:product}

\subsection{Bayesian-like updating coherence}

Here I'll present the previously-advertised justification for my choice to define the regularized e-process in Equation~(10) as a product of the regularizer and the original e-process.  Actually, I'll present two such justifications, the first of which is based on an analogy to the updating coherence property familiar in Bayesian inference, i.e., ``today's posterior is tomorrow's prior.''  To see this, let $z^n \equiv z^{1:n}$ be the data, processed as $\eval_\theta(z^{1:n})$, and $\rho(\theta)$ the regularizer.  Combining these according to Equation~\eqref{eq:reg.eprocess} gives the regularized e-process $\eval^\text{reg}(z^n, \theta)$.  Now, suppose that more data $z^{n:(n+m)}$ becomes available.  Since e-processes are typically combined via multiplication, I can write 
\[ \eval^\text{reg}(z^{1:(n+m)}, \theta) = \eval_\theta(z^{1:(n+m)}) \times \rho(\theta) = \eval_\theta(z^{n:(n+m)}) \times \underbrace{\eval_\theta(z^{1:n}) \times \rho(\theta)}_{\eval^\text{reg}(z^{1:n}, \theta)}. \]
That is, the regularized e-process based on $z^{1:n}$ becomes the updated, old-data-dependent version of the regularizer that's combined with a new-data-dependent e-process according to the rule in  Equation~\eqref{eq:reg.eprocess} of the main paper. 
%\eqref{eq:reg.eprocess}. 


\subsection{Formal dominance}

The second justification is based on a more formal demonstration that the product form dominates other strategies.  Frankly, this justification isn't much more compelling than the fact that multiplying the two ingredients is clearly the most natural way to merge them.  There are, however, some other reasonable options, e.g., averaging, so the result below adds some valuable insight.  
%It's just that the separation between the product rule in \eqref{eq:reg.eprocess} and other reasonable merging strategies provided Proposition~\ref{prop:dominates} below isn't substantial.  
The reader may have ideas on how to strengthen this result in one way or another.  

Throughout this section, to simplify notation, etc., I'll assume that $f$ is the identity function, so that the quantity of interest $\Theta$ corresponds exactly with $\Omega$.  This helps because it allows me to drop $f$ (and $\Omega$) from the notation, to drop the set of pullback measures in the theoretical formulation, and to express the model as ``$\prob_\theta$,'' directly in terms of $\theta$.  I'm also going to drop the explicit mention of a generic stopping time, and write ``$Z$'' for the observable data.  Since the main result of this section is conceptual in nature, these simplifications don't affect the take-away message.  

Following \citet{vovk.wang.2021}, define a function $(r,e) \mapsto m(r,e)$ to be a {\em re-merging function}---pronounced ``R-E-merging'' because it merges a {\bf r}egularizer and an {\bf e}-process---if, for  given prior information $\uprior$ about $\Theta=f(\Omega)$ and a corresponding regularizer $\rho$, the merged variable $m(\rho, \eval)$ satisfies two key properties:
\begin{itemize}
\item It treats the data as sovereign in the sense that 
\begin{equation}
\label{eq:sovereign}
e \geq 1 \implies m(r, e) \geq r, \quad \text{for all $r \in (0,\infty)$}, 
\end{equation}
which, in words, means that if the data-dependent component, $\eval_\theta(\cdot)$, offers evidence that doesn't favor hypothesis ``$\Theta=\theta$,'' then the merged e-process will show less support for that hypothesis than the prior information alone did.
\vspace{-2mm}
\item It's a regularized e-process in the sense that the property advertised in Equation~\eqref{eq:reg.ville0} holds for any input $\eval$, i.e., 
\begin{equation}
\label{eq:merging}
\uuprob\bigl[ m\{ \rho(\Theta), \eval_{\Theta}(Z) \} \bigr] \leq 1 \quad \text{for all e-processes $\eval$}. 
\end{equation}
\end{itemize} 
Recall that $\uuprob$ is determined by the model and by the prior information, so $\uuprob$ is fixed by the context of the problem---all that's free to vary in these considerations is the input e-process $\eval$ and the merger function $m$. The class of re-merging functions is non-empty, since the product mapping is an re-merger.  My claim is that, in a sense to be described below, the product merger in Equation~\eqref{eq:reg.eprocess} is ``best'' among all the re-merging functions.

Continuing to follow \citet{vovk.wang.2021}, I'll say that an re-merging function $m$ weakly dominates another re-merging function $m'$ if 
\begin{equation}
\label{eq:dominance}
(r,e) \in [1, \infty) \times [1, \infty) \implies m(r,e) \geq m'(r,e). 
\end{equation}
The idea is that, in the case where neither the data nor the prior show signs of compatibility with a given value $\theta$, then merging based on $m$ is more aggressive, i.e., shows no less evidence against $\theta$, than merging based on $m'$.  The complement to an e-process's anytime validity property is its efficiency, and efficiency requires that the e-process take large values when evidence is incompatible with a hypothesis in question; so, an re-merging function that returns larger regularized e-process values is preferred.  

The present case differs in many ways from that in \citet{vovk.wang.2021}, mainly due to the presence of the prior information, so I'll need some additional control.  Towards this, I'll say that an re-merging function $m$ {\em $\uuprob$-strictly weakly dominates} another re-merging function $m'$ if $m$ weakly dominates $m'$ in the sense of \eqref{eq:dominance} above, and if
\begin{align}
\uuprob \bigl[ m\{ \rho(\Theta), \eval_\Theta(Z) \} > m'\{ \rho(\Theta), \eval_\Theta(Z) \}, \; & \notag \\
\rho(\Theta) \geq 1, \, \eval_\Theta(Z) \geq 1 & \bigr] > 0, \quad \text{for all e-processes $\eval$}. \label{eq:uuprob.dominance}
\end{align}
Vovk and Wang's ``weak dominance'' in \eqref{eq:dominance} allows $m$ and $m'$ to be the same, so roughly all that \eqref{eq:uuprob.dominance} adds is that there exists a joint distribution for $(Z,\Theta)$---corresponding to a prior $\prior$ in Your credal set $\credal$---with respect to which the ``strict inequality'' event on the right-hand side of \eqref{eq:dominance} has positive probability.  Therefore, $\uuprob$-strict weak dominance rules out the possibility that $m$ and $m'$ differ only in an insignificant way relative to $\uuprob$.  In the special case where the prior information is vacuous, like in Vovk and Wang, strict weak dominance holds if $m$ weakly dominates $m'$ and if, for each e-process $\eval$, there exists a $\theta=\theta(\eval)$ such that $\rho(\theta) \geq 1$ and
\[ \prob_\theta\bigl[ m\{ \rho(\theta), \eval_\theta(Z) \} > m'\{ \rho(\theta), \eval_\theta(Z) \}, \, \eval_\theta(Z) \geq 1 \bigr] > 0, \]
i.e., if roughly strict inequality holds with positive model probability for some $\theta$.  

In the case of vacuous prior information, I'll require $\rho \equiv 1$, as was suggested in the main paper.  When the prior information is non-vacuous, the result below is restricted to {\em non-trivial} regularizers, which was loosely defined earlier as a regularizer that's not upper bounded by 1.  Here, however, I need to be more specific about what non-trivial means: specifically, $\rho$ is non-trivial (relative to the partial prior information) if $\uprior\{ \rho(\Theta) > 1 \} > 0$.  In words, non-triviality means that there exists a probability $\prior$ such that $\rho(\Theta)$ isn't $\prior$-almost surely upper bounded by 1. Finally, I'll also say that a regularizer is {\em admissible} if $\uprior\rho=1$, that is, if it can't be made larger in any substantive way without violating the upper bound presented in Definition~\ref{def:regularizer}. 

The following is similar to Proposition~4.2 in \citet{vovk.wang.2021} on what they refer to as {\em ie-merging} functions for merging independent e-values.  I show that the product mapping isn't strictly weakly dominated by any other re-merging functions.  

\begin{prop}
\label{prop:dominates}
Given $\uprior$, fix a non-trivial, admissible regularizer $\rho$.  Then the product rule in Equation~\eqref{eq:reg.eprocess} isn't $\uuprob$-strictly weakly dominated by any other re-merging function. 
\end{prop}

\begin{proof}
The proof is by contradiction; that is, I'll assume that the product rule defined in Equation~\eqref{eq:reg.eprocess} is $\uuprob$-strictly weakly dominated in the sense above and show that this leads to a contradiction.  Let $m$ denote this assumed-to-exist dominant re-merging function. 

The assumed non-triviality of the regularizer $\rho$ implies existence of points $\theta$ such that $\rho(\theta) > 1$ and such that $\rho(\theta) \leq 1$, and neither of these sets have $\uprior$-probability 0.  Some of the $\theta$'s in the first set could have $\rho(\theta) = \infty$, but, those must have  $\uprior$-probability 0 for, otherwise, the property that $\uprior\,\rho \leq 1$ would be violated; recall that the assumed admissibility of $\rho$ means that $\uprior\,\rho=1$.  

From the assumed strict weak dominance of $m$, and from \eqref{eq:sovereign}, I can deduce the following bound for generic inputs $(r,e)$:
\begin{align*}
m(r,e) & = m(r,e) \, \bigl( 1_{r < 1, e < 1} + 1_{r \geq 1, e < 1} + 1_{r < 1, e \geq 1} + 1_{r \geq 1, e \geq 1} \bigr) \\
& \geq  m(r,e) \, 1_{r < 1, e < 1} + m(r,e) \, 1_{r \geq 1, e < 1} + r \, 1_{r < 1, e \geq 1} + re \, 1_{r \geq 1, e \geq 1},
\end{align*}
where the ``$r$'' factor in the third term is by \eqref{eq:sovereign} and the ``$re$'' factor in the fourth term is by the assumed weak dominance.  The first two terms depend on details of the particular choice of $m$ on the respective ranges of $(r,e)$, details that can't be controlled with only the information provided.  I can apply the trivial non-negativity bound, however, and, from the above display, conclude that 
\[ m(r,e) \geq r \, 1_{r < 1, e \geq 1} + re \, 1_{r \geq 1, e \geq 1}. \]
(In fact, with the input e-process to be constructed next, those two terms I lower-bounded by 0 would typically be equal to 0, so the above inequality isn't loose.)  The not-strict equality in the above display is pointwise in $(r,e)$, i.e., I can't rule out equality above for any given pair $(r,e)$.  But a pointwise lower bound isn't the goal---I'm aiming for a lower bound in upper expectation.  For this latter goal, I'll apply \eqref{eq:uuprob.dominance} to flip the not-strict inequality ``$\geq$'' in the above display a strict inequality; more on this below.  

Towards establishing a contradiction, I only need to produce one example of an input e-process such that the conclusion is problematic, and I'll do this with an incredibly simple e-process.  Specifically, I define the input e-process $\eval^\star$ as follows:
\begin{itemize}
\item if $\theta$ is such that $\rho(\theta) < 1$, then $\eval_\theta^\star(Z) \equiv 1$, and 
\vspace{-2mm}
\item if $\theta$ is such that $\rho(\theta) \geq 1$, then  
\[ \eval_\theta^\star(Z) = \begin{cases} 2 & \text{if $Z \in \mathcal{E}_\theta$} \\ 0 & \text{otherwise}, \end{cases} \]
where $\mathcal{E}_\theta$ is an event with $\prob_\theta(\mathcal{E}_\theta) = \frac12$.
\end{itemize} 
It's easy to check that $\E_\theta\{ \eval_\theta^\star(Z) \} = 1$ for all $\theta$, so $\eval^\star$ is a genuine e-process.  What's important about this particular e-process, as it pertains to the present proof, is that 
\[ 1_{\rho(\theta) < 1, \eval_\theta^\star \geq 1} = 1_{\rho(\theta) < 1} \quad \text{and} \quad \eval_\theta^\star \, 1_{\rho(\theta) \geq 1, \eval_\theta^\star \geq 1} = \eval_\theta^\star \, 1_{\rho(\theta) \geq 1}. \]
Therefore, from the pointwise analysis above, 
\begin{align*}
m\{ \rho(\Theta), \eval_\Theta^\star(Z) \} & \overset{\text{\tiny (+)}}{\geq} \rho(\Theta) \, 1_{\rho(\Theta) < 1, \eval_\Theta^\star(Z) \geq 1} + \eval_\Theta^\star(Z) \, \rho(\Theta) \, 1_{\rho(\Theta) \geq 1, \eval_\Theta^\star(Z) \geq 1} \\
& = \rho(\Theta) \, 1_{\rho(\Theta) < 1} + \eval_\Theta^\star(Z) \, \rho(\Theta) \, 1_{\rho(\Theta) \geq 1}. 
\end{align*}
The (+) symbol above is to remind the reader that, while ``$\geq$'' holds pointwise, there's actually more that can be said.  That is, in addition to ``$\geq$'' for every $(z,\theta)$ pair, there exists a joint distribution for $(Z,\Theta)$, compatible with the available prior information, such that strict inequality holds with positive probability.  This implies that the inequality ``$\geq$'' highlighted with (+) is {\em strict inequality} ``$>$'' in expectation with respect to the aforementioned joint distribution.  Then 
\begin{align}
\uuprob\bigl[ m\{ \rho(\Theta), \eval_\Theta^\star(Z) \} \bigr] & > \uuprob\bigl\{ \rho(\Theta) \, 1_{\rho(\Theta) < 1} + \eval_\Theta^\star(Z) \, \rho(\Theta) \, 1_{\rho(\Theta) \geq 1} \bigr\} \notag \\
& = \sup_{\prior \in \credal} \E^{(Z,\Theta) \sim \prob_\bullet \otimes \prior} \bigl\{ \rho(\Theta) \, 1_{\rho(\Theta) < 1} + \eval_\Theta^\star(Z) \, \rho(\Theta) \, 1_{\rho(\Theta) \geq 1} \bigr\} \notag \\
& = \uprior\bigl\{ \rho(\Theta) \, 1_{\rho(\Theta) < 1} + \rho(\Theta) \, 1_{\rho(\Theta) \geq 1} \bigr\} \label{eq:dominance.a} \\
& = \uprior\,\rho \notag \\
& = 1, \label{eq:dominance.b}
\end{align}
where \eqref{eq:dominance.a} holds because $\eval^\star$ is an e-process relative to the model, i.e., $\E_\theta(\eval_\theta^\star) = 1$ for all $\theta$, and \eqref{eq:dominance.b} holds by the assumed admissibility of the regularizer $\rho$.  Therefore, the defining property \eqref{eq:merging} of an re-merging fails for the chosen $m$; that is, I constructed an e-process $(\eval^\star)$ such that $\uuprob[ m\{ \rho(\Theta), \eval_\Theta^\star(Z) \} ] > 1$.  Since $m$ was assumed to be re-merging, this creates the desired contradiction.  So, I conclude that, as claimed, there's no re-merging function $m$ that $\uuprob$-strictly weakly dominates the product rule in Equation~\eqref{eq:reg.eprocess}.  
\end{proof}



\section{More efficiency gains}
\label{S:emp.efficiency}

To follow up on the illustrations given in the main text, here I'll consider two different types of prior information and associated prior possibility contours.  As before, these both will be indexed by a parameter $K$, although the meaning of $K$ will be different in each case.  Consequently, the corresponding regularized e-processes will not be comparable between the two types of prior information here, but they will be comparable as $K$ varies within the types.  Here, $K \in \{0.1, 0.2, 0.4, 0.8\}$ and the two types of prior information I'll consider are below; first an explanation in words and then a mathematical description. 
\begin{enumerate}
\item ``You expect that $|\Theta| \leq K$.'' This is pretty clear in words but, mathematically, this corresponds to a credal set $\credal$ that contains exactly those distributions $\prior$ such that $\E^{\Theta \sim \prior}|\Theta| \leq K$.  Of course, this includes certain Gaussian, uniform, and even some heavier-tailed priors. As shown in \citet{dubois.etal.2004} and elsewhere, this prior information can be described mathematically via the possibility contour 
\[ q(\theta) = 1 \wedge K|\theta|^{-1}, \quad \theta \in \TT. \]
\item ``You're at most $100K/5\%$ sure that $|\Theta| > 2K$.'' This means that You can't rule out the possibility of $|\Theta| \leq 2K$, but that You judge the probability of the event $|\Theta| > 2K$ to be upper-bounded by $K/5$.  Mathematically, this can be described easily by the possibility contour 
\[ q(\theta) = \tfrac{K}{5} \, 1(|\theta| > 2K) + 1( |\theta| \leq 2K ), \quad \theta \in \TT. \]
I'm using the $K$-dependent weights only so that the different lines in the plots shown in Figure~\ref{fig:prob} don't overlap on a large part of the range of $\theta$ values. 
%\item ``For each confidence level $c \in [0,1]$, You're $100(1-c)\%$ sure that $\Theta$ is in the central $1-c$ interval corresponding to a $\nm(0, K)$ distribution.''  Note that this is different from, and considerably weaker than, a belief that $\Theta$ is $\nm(0,K)$ distributed.  In fact, this just means that $\nm(0,K)$ is a (most diffuse) element of the prior credal set $\credal$.  This is described mathematically by the probability-to-possibility transform of the dominating Gaussian distribution:
%\begin{align*}
%q(\theta) = \prob^{\Theta \sim \nm(0,K)}\{ {\tt dnorm}_K(\Theta) \leq {\tt dnorm}_K(\theta) \} = 1 - {\tt pchisq}(\theta^2 / K), 
%\end{align*}
%where ${\tt dnorm}_K$ is the $\nm(0,K)$ density and ${\tt pchisq}$ is the $\chisq(1)$ CDF.  
\end{enumerate} 

These two kinds of prior information are rather weak, only based on some very basic judgments about a first moment and the probability of a single event.  Figures~\ref{fig:mean} and \ref{fig:prob} plot the regularized and unregularized log-transformed e-process as a function of $\theta$ for the three types of priors, respectively, and for three different values of the observed sample mean $\bar z$ as in the main text, based on a sample of size $n=5$.
%\begin{itemize}
%\item $\bar z=0.25$ is pretty consistent with all the priors;
%\vspace{-2mm}
%\item $\bar z=0.5$ is at least marginally inconsistent with the priors, and;
%\vspace{-2mm}
%\item $\bar z=1$ is rather inconsistent with the priors.
%\end{itemize}
The black line corresponds to the unregularized e-process, and the four colored lines correspond to the different values of $K$; the dashed horizontal line corresponds to $-\log0.05 \approx 3$, which is the cutoff that determines the (regularized) e-process's 95\% confidence interval.  Not surprisingly, given that prior Types~1--2 are rather weak, the effect of regularization is difficult to detect.  There is a small efficiency gain in Figures~\ref{fig:mean}--\ref{fig:prob}, more so in the latter, since the prior does more than strictly constrain $\Theta$ to an interval around the origin.  

\begin{figure}[t]
\begin{center}
\subfigure[$\bar z=0.25$]{\scalebox{0.41}{\includegraphics{nm_mean_x025}}}
\subfigure[$\bar z=0.5$]{\scalebox{0.41}{\includegraphics{nm_mean_x05}}}
\subfigure[$\bar z=1$]{\scalebox{0.41}{\includegraphics{nm_mean_x1}}}
\end{center}
\caption{Plot of $\theta \mapsto \eval^\text{reg}(z^n, \theta)$ for three different data sets $z^n$ based on prior Type~1.}
\label{fig:mean}
\end{figure}

\begin{figure}[t]
\begin{center}
\subfigure[$\bar z=0.25$]{\scalebox{0.41}{\includegraphics{nm_prob_x025}}}
\subfigure[$\bar z=0.5$]{\scalebox{0.41}{\includegraphics{nm_prob_x05}}}
\subfigure[$\bar z=1$]{\scalebox{0.41}{\includegraphics{nm_prob_x1}}}
\end{center}
\caption{Plot of $\theta \mapsto \eval^\text{reg}(z^n, \theta)$ for three different data sets $z^n$ based on prior Type~2.}
\label{fig:prob}
\end{figure}




\section{Anytime validity implies no-sure-loss}
\label{S:euq.coherence}

When data $Z^N=z^n$ is fixed, the e-possibilistic IM typically defines an  imprecise probability model where $\uPi_{z^n}^{\eval \times \rho}$ has the mathematical properties of a possibility measure.  Provided that the function $\theta \mapsto \pi_{z^n}^{\eval \times \rho}(\theta)$ isn't bounded away from 1, which is typically the case, the function $H \mapsto \uPi_{z^n}^{\eval \times \rho}(H)$ determined by maximizing a possibility contour function over the set $H \subseteq \TT$ as in \eqref{eq:upper.p.poss} is a genuine possibility measure.  This implies that its credal set is non-empty and, in turn, that the IM output is {\em coherent} in the sense of \citet[][Sec.~2.5]{walley1991}.  For the present purposes, it's enough to understand coherence as a stronger version of {\em no-sure-loss}, which can be described as follows. Suppose beliefs concerning the uncertain value $\Theta$ are assessed via the buying/selling prices they consider acceptable for certain gambles about $\Theta$.  Let $z^n$ be one of the aforementioned typical data sets and let Your buying/selling prices be determined by the IM output $(\lPi_{z^n}^{\eval \times \rho}, \uPi_{z^n}^{\eval \times \rho})$ according to the interpretation in \eqref{eq:prices}.  Then coherence of Your IM implies that
\[ \sup_{\theta \in \TT} \sum_{k=1}^K \{ \uPi_{z^n}^{\eval \times \rho}(H_k) - 1(\theta \in H_k)\} \geq 0 \quad \text{for all $K$ and $(H_1,\ldots,H_K)$ combos}. \] 
To get the intuition, suppose the above condition fails.  Then there exists a combination $K$ and $(H_1,\ldots,H_K)$, and a sufficiently small $\delta > 0$, such that 
\[ \sup_{\theta \in \TT} \sum_{k=1}^K \bigl[ \{ \uPi_{z^n}^{\eval \times \rho}(H_k) + \delta\} - 1(\theta \in H_k) \bigr] < 0. \]
Since $\uPi_{z^n}^{\eval \times \rho}(H_k)$ is, by definition, Your infimum selling price for the gamble $1(\Theta \in H_k)$, the transactions where You accept payment of $\uPi_{z^n}^{\eval \times \rho}(H_k) + \delta$ dollars for \$$1(\Theta \in H_k)$ for each $k$ are all acceptable to You.  But then the above display reveals a troubling result: somehow, by only making transactions that are acceptable {\em a priori}, You end up with negative total earnings regardless of what value the uncertain $\Theta$ takes on.  This indicates a severe shortcoming in Your pricing scheme; fortunately, the e-possibilistic IM is typically free of this internal inconsistency.  

I said ``typically'' several times in the above paragraph, and the explanation here is exactly the same as in Remark~\ref{re:typically} above.  The IM output $\uPi_{z^n}^{\eval \times \rho}$ would fail to be a possibility measure if and only if $\pi_{z^n}^{\eval \times \rho}$ was bounded away from 1 on $\TT$; I've been calling that function a ``possibility contour'' but that's only legitimate if $\sup_{\theta \in \TT} \pi_{z^n}^{\eval \times \rho}(\theta) = 1$.  For $\pi_{z^n}^{\eval \times \rho}(\theta)$ to be strictly less than 1, and hence the IM output determines an incoherent imprecise probability, would require that $\theta \mapsto \eval^\text{reg}(z^n, \theta)$ also be bounded strictly greater than 1.  But the regularized Ville's inequality implies that $\uuprob(\eval^\text{reg}) \leq 1$, i.e., $\eval^\text{reg}(Z^N,\Theta)$ ``tends'' to be less than 1, so a data set $z^n$ could indeed be called atypical if it were such that $\eval^\text{reg}(z^n, \theta)$ were strictly greater than 1 for all $\theta$.  

In addition the fixed-data behavioral considerations, it's natural to interpret $(\credal, z^n) \mapsto (\lPi_{z^n}^{\eval \times \rho}, \uPi_{z^n}^{\eval \times \rho})$ as a rule by which ``prior'' information is {\em updated} in light of data ($z^n$) to a ``posterior'' quantification of uncertainty.  More familiar notions of imprecise-probabilistic updating include generalized Bayes rule \citep{walley1991, miranda.cooman.chapter} and Dempster's rule \citep[e.g.,][]{shafer1976, cuzzolin.book}.
%; this and other possibilistic IM constructions share similarities and differences with Bayes's and Dempster's rules, so I think it's fair to interpret them as updating rules.  
With this ``updating rule'' interpretation comes further questions about the IM's ability to protect You from sure loss, etc.  What's different here is that there's a temporal component: can I force You into transactions such that, no matter what data is observed, You lose money?  If so, then there's a serious issue with Your assessments. Mathematically, this {\em sure loss} property---see, e.g., \citet[][Sec.~2.4.1]{walley1991} and \citet[][Def.~3.3]{gong.meng.update}---corresponds to existence of a hypothesis $H \subset \TT$ such that 
\begin{equation}
\label{eq:sure.loss}
\sup_{z^n} \uPi_{z^n}^{\eval \times \rho}(H) < \lprior(H) \quad \text{or} \quad  \inf_{z^n} \lPi_{z^n}^{\eval \times \rho}(H) > \uprior(H). 
\end{equation}
For intuition, consider the first of the above two inequalities.  If this inequality holds, then, for any pair of positive numbers $(\eps, \delta)$, You'd be willing to buy the gamble \$$1(\Theta \in H)$ from me for $\$\{\lprior(H) - \eps\}$ and then sell the same gamble back to me, after observing $z^n$, for \$$\{\sup_{z^n} \uPi_{z^n}^{\eval \times \rho}(H) + \delta\}$.  No matter whether $\Theta \in H$ or $\Theta \not\in H$, the payoff You receive from this sequence of transactions is 
\[ \Bigl\{ \sup_{z^n} \uPi_{z^n}^{\eval \times \rho}(H) + \delta \Bigr\} - \{\lprior(H) - \eps\} = \Bigl\{ \underbrace{\sup_{z^n} \uPi_{z^n}^{\eval \times \rho}(H) - \lprior(H)}_{\text{$<0$, by \eqref{eq:sure.loss}}} \Bigr\} + (\eps + \delta). \]
Since both individual transactions are acceptable to You, and there exists pairs $(\eps, \delta)$ such that You net payoff is strictly negative, it follows that You can be made a sure loser.  

Fortunately, the possibilistic IM provably avoids an even less severe internal inconsistency, which I call {\em one-sided contraction}:
\begin{equation}
\label{eq:contraction}
\sup_{z^n} \uPi_{z^n}^{\eval \times \rho}(H) < \uprior(H) \quad \text{or} \quad \inf_{z^n} \lPi_{z^n}^{\eval \times \rho}(H) > \lprior(H). 
\end{equation}
One-sided contraction is less concerning than sure-loss, but still problematic.  To see this, suppose that the first of the two inequalities in \eqref{eq:contraction} holds for a given $H$.  If I want to buy \$$1(\Theta \in H)$ from You {\em a priori}, then it's apparent that I can wait until data $z^n$ is revealed and purchase the gamble for a lower price.  This doesn't imply that You lose money, only that You're systematically giving up opportunities to earn more money.  This can't happen in the familiar, precise Bayesian case: since the expected value of a posterior probability is the prior probability, it's impossible for each posterior/conditional probability to be less than the prior/marginal probability.  

The stronger notion of (two-sided) {\em contraction} corresponds to replacing the ``or'' in the above display with ``and,'' and the dual notion of {\em dilation} corresponds to flipping both inequalities and replacing ``or'' with ``and.'' Both contraction and dilation are problematic in their own respects, but contraction is generally more serious.  Of course, if an e-possibilistic IM avoids one-sided contraction, as the next theorem establishes, then it necessarily avoids both sure-loss and (two-sided) contraction.  

\begin{thm}
\label{thm:no.sure.loss}
The e-possibilistic IM avoids one-sided contraction, i.e., there are no $H$ such that \eqref{eq:contraction} holds for the updating rule $(\credal, z^n) \mapsto (\lPi_{z^n}^{\eval \times \rho}, \uPi_{z^n}^{\eval \times \rho})$. 
\end{thm}

\begin{proof}
Fix any hypothesis $H \subseteq \TT$.  I'll focus on proving that the first inequality in \eqref{eq:contraction} doesn't hold, i.e., that 
\begin{equation}
\label{eq:no.contraction}
\sup_{n,z^n} \uPi_{z^n}^{\eval \times \rho}(H) \geq \uprior(H). 
\end{equation}
The version involving lower probabilities is proved similarly.  To start, note that 
\[ \inf_{n, z^n} \eval_\theta(z^n) \leq \sup_N \sup_{\omega: f(\omega)=\theta} \E_\omega \{ \eval_\theta(Z^N) \} \leq 1, \]
where the first inequality is because an average is never smaller than the minimum, and the second inequality by Ville.  That is, for any $\theta$, there exists $(n,z^n)$ such that $\eval_\theta(z^n)$ is no more than 1.  Now write out the left-hand of \eqref{eq:no.contraction} as follows:
\[ \sup_{n,z^n} \uPi_{z^n}^{\eval \times \rho}(H) = \sup_{n,z^n} \sup_{\theta \in H} \pi_{z^n}^{\eval \times \rho}(\theta) = 1 \wedge \Bigl\{ \inf_{n,z^n} \inf_{\theta \in H} \rho(\theta) \, \eval_\theta(z^n) \Bigr\}^{-1}. \]
It follows from the ``first observation'' above that 
\[ \sup_{n,z^n} \uPi_{z^n}^{\eval \times \rho}(H) \geq 1 \wedge \Bigl\{ \inf_{\theta \in H} \rho(\theta) \Bigr\}^{-1}. \]
Since $\rho(\theta) \geq \{ \inf_{\theta \in H} \rho(\theta) \} \times 1(\theta \in H)$, it's easy to see that 
\[ \Bigl\{ \inf_{\theta \in H} \rho(\theta) \Bigr\} \times \uprior(H) \leq \uprior\,\rho \leq 1, \]
where the first inequality is by monotonicity of the upper expectation and the second by definition of the regularizer $\rho$.  Plugging this bound into that above gives 
\[ \sup_{n,z^n} \uPi_{z^n}^{\eval \times \rho}(H) \geq 1 \wedge \uprior(H) = \uprior(H), \]
which completes the proof of the theorem. 
\end{proof}

An even stronger internal consistency property, a similarly-temporal version of {\em coherence}, might be desired.  Suffice it to say that the e-possibilistic IM generally satisfies only one of the two necessary and sufficient conditions \citet[][Sec.~6.5.2]{walley1991} for coherence.  There are cases where coherence can be achieved, but it can't be achieved in general; see \citet[][Sec.~3.3]{martin.partial} for more discussion on this point. 
%\citet[][Sec.~6.5.2]{walley1991} presents two conditions that are necessary and sufficient for the kind of coherence that I have in mind here.  The first of those conditions is precisely no-one-sided-contraction, so it'd not be out of line to say that the e-possibilistic IM is {\em half-coherent} based on the result in Theorem~\ref{thm:no.sure.loss}.  The second of Walley's conditions is more complicated and not worth writing down here.  Since there are cases in which the only coherent prior-to-posterior updating rule is generalized Bayes \citep[e.g.,][Sec.~6.5.4]{walley1991}, it's not possible for the proposed IM to be coherent in general.  There are cases, however, where the IM is coherent: in the fixed sample size case, I argued in \citet[][Sec.~3.3]{martin.partial} that coherence holds provided that ``$z \mapsto \uPi_z(H)$'' is continuous for each $H$.  The fact that the sample size here isn't fixed creates some technical complications, so I'll postpone this investigation into coherence for now.



%Consider a situation in which the goal is to choose among a collection of possible actions, some of which are better in terms of utility, say, than others; perhaps there's even a single ``best'' action that maximizes utility.  The challenge is that this partial ordering of the possible actions depends on certain features of the world that are uncertain.  The classical approach to this problem, dating back to von Neumann and Morganstern, is to quantify uncertainty about the world using a precise probability distribution, find the expected utility with respect to this probability distribution, and then pick the action that maximizes this expected utility.  While there are strong mathematical justifications for this, it has been demonstrated that individuals' behavior is often inconsistent with maximizing expected utility in this sense.  So, there have been efforts to generalize von Neumann and Morganstern's formulation, including the use of imprecise lower/upper probabilities for evaluating the expected utility; see, e.g., \citet{denoeux.decision.2019} for a review.  Here I propose to use the e-possibilistic IM's lower and upper probabilities to assess the quality of various actions and, in particular, to make data-driven choices of optimal actions, very much in the spirit of classical Bayesian decision theory.  What's notable is that I can leverage the strong reliability properties of the e-possibilistic IM to establish new decision-theoretic reliability guarantees.  





\section{Proofs from the main paper}

\subsection{Proof of Proposition~\ref{prop:prior.reg}}

%\begin{proof}
By definition of the $\uprior$-upper expectation, 
\[ \uprior \, \rho = \int_0^1 \Bigl\{ \sup_{\theta: q(\theta) > s} \rho(\theta) \Bigr\} \, ds = \int_0^1 \Bigl\{ \sup_{\theta: q(\theta) > s} \frac{1}{\gamma \circ q(\theta)} \Bigr\} \, ds. \]
Since $\gamma$ is non-decreasing, 
\[ q(\theta) > s \implies \frac{1}{\gamma \circ q(\theta)} \leq \frac{1}{\gamma(s)}, \]
and from here the claim follows immediately from Equation~\eqref{eq:calibrate}.  
%\end{proof}


\subsection{Proof of Proposition~\ref{prop:pullback}}

%\begin{proof}
Take any fixed $\prior \in \credal$.  Non-emptiness of $\pbcred_\prior$ is a consequence of the classical results mentioned above.  To prove convexity, take two elements $\pb_1$ and $\pb_2$ in $\pbcred_\prior$ and a constant $\tau \in [0,1]$, and then define the mixture $\pb^\text{mix} = (1-\tau) \pb_1 + \tau \pb_2$; the goal is to show that $\pb^\text{mix} \in \pbcred_\prior$.  If $\Omega \sim \pb^\text{mix}$, then for any event $H$ in $\TT = f(\OO)$, 
\begin{align*}
\pb^\text{mix}\{ f(\Omega) \in H \} & = (1-\tau) \, \pb_1\{ f(\Omega) \in H\} + \tau \, \pb_2\{ f(\Omega) \in H \} \\
& = (1-\tau) \, \prior(H) + \tau \, \prior(H) \\
& = \prior(H).
\end{align*}
This implies that $\pb^\text{mix} \in \pbcred_\prior$ if $\pb_1$ and $\pb_2$ are and, therefore, that $\pbcred_\prior$ is convex.  Finally, to prove that $\pbcred_\prior$ is closed with respect to the weak topology, consider a sequence $(\pb_t: t \geq 1)$ in $\pbcred_\prior$ with a weak limit $\pb_\infty$; the goal is to show that $\pb_\infty \in \pbcred_\prior$.  By the continuous mapping theorem, if $\Omega_t \sim \pb_t$, then $f(\Omega_t) \to f(\Omega_\infty)$ in distribution as $t \to \infty$.  But the distribution of $f(\Omega_t)$ is $\prior$ for all $t$ and, consequently, the distribution of $f(\Omega_\infty)$ is also $\prior$.  This implies $\pb_\infty$ is contained in $\pbcred_\prior$ and, hence, the latter collection is closed. 
%\end{proof} 


\subsection{Proof of Theorem~\ref{thm:reg.ville}} 

%\begin{proof}
The first claim is just a direct computation using Equation~\eqref{eq:uuprob} when the function $g = \eval^\text{reg}$ factors as $g(\cdot, \omega) = \rho(f(\omega)) \, \eval_{f(\omega)}(\cdot)$:
\begin{align*}
\uuprob \, \eval^\text{reg} & = \sup_{\prior \in \credal} \sup_{\pb \in \pbcred_\prior} \E^{\Omega \sim \pb} \bigl[ \E^{Z \sim \prob_\Omega}\{ \eval^\text{reg}(Z^N, f(\Omega)) \} \bigr] \\
& = \sup_{\prior \in \credal} \sup_{\pb \in \pbcred_\prior} \E^{\Omega \sim \pb} \bigl[ \rho( f(\Omega) ) \,  \underbrace{\E^{Z \sim \prob_\Omega}\{ \eval_{f(\Omega)}(Z^N) \}}_{\text{$\leq 1$, by Equation \eqref{eq:eval.bound}}} \bigr] \\
& \leq \sup_{\prior \in \credal} \sup_{\pb \in \pbcred_\prior} \E^{\Omega \sim \pb} \bigl[ \rho\{ f(\Omega) \} \bigr] \\
& = \sup_{\prior \in \credal} \E^{\Theta \sim \prior} \{ \rho(\Theta) \} \\
& \leq 1,
\end{align*}
where the last inequality follows by definition of the regularizer $\rho$, and the penultimate equality follows by definition of $\pbcred_\prior$: if $\Omega \sim \pb \in \pbcred_\prior$, then the distribution of $f(\Omega)$ is the same as that of $\Theta$ under $\prior$.  The second claim, the regularized Ville's inequality in Equation~\eqref{eq:reg.ville}, follows from the first claim and an application of Markov's inequality inside the upper-probability calculation.  If I write ``$(Z,\Omega) \sim \prob_\bullet \otimes \pb$'' to represent the joint distribution of $(Z,\Omega)$ under the model where $\Omega \sim \pb$ and $(Z \mid \Omega=\omega) \sim \prob_\omega$, then:
\begin{align*}
\uuprob \bigl[ \eval^\text{\rm reg}\{Z^N, f(\Omega)\} > \alpha^{-1} \bigr] & = \sup_{\prior \in \credal} \sup_{\pb \in \pbcred_\prior} \prob^{(Z,\Omega) \sim \prob_\bullet \otimes \pb}\bigl[ \eval^\text{reg}\{Z^N, f(\Omega) \} > \alpha^{-1} \bigr] \\
& \leq \sup_{\prior \in \credal} \sup_{\pb \in \pbcred_\prior} \alpha \, \E^{\Omega \sim \pb} \bigl[ \E^{Z \sim \prob_\Omega}\{ \eval^\text{reg}(Z^N, f(\Omega)) \} \bigr] \\
& = \alpha \, \uuprob \, \eval^\text{reg} \\
& \leq \alpha,
\end{align*}
where the first ``$\leq$'' above is by the usual Markov's inequality and the last line is by the claim in Equation~\eqref{eq:reg.ville0} proved above.  
%\end{proof}


\subsection{Proof of Theorem~\ref{thm:valid.vac}}

%\begin{proof}
Since the IM with contour $\pi^\eval$ dominates that with contour $\pi$ defined earlier, the anytime validity of the latter implies that of the former.  For concreteness, however, I'll give a direct proof of anytime validity of $\uPi^\eval$.  

For any data set $z^n$, the IM's possibility measure output $H \mapsto \uPi_{z^n}(H)$ is monotone.  Therefore, if $\omega$ is such that $f(\omega) \in H$, then 
\[ \uPi_{z^n}^\eval(H) \geq \uPi_{z^n}^\eval(\{ f(\omega) \}) = \pi_{z^n}^\eval(f(\omega)), \]
and, consequently, for any $\alpha \in [0,1]$, 
\[ \uPi_{z^n}(H) \leq \alpha \implies \pi_{z^n}^\eval(f(\omega)) \leq \alpha \iff \eval_{f(\omega)}(z^n) \geq \alpha^{-1}, \]
where the right-most property is by definition of the IM's contour function in terms of the e-process's reciprocal.  It follows that 
\[ \sup_{\omega: f(\omega) \in H} \prob_\omega\{ \uPi_{Z^N}^\eval(H) \leq \alpha \} \leq \sup_{\omega: f(\omega) \in H} \prob_\omega\{ \eval_{f(\omega)}(Z^N) \geq \alpha^{-1} \} \leq \alpha, \]
with the last inequality due to Equation~\eqref{eq:ville}, thus completing the proof.
%\end{proof}


\subsection{Proof of Corollary~\ref{cor:uniform.vac}}

%\begin{proof}
The proof is based on the following observation:
\begin{equation}
\label{eq:equivalence} 
\uPi_{Z^N}^\eval(H) \leq \alpha \text{ for some $H$ with $H \ni f(\omega)$} \iff \pi_{Z^N}^\eval(f(\omega)) \leq \alpha. 
\end{equation}
The ``$\Longleftarrow$'' direction is obvious since $H := \{f(\omega)\}$ is a hypothesis that contains $f(\omega)$.  The ``$\Longrightarrow$'' is similarly obvious by the monotonicity property as stated in Equation~\eqref{eq:monotone}. Since Corollary~1 says the right-most event in the above display has probability no more than $\alpha$, uniformly in $\omega$ and in $N$, the same must be true of the equivalent left-most event.  
%\end{proof}


%\subsection{Proof of Theorem~4}
%
%%\begin{proof}
%Fix any hypothesis $H \subseteq \TT$.  I'll focus on proving that the first inequality in Equation~(28) doesn't hold, i.e., that 
%\begin{equation}
%\label{eq:no.contraction}
%\sup_{n,z^n} \uPi_{z^n}^{\eval \times \rho}(H) \geq \uprior(H). 
%\end{equation}
%The version involving lower probabilities is proved similarly.  To start, note that 
%\[ \inf_{n, z^n} \eval_\theta(z^n) \leq \sup_N \sup_{\omega: f(\omega)=\theta} \E_\omega \{ \eval_\theta(Z^N) \} \leq 1, \]
%where the first inequality is because an average is never smaller than the minimum, and the second inequality by Equation~(1).  That is, for any $\theta$, there exists $(n,z^n)$ such that $\eval_\theta(z^n)$ is no more than 1.  Now write out the left-hand of \eqref{eq:no.contraction} as follows:
%\[ \sup_{n,z^n} \uPi_{z^n}^{\eval \times \rho}(H) = \sup_{n,z^n} \sup_{\theta \in H} \pi_{z^n}^{\eval \times \rho}(\theta) = 1 \wedge \Bigl\{ \inf_{n,z^n} \inf_{\theta \in H} \rho(\theta) \, \eval_\theta(z^n) \Bigr\}^{-1}. \]
%It follows from the ``first observation'' above that 
%\[ \sup_{n,z^n} \uPi_{z^n}^{\eval \times \rho}(H) \geq 1 \wedge \Bigl\{ \inf_{\theta \in H} \rho(\theta) \Bigr\}^{-1}. \]
%Since $\rho(\theta) \geq \{ \inf_{\theta \in H} \rho(\theta) \} \times 1(\theta \in H)$, it's easy to see that 
%\[ \Bigl\{ \inf_{\theta \in H} \rho(\theta) \Bigr\} \times \uprior(H) \leq \uprior\,\rho \leq 1, \]
%where the first inequality is by monotonicity of the upper expectation and the second by definition of the regularizer $\rho$.  Plugging this bound into that above gives 
%\[ \sup_{n,z^n} \uPi_{z^n}^{\eval \times \rho}(H) \geq 1 \wedge \uprior(H) = \uprior(H), \]
%which completes the proof of the theorem. 
%%\end{proof}


\subsection{Proof of Theorem~\ref{thm:imdec}}

%\begin{proof}
Write $\Theta$ instead of $f(\Omega)$ for now.  First, observe that the integrand 
\[ s \mapsto \sup\{\ell_a(\theta): \pi_{Z^N}^{\eval \times \rho}(\theta) \geq s\} \]
in $\uPi_{Z^N}^{\eval \times \rho}(\ell_a)$ is non-decreasing.  Next, I proceed by considering two separate cases.  
\begin{enumerate}
\item Case: $\eval^\text{reg}(Z^N,\Theta) > 1$.  In this case,  $\pi_{Z^N}^{\eval \times \rho}(\Theta) = 1$ so, by the monotonicity property mentioned above, the pair $(Z^N,\Theta)$ is such that 
\[ \sup_{\theta: \pi_{Z^N}^{\eval \times \rho}(\theta) \geq s} \ell_a(\theta) \geq \ell_a(\Theta) \quad \text{for all $s \in [0,1]$}. \]
Of course, if the integrand is bounded below, then the integral, over all of $[0,1]$, is also bounded below by the same value.  Therefore, 
\[ \uPi_{Z^N}^{\eval \times \rho}(\ell_a) \geq \ell_a(\Theta), \]
which implies 
\begin{equation}
\label{eq:risk.bound}
\frac{\ell_a(\Theta)}{\uPi_{Z^N}^{\eval \times \rho}(\ell_a)} \leq 1 < \eval^\text{reg}(Z^N,\Theta). 
\end{equation}
\item Case: $\eval^\text{reg}(Z^N,\Theta) \leq 1$.  In this case, $\pi_{Z^N}^{\eval \times \rho}(\Theta) \leq 1$, so I can lower-bound the upper expected loss by truncating the range of integration as follows:
\begin{align*}
\uPi_{Z^N}^{\eval \times \rho}(\ell_a) & = \int_0^1 \Bigl\{ \sup_{\theta: \pi_{z^n}^{\eval \times \rho}(\theta) \geq s} \ell_a(\theta) \Bigr\} \, ds \\
& \geq \int_0^{\pi_{Z^N}^{\eval \times \rho}(\Theta)} \Bigl\{ \sup_{\theta: \pi_{z^n}^{\eval \times \rho}(\theta) \geq s} \ell_a(\theta) \Bigr\} \, ds \\
& \geq \pi_{Z^N}^{\eval \times \rho}(\Theta) \, \ell_a(\Theta),
\end{align*}
where the last inequality is again by the monotonicity property highlighted above.  On rearranging, and using the fact that $\pi_{Z^N}^{\eval \times \rho}(\Theta) = \eval_\Theta(Z^N) \, \rho(\Theta)$ in this case, I get 
\[ \frac{\ell_a(\Theta)}{\uPi_{Z^N}^{\eval \times \rho}(\ell_a)} \leq \eval^\text{reg}(Z^N,\Theta), \]
which is the same bound as in \eqref{eq:risk.bound}.  
\end{enumerate} 
Next, it's clear that the common bound derived in the two separate cases above holds uniformly in the actions, i.e., 
\[ \sup_{a \in \AA} \frac{\ell_a(\Theta)}{\uPi_{Z^N}^{\eval \times \rho}(\ell_a)} \leq \eval^\text{reg}(Z^N,\Theta). \]
Plugging $f(\Omega)$ back in for $\Theta$, taking $\uuprob$-expectation on both sides, and then applying the property in Equation~\eqref{eq:reg.ville0}, establishes the bound in Equation~\eqref{eq:dec.bound}. 
%\end{proof}




\bibliographystyle{apalike}
\bibliography{/Users/rgmarti3/Dropbox/Research/mybib}


\end{document}












